% Orthogonalité dans l’espace — Mathématiques 1re Tech — Cours signé OnBuch+
% Programme officiel MINESEC. Compiler avec : tectonic N19-orthogonalite-espace.tex
\newcommand{\DOCMATIERE}{Mathématiques}
\newcommand{\DOCNIVEAU}{1re Tech}
\newcommand{\DOCMODULE}{Mathématiques – classe de Première technique}
\newcommand{\DOCLECON}{N19}
\newcommand{\DOCTITRE}{Orthogonalité dans l’espace}
% OnBuch+ — préambule « cours signés » ENRICHI (compilé avec tectonic / XeLaTeX)
% Système visuel « Pop » d'OnBuch+ : fond crème (jamais blanc), bordures encre
% épaisses, ombres « dures » décalées, titres Archivo Black en capitales.
% Version MATHÉMATIQUES : graphes (pgfplots/tikz), boîtes pédagogiques
% (définition, propriété, méthode, attention, le savais-tu, exercice résolu,
% exercice, TP, bilan), page de garde et signature OnBuch+ ancrée en bas.
%
% Macros attendues dans main.tex AVANT \input{preamble} :
%   \DOCMATIERE  \DOCNIVEAU  \DOCMODULE  \DOCLECON (numéro)  \DOCTITRE
\documentclass[11pt]{article}

\usepackage{fontspec}
\usepackage[french]{babel}
\usepackage[a4paper,margin=2cm,top=3.4cm,bottom=2.8cm,headheight=1.4cm,footskip=1.3cm]{geometry}
\usepackage{amsmath,amssymb,amsfonts}
\usepackage{mathtools} % \xmapsto, \coloneqq... souvent utilisés par les modèles
\usepackage{cancel} % simplifications barrées
% \abs{z} (module, valeur absolue) et \norm{u} (norme), utilisés par les modèles
\providecommand{\abs}{}\let\abs\relax\DeclarePairedDelimiter{\abs}{\lvert}{\rvert}
\providecommand{\norm}{}\let\norm\relax\DeclarePairedDelimiter{\norm}{\lVert}{\rVert}
\usepackage[table]{xcolor} % [table] : fournit aussi \rowcolors (alternance de lignes)
\usepackage{pagecolor}
\usepackage{tikz}
\usetikzlibrary{arrows.meta,positioning,calc,shapes.geometric,shapes.misc,shapes.arrows,decorations.pathmorphing,decorations.pathreplacing,decorations.markings,patterns,shadows,mindmap,trees,fit,backgrounds,matrix,intersections,angles,quotes}
\usepackage{pgfplots}
\usepackage{pgf-pie} % diagrammes circulaires \pie (statistiques)
\pgfplotsset{compat=1.18}
\usepgfplotslibrary{fillbetween,groupplots} % aires entre courbes (calcul intégral)
\usepackage{enumitem}
\usepackage{fancyhdr}
% [most] charge toutes les bibliothèques tcolorbox (listings, minted, raster,
% documentation...) et gonfle inutilement la mémoire TeX sur un document long
% (~300 boîtes) : on ne charge que ce qui est réellement utilisé.
\usepackage[skins,breakable]{tcolorbox}
\usepackage{graphicx}
\usepackage{colortbl}
\usepackage{booktabs}
\usepackage{tabularx}
\usepackage{diagbox} % case d'en-tête diagonale des tableaux à double entrée
% Colonnes extensibles centrée / gauche / droite (C, L, R), souvent utilisées par les modèles
\newcolumntype{C}{>{\centering\arraybackslash}X}
\newcolumntype{L}{>{\raggedright\arraybackslash}X}
\newcolumntype{R}{>{\raggedleft\arraybackslash}X}
\usepackage{array}
\usepackage{multicol}
\usepackage{siunitx}
% parse-numbers=false : accepte aussi \SI{2,5 \times 10^{-3}}{...}
\sisetup{locale=FR,output-decimal-marker={,},group-minimum-digits=4,parse-numbers=false}
\DeclareSIUnit\ohm{\ensuremath{\symup{\Omega}}} % U+2126 absent de la police
\usepackage{qrcode}
% \degree : commande fréquemment utilisée par les modèles pour un angle ou une
% température en dehors de \SI{}{} (non fournie par les paquets chargés).
\providecommand{\degree}{^{\circ}}
% \qed (fin de démonstration), utilisé par les modèles sans amsthm
\providecommand{\qed}{\hfill$\square$}
% \barre : double barre des valeurs interdites dans un tableau de variations (tkz-tab)
\providecommand{\barre}{\ensuremath{\|}}
% environnement proof (amsthm non chargé) utilisé par les modèles
\ifdefined\proof\else\newenvironment{proof}[1][Démonstration]{\par\noindent\textit{#1.}\ }{\qed\par}\fi
\usepackage{lastpage}
\usepackage{hyperref}
\hypersetup{hidelinks}

% ============================================================
% Palette « Pop » OnBuch+ — mêmes hex que la classe P (plus_ui.dart)
% ============================================================
\definecolor{poporange}{HTML}{F59321}
\definecolor{popdark}{HTML}{E07A0C}
\definecolor{popcream}{HTML}{FFF6E8}
\definecolor{popink}{HTML}{1C1714}
\definecolor{poppurple}{HTML}{7A5AE0}
\definecolor{popgreen}{HTML}{2E9E6A}
\definecolor{poppink}{HTML}{E0457B}
\definecolor{popblue}{HTML}{2F7DE1}
\definecolor{popgold}{HTML}{C98A12}
\definecolor{popmuted}{HTML}{8A7F76}
\definecolor{popline}{HTML}{EADFD2}
\definecolor{popgray}{HTML}{8A7F76}
\colorlet{popgrey}{popgray}
% Alias tolérants (le modèle nomme parfois les couleurs différemment)
\colorlet{popwhite}{white}
\colorlet{popblack}{popink}
\colorlet{popred}{poppink}
\colorlet{popyellow}{popgold}
% Teintes claires pour fonds de tableaux / zones
\colorlet{popblueL}{popblue!10!white}
\colorlet{poporangeL}{poporange!14!white}
\colorlet{popgreenL}{popgreen!12!white}
\colorlet{poppurpleL}{poppurple!12!white}
\colorlet{poppinkL}{poppink!12!white}
\colorlet{popgoldL}{popgold!14!white}

\pagecolor{popcream}

% ============================================================
% Typographie
% ============================================================
\defaultfontfeatures{Ligatures=TeX}
\setmainfont{PlusJakartaSans-Medium.ttf}[Path=../../../fonts/,BoldFont=PlusJakartaSans-Bold.ttf]
% Maths en sans-serif (Fira Math) avec chiffres et lettres droites de Plus
% Jakarta, pour que les formules se fondent dans le texte.
\usepackage[math-style=ISO,bold-style=ISO]{unicode-math}
\setmathfont{FiraMath-Regular.otf}[Path=../../../fonts/]
\setmathfont{PlusJakartaSans-Medium.ttf}[Path=../../../fonts/,range={up/num,up/{Latin,latin},\mathup}]
\usepackage{icomma} % virgule décimale sans espace en mode math
% Caractères mathématiques Unicode absents des polices de texte (Plus Jakarta,
% Archivo Black) : rendus par leur équivalent mathématique, en texte comme en
% formule — sinon ils s'affichent en carré vide (ex. « Arithmétique dans ℤ »).
\usepackage{newunicodechar}
\newunicodechar{ℝ}{\ensuremath{\mathbb{R}}}
\newunicodechar{ℤ}{\ensuremath{\mathbb{Z}}}
\newunicodechar{ℂ}{\ensuremath{\mathbb{C}}}
\newunicodechar{ℕ}{\ensuremath{\mathbb{N}}}
\newunicodechar{ℚ}{\ensuremath{\mathbb{Q}}}
\newunicodechar{𝔻}{\ensuremath{\mathbb{D}}}
\newunicodechar{α}{\ensuremath{\alpha}}
\newunicodechar{β}{\ensuremath{\beta}}
\newunicodechar{γ}{\ensuremath{\gamma}}
\newunicodechar{δ}{\ensuremath{\delta}}
\newunicodechar{ε}{\ensuremath{\varepsilon}}
\newunicodechar{θ}{\ensuremath{\theta}}
\newunicodechar{λ}{\ensuremath{\lambda}}
\newunicodechar{μ}{\ensuremath{\mu}}
\newunicodechar{σ}{\ensuremath{\sigma}}
\newunicodechar{τ}{\ensuremath{\tau}}
\newunicodechar{φ}{\ensuremath{\varphi}}
\newunicodechar{ψ}{\ensuremath{\psi}}
\newunicodechar{ω}{\ensuremath{\omega}}
\newunicodechar{Σ}{\ensuremath{\Sigma}}
% majuscules produites par \MakeUppercase dans les titres (α-aminés -> Α)
\newunicodechar{Α}{\ensuremath{\alpha}}
\newunicodechar{Β}{\ensuremath{\beta}}
\newunicodechar{Γ}{\ensuremath{\Gamma}}
\newunicodechar{Δ}{\ensuremath{\Delta}}
\newunicodechar{Ε}{\ensuremath{\varepsilon}}
\newunicodechar{Θ}{\ensuremath{\Theta}}
\newunicodechar{Λ}{\ensuremath{\Lambda}}
\newunicodechar{Μ}{\ensuremath{\mu}}
\newunicodechar{Τ}{\ensuremath{\tau}}
\newunicodechar{Φ}{\ensuremath{\Phi}}
\newunicodechar{Ψ}{\ensuremath{\Psi}}
\newunicodechar{Ω}{\ensuremath{\Omega}}
\newunicodechar{↦}{\ensuremath{\mapsto}}
\newunicodechar{∈}{\ensuremath{\in}}
\newunicodechar{∉}{\ensuremath{\notin}}
\newunicodechar{∧}{\ensuremath{\wedge}}
\newunicodechar{⇔}{\ensuremath{\Leftrightarrow}}
\newunicodechar{⟺}{\ensuremath{\Longleftrightarrow}}
\newunicodechar{⇒}{\ensuremath{\Rightarrow}}
\newunicodechar{⟹}{\ensuremath{\Longrightarrow}}
\newunicodechar{∘}{\ensuremath{\circ}}
\newunicodechar{∪}{\ensuremath{\cup}}
\newunicodechar{∩}{\ensuremath{\cap}}
\newunicodechar{≡}{\ensuremath{\equiv}}
\newunicodechar{⊂}{\ensuremath{\subset}}
\newunicodechar{⊥}{\ensuremath{\perp}}
\newunicodechar{∃}{\ensuremath{\exists}}
\newunicodechar{∀}{\ensuremath{\forall}}
\newunicodechar{∛}{\ensuremath{\sqrt[3]{\,}}}
\newunicodechar{∜}{\ensuremath{\sqrt[4]{\,}}}
\newunicodechar{⃗}{\ensuremath{{}^{\to}}}
\newunicodechar{ⁿ}{\ifmmode^{n}\else\textsuperscript{\ensuremath{n}}\fi}
\newunicodechar{ˣ}{\ifmmode^{x}\else\textsuperscript{\ensuremath{x}}\fi}
\newunicodechar{ᵇ}{\ifmmode^{b}\else\textsuperscript{\ensuremath{b}}\fi}
\newunicodechar{ʳ}{\ifmmode^{r}\else\textsuperscript{\ensuremath{r}}\fi}
\newunicodechar{ᵘ}{\ifmmode^{u}\else\textsuperscript{\ensuremath{u}}\fi}
\newunicodechar{ʸ}{\ifmmode^{y}\else\textsuperscript{\ensuremath{y}}\fi}
\newunicodechar{ᵃ}{\ifmmode^{a}\else\textsuperscript{\ensuremath{a}}\fi}
\newunicodechar{ᵉ}{\ifmmode^{e}\else\textsuperscript{\ensuremath{e}}\fi}
\newunicodechar{ᵏ}{\ifmmode^{k}\else\textsuperscript{\ensuremath{k}}\fi}
\newunicodechar{ᵖ}{\ifmmode^{p}\else\textsuperscript{\ensuremath{p}}\fi}
\newunicodechar{ᴺ}{\ifmmode^{N}\else\textsuperscript{\ensuremath{N}}\fi}
\newunicodechar{ᶜ}{\ifmmode^{c}\else\textsuperscript{\ensuremath{c}}\fi}
\newunicodechar{ʰ}{\ifmmode^{h}\else\textsuperscript{\ensuremath{h}}\fi}
\newunicodechar{ᵠ}{\ifmmode^{q}\else\textsuperscript{\ensuremath{q}}\fi}
\newunicodechar{ₙ}{\ifmmode_{n}\else\textsubscript{\ensuremath{n}}\fi}
\newunicodechar{ₐ}{\ifmmode_{a}\else\textsubscript{\ensuremath{a}}\fi}
\newunicodechar{ᵢ}{\ifmmode_{i}\else\textsubscript{\ensuremath{i}}\fi}
\newunicodechar{ᵣ}{\ifmmode_{r}\else\textsubscript{\ensuremath{r}}\fi}
\newunicodechar{ₖ}{\ifmmode_{k}\else\textsubscript{\ensuremath{k}}\fi}
\newunicodechar{ᵦ}{\ifmmode_{\beta}\else\textsubscript{\ensuremath{\beta}}\fi}
\newunicodechar{ₓ}{\ifmmode_{x}\else\textsubscript{\ensuremath{x}}\fi}
\newfontfamily\popdisplay{ArchivoBlack-Regular.ttf}[Path=../../../fonts/]
\newfontfamily\poplogo{SpaceGrotesk-Bold.ttf}[Path=../../../fonts/]
\newfontfamily\popsemi{PlusJakartaSans-SemiBold.ttf}[Path=../../../fonts/]
\newfontfamily\popxbold{PlusJakartaSans-ExtraBold.ttf}[Path=../../../fonts/]
\newfontfamily\popxboldls{PlusJakartaSans-ExtraBold.ttf}[Path=../../../fonts/,LetterSpace=3.0]

\setlist[enumerate]{leftmargin=1.7em,itemsep=5pt,topsep=4pt}
\setlist[itemize]{leftmargin=1.7em,itemsep=4pt,topsep=4pt,label={\color{poporange}\textbullet}}
\setlength{\parindent}{0pt}
\setlength{\parskip}{5pt}
\renewcommand{\arraystretch}{1.25}

% pgfplots : style Pop par défaut
\pgfplotsset{
  every axis/.append style={axis line style={popink,line width=1pt},tick style={popink},
    grid style={popline,line width=0.5pt},label style={font=\small},tick label style={font=\footnotesize}},
  popcurve/.style={poporange,line width=1.8pt,no markers,smooth},
}
% Même style disponible dans un \draw TikZ simple (hors pgfplots)
\tikzset{popcurve/.style={poporange,line width=1.8pt}}

% ============================================================
% Pill / squiggle
% ============================================================
\newcommand{\poppill}[3]{%
  \tikz[baseline=-0.55ex]\node[draw=popink,line width=1.2pt,rounded corners=2.6pt,%
    fill=#2,inner xsep=6.5pt,inner ysep=3pt]{%
    \popxboldls\fontsize{7}{7}\selectfont\color{#3}\MakeUppercase{#1}};%
}
\newcommand{\popsquiggle}[1]{%
  \tikz[baseline=-0.4ex]{
    \draw[#1,line width=2.6pt,line cap=round]
      plot [smooth] coordinates {(0,0.09) (0.34,0.01) (0.68,0.21) (1.02,0.01) (1.36,0.10)};
  }%
}
% Mot-clé surligné (marqueur orange sous le texte)
\newcommand{\cle}[1]{{\popxbold\color{popdark}#1}}

% ============================================================
% En-tête / pied
% ============================================================
\pagestyle{fancy}
\fancyhf{}
\renewcommand{\headrulewidth}{0pt}
\renewcommand{\footrulewidth}{0pt}
\lhead{\raisebox{-0.15ex}{\poplogo\fontsize{13}{13}\selectfont\color{popink}OnBuch\color{poporange}+}}
\rhead{\poppill{\DOCMATIERE\ \textbullet\ \DOCNIVEAU\ \textbullet\ Leçon \DOCLECON}{white}{popink}}
\lfoot{\popxbold\fontsize{7.6}{7.6}\selectfont\color{popmuted}\MakeUppercase{Cours signé OnBuch+}\ \textbullet\ {\popsemi\DOCTITRE}}
\rfoot{\tikz[baseline=-0.6ex]\node[rounded corners=8.5pt,fill=popink,minimum height=17pt,minimum width=17pt,inner xsep=5pt,inner sep=0]{\popxbold\fontsize{8}{8}\selectfont\color{popcream}\thepage\,/\,\pageref*{LastPage}};}

% ============================================================
% Titres
% ============================================================
\newcommand{\chaptitre}[2]{%
  \begin{tcolorbox}[enhanced,colback=poporange,colframe=popink,boxrule=2.6pt,arc=11pt,
      drop shadow=popink,left=15pt,right=15pt,top=12pt,bottom=11pt,before skip=3pt,after skip=18pt]
    \poppill{#2}{popink}{popcream}\par\vspace{9pt}
    {\raggedright\hyphenpenalty=10000\exhyphenpenalty=10000\popdisplay\fontsize{22}{24}\selectfont\color{popcream}\MakeUppercase{#1}\par}
  \end{tcolorbox}
}
% Section numérotée : pastille orange avec le numéro + titre Archivo
\newcounter{popsec}
\newcommand{\coursec}[1]{%
  \stepcounter{popsec}\setcounter{popsub}{0}%
  \par\vspace{16pt}\noindent
  \begin{minipage}{\linewidth}
  \tikz[baseline=-0.7ex]\node[draw=popink,line width=1.6pt,fill=poporange,rounded corners=4pt,minimum size=22pt,inner sep=0]{\popdisplay\fontsize{12}{12}\selectfont\color{popcream}\thepopsec};%
  \hspace{8pt}{\popdisplay\fontsize{15}{17}\selectfont\color{popink}\MakeUppercase{#1}}\par
  \vspace{4pt}\noindent{\color{poporange}\rule{36pt}{3.6pt}}
  \end{minipage}\par\nopagebreak\vspace{8pt}
}
\newcounter{popsub}
\newcommand{\courssub}[1]{%
  \stepcounter{popsub}%
  \par\vspace{9pt}\noindent{\popxbold\fontsize{11.5}{13}\selectfont\color{popink}{\color{poporange}\thepopsec.\thepopsub}\hspace{6pt}#1}\par\nopagebreak\vspace{4pt}
}

% ============================================================
% Boîtes pédagogiques — même recette : fond blanc, bordure épaisse colorée,
% ombre dure encre, badge plein. Titre optionnel [#1].
% ============================================================
\tcbset{popbase/.style={breakable,enhanced,colback=white,boxrule=2.3pt,arc=9pt,
  shadow={3pt}{-3pt}{0pt}{popink},left=14pt,right=14pt,top=11pt,bottom=11pt,before skip=11pt,after skip=15pt,
  fontupper=\popsemi\fontsize{10.6}{14.5}\selectfont\color{popink},
  fontlower=\popsemi\fontsize{10.6}{14.5}\selectfont\color{popink}}}
% Coche dessinée (le glyphe ✓ n'existe pas dans les polices du document)
\newcommand{\popcheck}[1]{\tikz[baseline=-0.5ex]\draw[#1,line width=1.6pt,line cap=round,line join=round](-0.13,0.0)--(-0.04,-0.09)--(0.13,0.1);}
\newcommand{\popboxhead}[3]{\poppill{#1}{#2}{white}\ifx\relax#3\relax\else\hspace{6pt}{\popxbold\fontsize{10.6}{12}\selectfont\color{popink}#3}\fi\par\vspace{7pt}}

\newtcolorbox{aretenir}[1][]{popbase,colframe=popgreen,before upper={\popboxhead{À retenir}{popgreen}{#1}}}
\newtcolorbox{exemplebox}[1][]{popbase,colframe=poporange,before upper={\popboxhead{Exemple}{poporange}{#1}}}
\newtcolorbox{definition}[1][]{popbase,colframe=popblue,before upper={\popboxhead{Définition}{popblue}{#1}}}
\newtcolorbox{propriete}[1][]{popbase,colframe=poppurple,before upper={\popboxhead{Propriété}{poppurple}{#1}}}
\newtcolorbox{methode}[1][]{popbase,colframe=popgold,before upper={\popboxhead{Méthode}{popgold}{#1}}}
\newtcolorbox{attention}[1][]{popbase,colframe=poppink,before upper={\popboxhead{Attention}{poppink}{#1}}}
\newtcolorbox{experience}[1][]{popbase,colframe=popblue,colback=popblueL,before upper={\popboxhead{Expérience}{popblue}{#1}}}
% « Le savais-tu ? » : carte violette pleine (contexte, culture, Cameroun)
\newtcolorbox{savaistu}[1][]{popbase,colback=poppurple,colframe=popink,
  fontupper=\popsemi\fontsize{10.4}{14.2}\selectfont\color{white},
  before upper={\poppill{Le savais-tu\,?}{white}{poppurple}\ifx\relax#1\relax\else\hspace{6pt}{\popxbold\color{white}#1}\fi\par\vspace{7pt}}}
% Exercice résolu : énoncé en haut, \tcblower puis la solution sur fond crème
\newcounter{popexo}\newcounter{popexr}
\newtcolorbox{exoresolu}[1][]{popbase,colframe=poporange,colbacklower=popcream,
  segmentation style={popink,line width=1.2pt,dashed},
  before upper={\stepcounter{popexr}\popboxhead{Exercice résolu \thepopexr}{poporange}{#1}},
  before lower={\poppill{Solution}{popink}{popcream}\par\vspace{6pt}}}
% Exercice (non résolu en place) — la correction est dans la partie Corrigés
\newtcolorbox{exercice}[1][]{popbase,colframe=popink,
  before upper={\stepcounter{popexo}\popboxhead{Exercice \thepopexo}{popink}{#1}}}
% Corrigé
\newtcolorbox{corrige}[1][]{popbase,colframe=popgreen,colback=popgreenL,
  before upper={\popboxhead{Corrigé}{popgreen}{#1}}}
% Objectifs / prérequis / bilan
\newtcolorbox{objectifs}{popbase,colframe=popink,colback=poporangeL,
  before upper={\popboxhead{Objectifs de la leçon}{popink}{}}}
\newtcolorbox{prerequis}{popbase,colframe=popink,colback=popblueL,
  before upper={\popboxhead{Prérequis}{popblue}{}}}
\newtcolorbox{bilan}{popbase,colframe=popink,colback=white,boxrule=2.8pt,
  before upper={\popboxhead{Fiche bilan}{poporange}{L'essentiel à connaître pour l'examen}}}
% Figure encadrée avec légende
\newtcolorbox{popfigure}[1][]{enhanced,breakable=false,colback=white,colframe=popline,boxrule=1.4pt,arc=8pt,
  left=10pt,right=10pt,top=10pt,bottom=8pt,before skip=10pt,after skip=12pt,halign=center,
  lower separated=false,halign lower=center,
  fontlower=\popsemi\fontsize{9}{11.5}\selectfont\color{popmuted},
  #1}
\newcommand{\legende}[1]{\par\vspace{6pt}{\popsemi\fontsize{9}{11.5}\selectfont\color{popmuted}#1}}

% ============================================================
% Signature OnBuch+
% ============================================================
\newcommand{\poplogoinline}[1]{{\poplogo\fontsize{#1}{#1}\selectfont\color{popink}OnBuch\color{poporange}+}}
\newcommand{\signatureonbuch}{%
  \begin{tcolorbox}[enhanced,colback=popink,colframe=popink,boxrule=2.2pt,arc=10pt,
      drop shadow=poporange,left=14pt,right=14pt,top=10pt,bottom=10pt,before skip=0pt,after skip=0pt]
    \tikz[baseline=-0.7ex]\node[circle,fill=popgreen,draw=popcream,line width=1.3pt,minimum size=22pt,inner sep=0]{\popcheck{white}};%
    \hspace{9pt}\begin{minipage}[c]{0.62\linewidth}
      {\popxbold\fontsize{12}{13}\selectfont\color{popcream}Signé\ \textbullet\ {\poplogo OnBuch\color{poporange}+}}\par\vspace{2pt}
      {\raggedright\popsemi\fontsize{8}{10.5}\selectfont\color{popline}\MakeUppercase{Équipe pédagogique OnBuch+}\\\MakeUppercase{Conforme au programme officiel MINESEC}\par}
    \end{minipage}\hfill
    {\poplogo\fontsize{22}{22}\selectfont\color{popcream}OnBuch\color{poporange}+}
  \end{tcolorbox}%
}

% ============================================================
% Page de garde
% #1 titre · #2 matière · #3 niveau · #4 module · #5 n° leçon · #6 durée ·
% #7 description / objectifs (texte ou itemize)
% ============================================================
\newcommand{\pagegarde}[7]{%
  \thispagestyle{empty}
  \newgeometry{a4paper,margin=1.7cm,top=1.6cm,bottom=1.4cm}%
  \begin{tikzpicture}[remember picture,overlay]
    \fill[poporange!18] ([xshift=-3.2cm,yshift=-3.6cm]current page.north east) circle (4.6cm);
    \fill[poppurple!14] ([xshift=2.4cm,yshift=7.6cm]current page.south west) circle (3.8cm);
  \end{tikzpicture}%
  {\poplogo\fontsize{30}{30}\selectfont\color{popink}OnBuch\color{poporange}+}\hfill
  \raisebox{0.6ex}{\poppill{Cours signé \textbullet\ Premium}{popink}{popcream}}\par
  \vspace{1pt}\popsquiggle{poppurple}\par
  \vspace{8pt}
  \begin{tcolorbox}[enhanced,colback=poporange,colframe=popink,boxrule=2.8pt,arc=13pt,
      drop shadow=popink,left=18pt,right=18pt,top=12pt,bottom=13pt,after skip=10pt]
    \poppill{#2 \textbullet\ #3}{popink}{popcream}\hspace{5pt}\poppill{#4}{white}{popink}\par\vspace{8pt}
    {\popdisplay\fontsize{14}{14}\selectfont\color{popink}LEÇON #5}\par\vspace{4pt}
    {\raggedright\hyphenpenalty=10000\exhyphenpenalty=10000\popdisplay\fontsize{25}{27}\selectfont\color{popcream}\MakeUppercase{#1}\par}
    \vspace{8pt}
    {\popxbold\fontsize{9.5}{9.5}\selectfont\color{popink}\MakeUppercase{Durée conseillée : #6}}
  \end{tcolorbox}
  \begin{tcolorbox}[enhanced,colback=white,colframe=popink,boxrule=2.2pt,arc=10pt,
      drop shadow=popink,left=16pt,right=16pt,top=10pt,bottom=10pt,after skip=8pt]
    \poppill{Dans cette leçon}{poppurple}{white}\par\vspace{5pt}
    \popsemi\fontsize{9.3}{12.6}\selectfont\color{popink}#7
  \end{tcolorbox}
  \noindent\begin{minipage}[t]{0.487\linewidth}
    \begin{tcolorbox}[enhanced,colback=popgreen,colframe=popink,boxrule=2.2pt,arc=10pt,
        drop shadow=popink,left=11pt,right=11pt,top=10pt,bottom=10pt,height=3.1cm,valign=center]
      \tcbox[colback=white,colframe=popink,boxrule=1.3pt,arc=5pt,boxsep=0pt,nobeforeafter,
        left=3pt,right=3pt,top=3pt,bottom=3pt,valign=center]{\qrcode[height=1.9cm]{https://play.google.com/store/apps/details?id=com.luvvix.onbuch&hl=fr}}%
      \hspace{8pt}\begin{minipage}[c]{0.5\linewidth}
        {\popdisplay\fontsize{11}{12.5}\selectfont\color{white}\MakeUppercase{Télécharge OnBuch}}\par\vspace{4pt}
        {\raggedright\popsemi\fontsize{8.4}{11}\selectfont\color{white}Gratuit \textbullet\ Google Play\\Tuteur IA, annales, concours, résultats\par}
      \end{minipage}
    \end{tcolorbox}
  \end{minipage}\hfill
  \begin{minipage}[t]{0.487\linewidth}
    \begin{tcolorbox}[enhanced,colback=poppurple,colframe=popink,boxrule=2.2pt,arc=10pt,
        drop shadow=popink,left=11pt,right=11pt,top=10pt,bottom=10pt,height=3.1cm,valign=center]
      \tikz[baseline=-0.7ex]\node[circle,fill=white,draw=popink,line width=1.3pt,minimum size=1.6cm,inner sep=0]{\popdisplay\fontsize{24}{24}\selectfont\color{poppurple}+};%
      \hspace{8pt}\begin{minipage}[c]{0.6\linewidth}
        {\popdisplay\fontsize{11}{12.5}\selectfont\color{white}\MakeUppercase{Passe à OnBuch+}}\par\vspace{4pt}
        {\raggedright\popsemi\fontsize{8.4}{11}\selectfont\color{white}Tous les cours signés, Tuteur IA illimité, fiches PDF — onglet OnBuch+ de l'app\par}
      \end{minipage}
    \end{tcolorbox}
  \end{minipage}
  \vspace{8pt}
  \popcontenu
  \vfill
  \signatureonbuch
  \restoregeometry
  \newpage
}

% Rangée « Ce cours contient » (page de garde)
\newcommand{\poptile}[3]{% #1 couleur #2 gros texte #3 légende
  \begin{tcolorbox}[enhanced,colback=#1,colframe=popink,boxrule=2pt,arc=9pt,shadow={2.5pt}{-2.5pt}{0pt}{popink},
      left=6pt,right=6pt,top=9pt,bottom=9pt,halign=center,valign=center,height=1.75cm,width=\linewidth,nobeforeafter]
    {\popdisplay\fontsize{12.5}{13}\selectfont\color{white}\MakeUppercase{#2}}\par\vspace{3pt}
    {\centering\popsemi\fontsize{7.8}{9.5}\selectfont\color{white}#3\par}
  \end{tcolorbox}}
\newcommand{\popcontenu}{%
  \par\noindent{\popxboldls\fontsize{7.5}{7.5}\selectfont\color{popmuted}CE COURS SIGNÉ CONTIENT}\par\vspace{7pt}
  \noindent\begin{minipage}[t]{0.235\linewidth}\poptile{popblue}{Cours}{complet, illustré, pas à pas}\end{minipage}\hfill
  \begin{minipage}[t]{0.235\linewidth}\poptile{poporange}{Méthodes}{et exercices résolus}\end{minipage}\hfill
  \begin{minipage}[t]{0.235\linewidth}\poptile{poppink}{Exercices}{type examen + corrigés}\end{minipage}\hfill
  \begin{minipage}[t]{0.235\linewidth}\poptile{popgreen}{Fiche bilan}{l'essentiel à retenir}\end{minipage}\par}

% Sommaire
\newtcolorbox{sommaire}{enhanced,colback=white,colframe=popink,boxrule=2.2pt,arc=10pt,
  drop shadow=popink,left=18pt,right=18pt,top=13pt,bottom=13pt,before skip=6pt,after skip=20pt,
  before upper={\poppill{Sommaire}{poppurple}{white}\par\vspace{9pt}\popsemi\fontsize{10.6}{14.5}\selectfont\color{popink}}}

% Signature de fin de cours — poussée tout en bas de la dernière page
\newcommand{\findecours}{%
  \par\vfill\nopagebreak
  \begin{center}\popsquiggle{poporange}\end{center}\vspace{4pt}
  \signatureonbuch
}
\AtBeginDocument{\def\llbracket{\mathopen{[\mkern-3.5mu[}}\def\rrbracket{\mathclose{]\mkern-3.5mu]}}} % intervalles d'entiers (glyphe ⟦ absent de la police)
% Glyphes absents de Fira Math : redéfinis à partir de glyphes disponibles
\AtBeginDocument{%
  \def\setminus{\mathbin{\backslash}}\let\smallsetminus\setminus
  \def\vdots{\mathord{\vbox{\baselineskip3.2pt\lineskiplimit0pt\kern4pt\hbox{.}\hbox{.}\hbox{.}}}}%
  \def\ddots{\mathinner{\mkern1mu\raise6.5pt\hbox{.}\mkern2mu\raise3.6pt\hbox{.}\mkern2mu\raise0.7pt\hbox{.}\mkern1mu}}%
  \def\triangle{\mathord{\scalebox{0.9}{$\symup{\Delta}$}}}%
  \def\nabla{\mathord{\raisebox{\depth}{\scalebox{1}[-1]{$\symup{\Delta}$}}}}%
  \def\checkmark{\popcheck{popdark}}%
  \def\star{\mathbin{\ast}}\def\bigstar{\mathord{\ast}}%
  \def\diamond{\mathbin{\scalebox{0.6}{$\Diamond$}}}%
  \def\bigcup{\mathop{\mathchoice{\raisebox{-0.3em}{\scalebox{1.7}{$\cup$}}}{\scalebox{1.25}{$\cup$}}{\cup}{\cup}}\displaylimits}%
  \def\bigcap{\mathop{\mathchoice{\raisebox{-0.3em}{\scalebox{1.7}{$\cap$}}}{\scalebox{1.25}{$\cap$}}{\cap}{\cap}}\displaylimits}%
  \def\lesseqqgtr{\mathrel{\lessgtr}}\def\bigoplus{\mathop{\oplus}\displaylimits}%
}
\providecommand{\ssi}{si et seulement si} % abréviation fréquente
\AtBeginDocument{\providecommand{\wideparen}{\overparen}} % arc de cercle

%% FIGFIX-BEGIN v5 — correctifs globaux des figures (docs/onbuch-generation/tools/figfix)
\usepackage{adjustbox}\usepackage{etoolbox}\tcbuselibrary{raster}
% toute figure TikZ/pgfplots plus large que la ligne est réduite à la largeur disponible
\BeforeBeginEnvironment{tikzpicture}{\begin{adjustbox}{max width=\linewidth}}
\AfterEndEnvironment{tikzpicture}{\end{adjustbox}}
% légendes pgfplots lisibles par-dessus les courbes
\pgfplotsset{every axis/.append style={legend style={fill=white,fill opacity=0.92,text opacity=1,draw=black!15,inner sep=3pt}}}
% étiquettes d'arêtes (syntaxe edge["..."]) avec fond blanc
\tikzset{every edge quotes/.append style={fill=white,inner sep=1.2pt}}
%% FIGFIX-END
\begin{document}

\pagegarde{Orthogonalité dans l’espace}{Mathématiques}{1re Tech}{Mathématiques – classe de Première technique}{N19}{3 séances (fiche de progression harmonisée)}{Cette leçon établit les fondements de l'orthogonalité en géométrie dans l'espace : droites, plans et leurs relations mutuelles. Elle outille l'élève pour démontrer des perpendicularités, déterminer des plans orthogonaux et modéliser des situations techniques concrètes.
\vspace{4pt}
\begin{multicols}{2}\small
\begin{itemize}
\item À la fin de cette leçon, je sais définir et reconnaître deux droites orthogonales dans l'espace, y compris le cas des droites non coplanaires.
\item Je sais énoncer, démontrer et appliquer le critère de perpendicularité d'une droite à un plan.
\item Je sais déterminer si deux plans sont orthogonaux en utilisant la définition ou la propriété du diedre.
\item Je sais construire la droite orthogonale à un plan passant par un point donné (dans le plan ou hors du plan).
\item Je sais justifier rigoureusement une orthogonalité en rédigeant une démonstration structurée (données, propriété, conclusion).
\item Je sais modéliser un problème concret (architecture, charpente, éclairage) en identifiant les droites et plans orthogonaux.
\end{itemize}\end{multicols}}

\begin{sommaire}
\begin{multicols}{2}
\begin{enumerate}
\item 1. Droites orthogonales dans l'espace
\item 2. Droite orthogonale à un plan : Critère et Propriété fondamentale
\item 3. Plans orthogonaux : Définition, Critère et Propriétés
\item 4. Construction et Calculs : Applications techniques
\item 5. Synthèse et Résolution de problèmes complexes
\item Activité d'intégration
\item Méthodes et exercices résolus
\item Exercices
\item Corrigés des exercices
\item Fiche bilan
\end{enumerate}
\end{multicols}
\end{sommaire}

\begin{objectifs}
\begin{itemize}
\item À la fin de cette leçon, je sais définir et reconnaître deux droites orthogonales dans l'espace, y compris le cas des droites non coplanaires.
\item Je sais énoncer, démontrer et appliquer le critère de perpendicularité d'une droite à un plan.
\item Je sais déterminer si deux plans sont orthogonaux en utilisant la définition ou la propriété du diedre.
\item Je sais construire la droite orthogonale à un plan passant par un point donné (dans le plan ou hors du plan).
\item Je sais justifier rigoureusement une orthogonalité en rédigeant une démonstration structurée (données, propriété, conclusion).
\item Je sais modéliser un problème concret (architecture, charpente, éclairage) en identifiant les droites et plans orthogonaux.
\item Je sais utiliser la relation d'orthogonalité pour calculer des longueurs ou des angles dans un tétraèdre rectangle ou un pavé droit.
\item Je sais distinguer l'orthogonalité dans le plan (2D) de l'orthogonalité dans l'espace (3D) et éviter la confusion 'non coplanaires = non orthogonales'.
\end{itemize}
\end{objectifs}

\begin{prerequis}
\begin{itemize}
\item Orthogonalité dans le plan (classe de 4ème/3ème) : définition, propriété du cercle, théorème de Pythagore et réciproque.
\item Notions de base de la géométrie dans l'espace (Seconde/Début 1re) : droites, plans, positions relatives (coplanaires, sécantes, parallèles), parallélisme droite/plan et plan/plan.
\item Repérage dans l'espace : coordonnées d'un vecteur dans un repère orthonormé (si option vecteurs vue avant ou en parallèle).
\item Calcul algébrique : développement, factorisation, résolution d'équations du premier et second degré.
\item Logique mathématique : structure d'une démonstration (si... alors...), contraposée, réciproque.
\end{itemize}
\end{prerequis}

\newpage

\coursec{Droites orthogonales dans l'espace}

Au stade omnisports de Yaoundé, les projecteurs sont orientés de manière à éclairer la pelouse sans éblouir les spectateurs dans les tribunes. Chaque faisceau lumineux part d'un mât, traverse l'air, et frappe le sol selon une inclinaison précise. Comment les ingénieurs ont-ils calculé l'inclinaison parfaite de chaque faisceau par rapport au sol et aux gradins ? Comment vérifier, par exemple, qu'un câble de suspension est bien « à angle droit » avec une poutre qu'il ne touche même pas ?

\textbf{Problématique :} comment définir, reconnaître et démontrer l'orthogonalité de deux droites dans l'espace, y compris lorsqu'elles ne se coupent pas ?

Cette leçon vous donne les outils mathématiques pour modéliser et résoudre ce type de problème. Nous commençons par le plus simple : l'orthogonalité entre droites.

\courssub{Rappel : l'orthogonalité dans le plan}

En classe de Seconde, vous avez appris que deux droites d'un même plan sont \cle{perpendiculaires} lorsqu'elles se coupent en formant un angle droit ($90^\circ$). On note $d \perp d'$.

\begin{aretenir}[Orthogonalité dans le plan]
Deux droites \cle{sécantes} d'un plan sont dites \cle{perpendiculaires} si elles forment un angle droit. Pour le prouver, on dispose de plusieurs outils :
\begin{itemize}
\item le théorème de Pythagore et sa réciproque ;
\item les propriétés des figures usuelles (diagonales d'un losange, médiatrice, tangente et rayon d'un cercle) ;
\item le produit scalaire (vu en Première) : $\vec{u} \cdot \vec{v} = 0$ si et seulement si $\vec{u}$ et $\vec{v}$ sont orthogonaux.
\end{itemize}
\end{aretenir}

Dans le plan, la situation est simple : deux droites sont soit sécantes, soit parallèles. Mais dans l'espace, une troisième possibilité apparaît, et c'est toute la richesse de cette leçon.

\courssub{Deux droites dans l'espace : la nouveauté des droites non coplanaires}

Dans l'espace, deux droites peuvent être :
\begin{itemize}
\item \cle{sécantes} : elles se coupent en un point (elles sont alors coplanaires) ;
\item \cle{parallèles} : elles sont dans un même plan et ne se coupent pas ;
\item \cle{non coplanaires} : il n'existe \emph{aucun plan} qui les contienne toutes les deux. Elles ne se coupent pas et ne sont pas parallèles.
\end{itemize}

Prenons un objet que vous connaissez tous : une caisse en bois d'atelier, qui a la forme d'un pavé droit $ABCDEFGH$. L'arête $[AB]$ (le bord avant du couvercle) et l'arête $[CG]$ (l'arête verticale arrière droite) ne sont dans aucun plan commun : elles sont non coplanaires. Pourtant, intuitivement, elles ont des directions « à angle droit » : l'une est horizontale, l'autre verticale. C'est exactement cette intuition que la définition suivante rend rigoureuse.

\begin{definition}[Droites orthogonales dans l'espace]
Deux droites $d$ et $d'$ de l'espace sont dites \cle{orthogonales} si, en translatant l'une d'elles (parallèlement à elle-même) pour qu'elle devienne sécante à l'autre, les deux droites obtenues forment un \cle{angle droit}.

On note $d \perp d'$.

Autrement dit : $d \perp d'$ si les parallèles à $d$ et $d'$ menées par un même point quelconque de l'espace sont perpendiculaires.
\end{definition}

\begin{attention}[Vocabulaire : orthogonal ou perpendiculaire ?]
Le mot \cle{perpendiculaires} est le plus souvent réservé aux droites \textbf{sécantes} qui forment un angle droit. Le mot \cle{orthogonales} est plus général : il couvre à la fois les droites sécantes perpendiculaires et les droites \textbf{non coplanaires} dont les directions forment un angle droit. Ainsi :
\begin{center}
\textbf{Perpendiculaires} $\Rightarrow$ orthogonales, mais \textbf{orthogonales} $\nRightarrow$ perpendiculaires.
\end{center}
Dans la pratique, le symbole $\perp$ s'utilise dans les deux cas.
\end{attention}

\courssub{Le cas des droites orthogonales non coplanaires}

\begin{methode}[Prouver que deux droites non coplanaires sont orthogonales]
Pour démontrer que deux droites $d$ et $d'$ non coplanaires sont orthogonales :
\begin{enumerate}
\item \textbf{Translater} l'une des deux droites parallèlement à elle-même, de sorte qu'elle devienne sécante à l'autre (on choisit un point de passage pratique, souvent un sommet de la figure) ;
\item \textbf{Se placer dans le plan} ainsi formé par les deux droites sécantes ;
\item \textbf{Prouver l'angle droit} dans ce plan, avec les outils du plan (Pythagore, propriétés des figures, produit scalaire) ;
\item \textbf{Conclure} : $d \perp d'$.
\end{enumerate}
\end{methode}

\begin{exemplebox}[Arêtes orthogonales d'un pavé droit]
Soit $ABCDEFGH$ un pavé droit. Montrons que les arêtes $(AB)$ et $(CG)$ sont orthogonales.

Les droites $(AB)$ et $(CG)$ sont non coplanaires. On translate $(CG)$ parallèlement à elle-même : la droite $(BF)$ est parallèle à $(CG)$ (ce sont deux arêtes verticales du pavé) et $(BF)$ est sécante à $(AB)$ au point $B$. Or le pavé est droit, donc la face $ABFE$ est un rectangle : $(AB) \perp (BF)$. Par définition, on conclut :
\[ (AB) \perp (CG). \]
\end{exemplebox}

\begin{popfigure}
\begin{tikzpicture}[scale=1.15, line join=round]
% Pavé droit ABCDEFGH
\coordinate (A) at (0,0);
\coordinate (B) at (3.6,0);
\coordinate (C) at (4.6,1.4);
\coordinate (D) at (1,1.4);
\coordinate (E) at (0,2.6);
\coordinate (F) at (3.6,2.6);
\coordinate (G) at (4.6,4);
\coordinate (H) at (1,4);
% Faces cachées en pointillés
\draw[dashed, popmuted] (A)--(D) (D)--(C) (D)--(H);
% Arêtes visibles
\draw[popdark] (A)--(B) (B)--(C) (C)--(G) (G)--(H) (H)--(E) (E)--(A) (E)--(F) (F)--(B) (F)--(G);
% Arêtes mises en évidence
\draw[line width=1.6pt, popblue] (A)--(B);
\draw[line width=1.6pt, poporange] (C)--(G);
\draw[line width=1.4pt, popgreen] (B)--(F);
% Angle droit en B entre AB et BF
\draw[popgreen] (0.45,0)--(0.45,0.32)--(0,0.32);
% Sommets
\foreach \p/\pos in {A/below left, B/below, C/right, D/above right, E/left, F/below right, G/above right, H/above left}
  \fill[popdark] (\p) circle (1.3pt) node[\pos, font=\small] {$\p$};
% Annotations
\node[popblue, font=\small] at (1.8,-0.35) {$(AB)$};
\node[poporange, font=\small] at (5.05,2.7) {$(CG)$};
\node[popgreen, font=\small] at (3.35,1.3) {$(BF)$};
\node[popgreen, font=\footnotesize, align=center] at (1.6,0.75) {translation de $(CG)$\\ angle droit en $B$};
\end{tikzpicture}
\legende{Pavé droit $ABCDEFGH$ : les arêtes $(AB)$ et $(CG)$ sont non coplanaires mais orthogonales. On translate $(CG)$ en $(BF)$, sécante à $(AB)$ en $B$, où l'angle droit apparaît.}
\end{popfigure}

\courssub{Propriété fondamentale : orthogonalité et parallélisme}

Voici une propriété très utile dans les démonstrations : l'orthogonalité « se transmet » par parallélisme.

\begin{propriete}[Orthogonalité et droites parallèles]
Si une droite $d$ est orthogonale à une droite $\Delta$, alors $d$ est orthogonale à \textbf{toute} droite $\Delta'$ parallèle à $\Delta$.
\[ d \perp \Delta \quad \text{et} \quad \Delta' \parallel \Delta \quad \Longrightarrow \quad d \perp \Delta'. \]
De même, si $d \perp \Delta$ et $d' \parallel d$, alors $d' \perp \Delta$.
\end{propriete}

\begin{experience}[Pourquoi c'est vrai : démonstration guidée]
Soit $d \perp \Delta$ et $\Delta' \parallel \Delta$. Raisonnons pas à pas.
\begin{enumerate}
\item Par définition de l'orthogonalité, menons par un point $O$ de $d$ la parallèle $\Delta_0$ à $\Delta$ : alors $d$ et $\Delta_0$ sont sécantes en $O$ et perpendiculaires.
\item Menons par le même point $O$ la parallèle $\Delta'_0$ à $\Delta'$. Comme $\Delta' \parallel \Delta$ et $\Delta_0 \parallel \Delta$, on a $\Delta'_0 \parallel \Delta_0$ (transitivité du parallélisme). Or deux droites parallèles passant par un même point sont \textbf{confondues} : $\Delta'_0 = \Delta_0$.
\item Donc la parallèle à $\Delta'$ menée par $O$ est perpendiculaire à $d$ : par définition, $d \perp \Delta'$. \hfill $\square$
\end{enumerate}
\textbf{Idée à retenir :} l'orthogonalité ne dépend que des \emph{directions} des droites, pas de leur position. C'est pourquoi la translation dans la définition est légitime.
\end{experience}

\begin{exemplebox}[Application immédiate]
Dans le pavé droit $ABCDEFGH$, on a montré que $(AB) \perp (CG)$. Comme $(CG) \parallel (DH)$ (arêtes verticales), la propriété donne directement $(AB) \perp (DH)$, sans refaire de démonstration. De même, $(AB) \perp (AE)$ et $(AB) \perp (BF)$ car ce sont des arêtes sécantes perpendiculaires du rectangle $ABFE$.
\end{exemplebox}

\courssub{Un repère de l'espace : trois droites deux à deux orthogonales}

Dans un \cle{repère orthonormé} de l'espace $(O\,;\vec{i},\vec{j},\vec{k})$, les trois axes $(Ox)$, $(Oy)$ et $(Oz)$ sont sécants en $O$ et \textbf{deux à deux orthogonaux}. C'est la situation idéale : chaque axe est perpendiculaire aux deux autres. C'est exactement la configuration du coin d'une pièce : les deux arêtes du sol et l'arête verticale du mur.

\begin{popfigure}
\begin{tikzpicture}[scale=1.3, >=Stealth]
\coordinate (O) at (0,0);
% Axes
\draw[->, line width=1.2pt, popblue] (O)--(3,0) node[below right] {$x$};
\draw[->, line width=1.2pt, poporange] (O)--(1.5,1.3) node[above right] {$y$};
\draw[->, line width=1.2pt, popgreen] (O)--(0,2.6) node[above left] {$z$};
% Angles droits
\draw[popdark] (0.35,0)--(0.35,0.28)--(0,0.28);
\draw[popdark] (0.17,0.15)--(0.02,0.42)--(-0.15,0.27);
% Vecteurs unitaires
\draw[->, line width=1.4pt, popblue] (O)--(1,0) node[midway, below, font=\small] {$\vec{i}$};
\draw[->, line width=1.4pt, poporange] (O)--(0.6,0.52) node[midway, right, font=\small] {$\vec{j}$};
\draw[->, line width=1.4pt, popgreen] (O)--(0,1) node[midway, left, font=\small] {$\vec{k}$};
\fill[popdark] (O) circle (1.5pt) node[below left] {$O$};
\end{tikzpicture}
\legende{Repère orthonormé de l'espace : les trois axes sont sécants en $O$ et deux à deux orthogonaux (perpendiculaires).}
\end{popfigure}

\begin{attention}[Erreur fréquente : non coplanaires ne veut pas dire non orthogonales]
Beaucoup d'élèves pensent : « ces deux droites ne se coupent pas, donc elles ne peuvent pas être orthogonales ». C'est \textbf{faux} ! Deux droites non coplanaires \emph{peuvent} être orthogonales : c'est le cas de $(AB)$ et $(CG)$ dans le pavé. Retenez :
\begin{itemize}
\item non coplanaires $\Rightarrow$ non sécantes, donc jamais \emph{perpendiculaires} au sens strict ;
\item mais non coplanaires $\nRightarrow$ non orthogonales : tout dépend des \textbf{directions}.
\end{itemize}
À l'inverse, deux droites non coplanaires ne sont pas toujours orthogonales : il faut le \emph{démontrer} par la méthode de translation.
\end{attention}

\courssub{Exemples concrets de la vie technique}

L'orthogonalité de droites non coplanaires est partout autour de nous :

\begin{itemize}
\item \textbf{Terrain de football :} le poteau vertical d'un but est perpendiculaire (sécant) aux lignes de touche horizontales qui partent de son pied. Mais il est aussi \emph{orthogonal} à la ligne de touche \emph{opposée}, qu'il ne rencontre jamais : leurs directions (verticale et horizontale) forment un angle droit.
\item \textbf{Atelier de mécanique :} l'axe vertical d'une perceuse à colonne est orthogonal à tout rail horizontal de la table, même à ceux qu'il ne croise pas.
\item \textbf{Bâtiment :} dans une salle de classe, l'arête verticale d'un mur est orthogonale à toutes les arêtes horizontales du sol et du plafond, y compris celles du mur opposé (non coplanaires).
\end{itemize}

\begin{savaistu}[Les ponts à haubans]
En architecture et en génie civil, les \cle{haubans} d'un pont à haubans — ces câbles inclinés qui relient le pylône au tablier — sont souvent conçus pour être orthogonaux à certaines poutres du tablier, sans pour autant être coplanaires avec elles. Cette orthogonalité garantit une meilleure répartition des forces de traction et évite des contraintes de cisaillement parasites. Les ingénieurs du pont de la Wouri à Douala, comme ceux des grands ponts modernes, utilisent en permanence ces calculs d'orthogonalité dans l'espace !
\end{savaistu}

\courssub{Exercice résolu : orthogonalité dans un cube}

\begin{exoresolu}[Diagonales orthogonales dans un cube]
\textbf{Énoncé.} Soit $ABCDEFGH$ un cube de côté $6$ cm, avec $ABCD$ la face du bas et $EFGH$ la face du haut ($E$ au-dessus de $A$, etc.). On considère la grande diagonale $[AG]$ et la diagonale $[BD]$ de la face du bas.
\begin{enumerate}
\item Montrer que les droites $(AG)$ et $(BD)$ sont non coplanaires.
\item Montrer que $(AG) \perp (BD)$.
\end{enumerate}
\tcblower
\textbf{Solution détaillée.}

\textbf{1. Non-coplanarité.}
La droite $(BD)$ est contenue dans le plan $(ABC)$ (face du bas). La droite $(AG)$ coupe ce plan au seul point $A$ (car $G$ n'est pas dans le plan du bas). Or $A \notin (BD)$ (les diagonales du carré $ABCD$ se coupent en leur centre, pas en $A$). Donc $(AG)$ n'est pas contenue dans un plan commun avec $(BD)$ : les droites $(AG)$ et $(BD)$ sont \textbf{non coplanaires}.

\textbf{2. Orthogonalité.}
\emph{Méthode 1 : Calcul vectoriel (rigoureux et direct).}
Munissons l'espace du repère $(A\,;\overrightarrow{AB},\overrightarrow{AD},\overrightarrow{AE})$, orthonormé car le cube a des arêtes perpendiculaires de même longueur. Alors :
\[ \overrightarrow{AG} = \overrightarrow{AB} + \overrightarrow{AD} + \overrightarrow{AE} \quad \text{donc} \quad \overrightarrow{AG}\,(6\,;6\,;6), \qquad \overrightarrow{BD} = \overrightarrow{AD} - \overrightarrow{AB} \quad \text{donc} \quad \overrightarrow{BD}\,(-6\,;6\,;0). \]
\[ \overrightarrow{AG} \cdot \overrightarrow{BD} = 6\times(-6) + 6\times 6 + 6\times 0 = -36 + 36 + 0 = 0. \]
Le produit scalaire est nul : $\boxed{(AG) \perp (BD)}$, bien que ces droites soient non coplanaires.

\emph{Méthode 2 : Raisonnement géométrique par translation (esquisse).}
Soit $O$ le centre du carré $ABCD$ (intersection des diagonales). On translate $(AG)$ parallèlement à elle-même pour la faire passer par $O$ ; cette parallèle, notée $d$, est contenue dans le plan $(ACG)$ car $O \in (AC)$. Dans le plan $(ACG)$, la droite $d$ passe par $O$. La droite $(BD)$ passe aussi par $O$ et est perpendiculaire à $(AC)$ (propriété du carré). De plus, $(BD)$ est horizontale (dans le plan de la base) tandis que la direction de $(CG)$ (verticale) est orthogonale à $(BD)$ ; or $d$ a même direction que $(AG)$ qui est combinaison de $(AC)$ (horizontal) et $(CG)$ (vertical). On en déduit que $d \perp (BD)$ en $O$, donc $(AG) \perp (BD)$. \textit{Remarque : cette démonstration géométrique complète utilise implicitement le critère « droite orthogonale à un plan » qui sera formalisé en Section 2.}
\end{exoresolu}

\begin{exercice}[À vous de jouer]
Soit $ABCDEFGH$ un pavé droit tel que $AB = 8$ cm, $AD = 5$ cm et $AE = 4$ cm.
\begin{enumerate}
\item Citer quatre droites orthogonales à $(AB)$, dont deux non coplanaires avec $(AB)$.
\item Montrer que $(EH) \perp (CG)$.
\item La droite $(AC)$ est-elle orthogonale à $(BF)$ ? Justifier par la méthode de translation.
\end{enumerate}
\end{exercice}

\begin{corrige}[Exercice : pavé droit]
\textbf{1.} Droites sécantes à $(AB)$ et perpendiculaires (dans les faces $ABCD$ et $ABFE$) : $(AD)$, $(AE)$, $(BC)$, $(BF)$. \\
Droites non coplanaires avec $(AB)$ et orthogonales (par translation des précédentes via la propriété des parallèles) : $(CG)$, $(DH)$, $(EH)$, $(FG)$. \\
On peut citer par exemple : $(AD)$, $(AE)$, $(CG)$, $(DH)$.

\textbf{2.} $(EH) \parallel (AD)$ (côtés opposés du rectangle $ADHE$). \\
Dans le pavé droit, la face $ADHE$ est un rectangle, donc $(AD) \perp (DH)$. \\
Or $(DH) \parallel (CG)$ (arêtes verticales). \\
Par la propriété « orthogonalité et parallélisme » : $(AD) \perp (DH)$ et $(DH) \parallel (CG)$ $\Rightarrow$ $(AD) \perp (CG)$. \\
À nouveau par la même propriété : $(EH) \parallel (AD)$ et $(AD) \perp (CG)$ $\Rightarrow$ $\boxed{(EH) \perp (CG)}$.

\textbf{3.} On translate $(BF)$ en $(AE)$ (parallèles, arêtes verticales), qui est sécante à $(AC)$ en $A$. \\
\textit{Rappel (Seconde) :} dans un pavé droit, une arête latérale est perpendiculaire au plan de la base, donc à toute droite de ce plan passant par son pied. Ainsi $(AE) \perp (ABCD)$, et comme $(AC) \subset (ABCD)$ et passe par $A$, on a $(AE) \perp (AC)$. \\
Par définition de l'orthogonalité dans l'espace (translation) : $\boxed{(AC) \perp (BF)}$, alors qu'elles sont non coplanaires.
\end{corrige}

\begin{aretenir}[Bilan de la section]
\begin{itemize}
\item Deux droites de l'espace sont \cle{orthogonales} si, après translation de l'une pour les rendre sécantes, elles forment un angle droit. Notation : $d \perp d'$.
\item \cle{Perpendiculaires} = orthogonales \emph{et} sécantes. Des droites non coplanaires peuvent être orthogonales.
\item L'orthogonalité se transmet par parallélisme : si $d \perp \Delta$ et $\Delta' \parallel \Delta$, alors $d \perp \Delta'$.
\item La méthode reine : \textbf{translater}, se ramener à un plan, prouver l'angle droit dans ce plan (Pythagore, produit scalaire, figures usuelles).
\item Le calcul vectoriel (produit scalaire nul) est une preuve universelle, valable même pour des droites non coplanaires.
\end{itemize}
Dans la prochaine section, nous passerons du couple « droite-droite » au couple « droite-plan » : quand une droite est-elle orthogonale à tout un plan, comme le mât d'un projecteur par rapport à la pelouse du stade ?
\end{aretenir}

\coursec{Droite orthogonale à un plan : Critère et Propriété fondamentale}

Dans la section précédente, nous avons appris à reconnaître deux droites orthogonales dans l'espace : deux droites sont orthogonales lorsque la parallèle à l'une menée par un point de l'autre leur est perpendiculaire. Nous allons maintenant franchir une étape décisive : au lieu de comparer une droite à une autre droite, nous allons comparer une \cle{droite} à un \cle{plan} tout entier. C'est la notion de \cle{droite orthogonale à un plan}, sans doute la plus utile de toute la géométrie dans l'espace : c'est elle qui permet de calculer des hauteurs, des volumes et des distances, et c'est elle qui garantit qu'un poteau est bien vertical sur un chantier.

\courssub{Intuition : quand une droite est-elle « debout » sur un plan ?}

Imagine un piquet planté dans le sol. Le sol est un plan $(P)$, le piquet est une droite $(d)$ qui le perce en un point $H$. Quand dit-on que le piquet est « bien droit » ? L'ouvrier du chantier le vérifie avec un niveau à bulle ou un fil à plomb : le piquet doit faire un angle droit avec le sol \emph{dans toutes les directions}. Autrement dit, quelle que soit la droite tracée au sol et passant par le pied du piquet, le piquet doit être perpendiculaire à cette droite.

C'est exactement cette idée que les mathématiques retiennent : une droite est orthogonale à un plan si elle est orthogonale à \textbf{toutes} les droites du plan passant par son point d'intersection avec le plan.

\begin{definition}[Droite orthogonale à un plan, pied de la perpendiculaire]
Soit $(d)$ une droite et $(P)$ un plan sécants en un point $H$.
On dit que $(d)$ est \cle{orthogonale} (ou \cle{perpendiculaire}) au plan $(P)$ lorsque $(d)$ est orthogonale à \textbf{toute droite} de $(P)$ passant par $H$.
On note $(d) \perp (P)$.
Le point $H$, intersection de $(d)$ et de $(P)$, est appelé le \cle{pied de la perpendiculaire} à $(P)$ issue de $(d)$.
\end{definition}

\begin{attention}[Orthogonalité et perpendicularité dans l'espace]
Dans l'espace, deux droites orthogonales ne sont pas forcément sécantes. Mais ici, comme $(d)$ coupe $(P)$ en $H$, toute droite de $(P)$ passant par $H$ est \emph{sécante} avec $(d)$ : l'orthogonalité se traduit alors concrètement par un angle droit visible dans un plan contenant $(d)$ et cette droite. On marque ces angles droits par le symbole usuel sur les figures.
\end{attention}

Le problème pratique est évident : un plan contient une \emph{infinité} de droites passant par $H$. Impossible de toutes les vérifier une par une ! Le théorème suivant, appelé \cle{critère de perpendicularité}, va nous sauver : il suffit d'en vérifier \textbf{deux}, à condition qu'elles soient sécantes.

\courssub{Le critère de perpendicularité droite--plan}

\begin{propriete}[Critère de perpendicularité d'une droite à un plan]
Soit $(d)$ une droite et $(P)$ un plan. Si $(d)$ est orthogonale à \textbf{deux droites sécantes} de $(P)$, alors $(d)$ est orthogonale au plan $(P)$, c'est-à-dire orthogonale à \textbf{toute} droite de $(P)$ passant par son pied.
\par
Ce théorème est admis au Probatoire technique, mais sa démonstration est formatrice et peut être demandée sous forme guidée.
\end{propriete}

\begin{experience}[Démonstration guidée du critère]
\textbf{Hypothèses :} $(d)$ coupe $(P)$ en $H$ ; $(D_1)$ et $(D_2)$ sont deux droites de $(P)$ sécantes en $H$ ; $(d) \perp (D_1)$ et $(d) \perp (D_2)$.
\textbf{Conclusion à prouver :} pour toute droite $(\Delta)$ de $(P)$ passant par $H$, $(d) \perp (\Delta)$.

\begin{enumerate}
\item \textbf{Mise en place.} Sur $(d)$, de part et d'autre de $H$, on place deux points $M$ et $M'$ tels que $HM = HM'$. Soit $(\Delta)$ une droite quelconque de $(P)$ passant par $H$, distincte de $(D_1)$ et $(D_2)$.
\item \textbf{Médiatrices.} Puisque $(d) \perp (D_1)$ en $H$ avec $HM = HM'$, la droite $(D_1)$ est la \cle{médiatrice} du segment $[MM']$ dans le plan $(d, D_1)$. Donc tout point de $(D_1)$ est équidistant de $M$ et $M'$. De même, tout point de $(D_2)$ est équidistant de $M$ et $M'$.
\item \textbf{Choix de deux points.} La droite $(\Delta)$ coupe le segment joignant un point $A \in (D_1)$ à un point $B \in (D_2)$ (on choisit $A$ et $B$ de part et d'autre de $(\Delta)$) en un point $C$. Par l'étape 2 : $AM = AM'$ et $BM = BM'$.
\item \textbf{Triangles isométriques.} Les triangles $ABM$ et $ABM'$ ont leurs trois côtés égaux deux à deux ($AB$ commun, $AM = AM'$, $BM = BM'$). Ils sont donc isométriques (cas SSS). Il en résulte que $\angle BAM = \angle BAM'$. Les triangles $ACM$ et $ACM'$ ont alors : $AC$ commun, $AM = AM'$, $\angle CAM = \angle CAM'$ ; ils sont isométriques (cas SAS), d'où $CM = CM'$.
\item \textbf{Conclusion.} Le triangle $MCM'$ est isocèle en $C$, et $H$ est le milieu de $[MM']$ : la médiane $(CH)$ est donc aussi la hauteur, c'est-à-dire $(\Delta) \perp (d)$ en $H$. Comme $(\Delta)$ était quelconque, $(d)$ est orthogonale à toute droite de $(P)$ passant par $H$ : $(d) \perp (P)$. $\blacksquare$
\end{enumerate}
\emph{Remarque (approche vectorielle) :} si $\vec{u}$ dirige $(d)$ et si $\vec{u} \cdot \vec{v}_1 = 0$ et $\vec{u} \cdot \vec{v}_2 = 0$ pour deux directions non colinéaires $\vec{v}_1, \vec{v}_2$ du plan, alors $\vec{u}$ est orthogonal à toute combinaison $a\vec{v}_1 + b\vec{v}_2$, donc à toute direction du plan.
\end{experience}

\begin{popfigure}
\begin{center}
\begin{tikzpicture}[scale=1.05, line cap=round, line join=round]
% Plan (P) en perspective cavalière
\fill[popblueL, opacity=0.55] (0,0) -- (4.6,0) -- (6.1,1.9) -- (1.5,1.9) -- cycle;
\draw[popblue, thick] (0,0) -- (4.6,0) -- (6.1,1.9) -- (1.5,1.9) -- cycle;
\node[popblue] at (5.6,1.55) {$(P)$};
% Pied H
\coordinate (H) at (3.1,0.95);
% Droites D1 et D2 dans le plan passant par H
\draw[popdark, thick] (1.9,0.35) -- (4.4,1.6) node[right] {\small $(D_1)$};
\draw[popdark, thick] (2.2,1.75) -- (4.1,0.2) node[right] {\small $(D_2)$};
% Droite (d) orthogonale au plan en H
\draw[poporange, very thick] (3.1,-0.35) -- (3.1,3.4) node[above] {$(d)$};
% Angles droits
\draw[popink] (3.1,0.95) -- (3.32,1.02) -- (3.34,0.82);
\draw[popink] (3.1,0.95) -- (3.02,1.14) -- (2.84,1.08);
% Point H
\fill[popink] (H) circle (1.6pt);
\node[popink, below left] at (H) {$H$};
\end{tikzpicture}
\end{center}
\legende{La droite $(d)$ est orthogonale aux deux droites sécantes $(D_1)$ et $(D_2)$ du plan $(P)$ : d'après le critère, $(d) \perp (P)$. Les angles droits sont marqués en $H$, pied de la perpendiculaire.}
\end{popfigure}

\begin{attention}[Le piège classique : une seule droite ne suffit pas]
Prouver que $(d)$ est orthogonale à \textbf{une seule} droite de $(P)$ ne permet \textbf{pas} de conclure que $(d) \perp (P)$. Contre-exemple : un piquet penché sur le sol peut très bien faire un angle droit avec \emph{une} ligne tracée au sol (celle qui est perpendiculaire à sa direction d'inclinaison) sans être vertical.
Il faut impérativement \textbf{deux droites sécantes} du plan. Et attention : deux droites \emph{parallèles} du plan ne conviennent pas non plus — elles ne « balayent » qu'une seule direction.
\end{attention}

\courssub{Propriété fondamentale}

Le critère sert à \emph{prouver} qu'une droite est orthogonale à un plan. La propriété suivante sert, elle, à \emph{utiliser} cette orthogonalité : c'est la réciproque immédiate de la définition, et c'est elle que l'on invoque dans presque tous les calculs.

\begin{propriete}[Propriété fondamentale]
Si une droite $(d)$ est orthogonale à un plan $(P)$ en un point $H$, alors $(d)$ est orthogonale (et même perpendiculaire, car sécante) à \textbf{toute droite} de $(P)$ passant par $H$.
En particulier, pour tout point $M$ de $(P)$, le triangle formé par $H$, $M$ et un point de $(d)$ est rectangle en $H$, ce qui permet d'appliquer le \textbf{théorème de Pythagore}.
\end{propriete}

\begin{aretenir}[Le duo à connaître par cœur]
\begin{itemize}
\item \textbf{Pour prouver} $(d) \perp (P)$ : je montre que $(d)$ est orthogonale à \emph{deux droites sécantes} de $(P)$ (critère).
\item \textbf{Pour utiliser} $(d) \perp (P)$ : j'en déduis que $(d)$ est perpendiculaire à \emph{toute} droite de $(P)$ passant par le pied $H$, et j'écris Pythagore dans un triangle rectangle (propriété fondamentale).
\end{itemize}
\end{aretenir}

\begin{methode}[Prouver qu'une droite est orthogonale à un plan]
Pour prouver que $(d) \perp (P)$ :
\begin{enumerate}
\item \textbf{Identifier} dans le plan $(P)$ deux droites \emph{sécantes} $(D_1)$ et $(D_2)$ (souvent deux arêtes d'une face d'un solide, ou une arête et une médiane/hauteur).
\item \textbf{Prouver} séparément $(d) \perp (D_1)$ puis $(d) \perp (D_2)$, en utilisant les hypothèses : triangles isocèles, faces carrées ou rectangles, Pythagore réciproque, orthogonalités déjà connues.
\item \textbf{Conclure} par le critère : « $(d)$ est orthogonale à deux droites sécantes de $(P)$, donc $(d) \perp (P)$ », en précisant le pied $H$.
\end{enumerate}
\end{methode}

\courssub{Distance d'un point à un plan}

\begin{definition}[Distance d'un point à un plan]
Soit $A$ un point et $(P)$ un plan. La perpendiculaire à $(P)$ passant par $A$ coupe $(P)$ en un unique point $H$, appelé \cle{projeté orthogonal} de $A$ sur $(P)$ (c'est le pied de la perpendiculaire).
La \cle{distance} du point $A$ au plan $(P)$ est la longueur $AH$. On la note $d(A, (P)) = AH$.
C'est la \textbf{plus courte} distance entre $A$ et un point de $(P)$ : pour tout point $M$ de $(P)$ distinct de $H$, le triangle $AHM$ est rectangle en $H$, donc $AM > AH$ (l'hypoténuse est le plus grand côté).
\end{definition}

Cette notion est fondamentale en technique : la distance d'un point à un plan, c'est la \emph{hauteur} d'un solide, la \emph{profondeur} d'un perçage, le \emph{dégagement} d'une pièce au-dessus d'un établi.

\begin{popfigure}
\begin{center}
\begin{tikzpicture}[scale=1.0, line cap=round, line join=round]
% Plan (P)
\fill[popgreenL, opacity=0.5] (0,0) -- (4.8,0) -- (6.3,1.8) -- (1.5,1.8) -- cycle;
\draw[popgreen, thick] (0,0) -- (4.8,0) -- (6.3,1.8) -- (1.5,1.8) -- cycle;
\node[popgreen] at (5.8,1.45) {$(P)$};
% Point A hors du plan et projeté H
\coordinate (A) at (2.6,3.6);
\coordinate (H) at (3.4,0.9);
% Perpendiculaire (construction)
\draw[poporange, very thick] (A) -- (H);
\draw[poporange, dashed] (H) -- (3.65,-0.25);
% angle droit
\draw[popink] (H) ++(0.05,0.22) -- ++(0.2,0.03) -- ++(-0.05,-0.22);
\fill[popink] (A) circle (1.6pt); \node[popink, above] at (A) {$A$};
\fill[popink] (H) circle (1.6pt); \node[popink, below right] at (H) {$H$};
% Un point M quelconque et segment AM
\coordinate (M) at (4.9,0.75);
\fill[popdark] (M) circle (1.4pt); \node[popdark, right] at (M) {$M$};
\draw[popdark, dashed] (A) -- (M);
\draw[popdark] (H) -- (M);
\node[poporange, left] at (3.0,2.3) {\small $d(A,(P)) = AH$};
\end{tikzpicture}
\end{center}
\legende{Le projeté orthogonal $H$ de $A$ sur $(P)$ : la distance $AH$ est la plus courte distance de $A$ au plan, car $AM > AH$ pour tout autre point $M$ de $(P)$ (triangle $AHM$ rectangle en $H$).}
\end{popfigure}

\courssub{Exemple chiffré : volume d'un tétraèdre}

\begin{exoresolu}[Hauteur, orthogonalité et volume d'un tétraèdre]
\textbf{Énoncé.} Soit $SABC$ un tétraèdre tel que $(SA) \perp (ABC)$. On donne $SA = 5$~cm, $AB = 3$~cm, $AC = 4$~cm et $(AB) \perp (AC)$.
\begin{enumerate}
\item Justifier que $[SA]$ est la hauteur du tétraèdre relative à la base $ABC$.
\item Calculer l'aire de la base $ABC$.
\item En déduire le volume du tétraèdre.
\item Calculer $SB$.
\end{enumerate}
\tcblower
\textbf{Solution détaillée.}
\begin{enumerate}
\item Par hypothèse, la droite $(SA)$ est orthogonale au plan $(ABC)$ et le coupe en $A$. Le segment $[SA]$ joint donc le sommet $S$ à son projeté orthogonal $A$ sur le plan de base : $[SA]$ est bien la \cle{hauteur} du tétraèdre relative à la base $ABC$, et $h = SA = 5$~cm.
\item D'après la \textbf{propriété fondamentale}, $(SA) \perp (ABC)$ entraîne que $(SA)$ est perpendiculaire à toute droite du plan passant par $A$, en particulier $(SA) \perp (AB)$ et $(SA) \perp (AC)$. De plus, l'hypothèse $(AB) \perp (AC)$ signifie que le triangle $ABC$ est \textbf{rectangle en $A$}. Son aire est donc :
\[ \mathcal{A}_{ABC} = \frac{AB \times AC}{2} = \frac{3 \times 4}{2} = 6 \text{ cm}^2. \]
\item Le volume d'un tétraèdre (pyramide) est $V = \dfrac{1}{3} \times \mathcal{A}_{\text{base}} \times h$ :
\[ V = \frac{1}{3} \times 6 \times 5 = 10 \text{ cm}^3. \]
\item Comme $(SA) \perp (AB)$, le triangle $SAB$ est rectangle en $A$. Le théorème de Pythagore donne :
\[ SB^2 = SA^2 + AB^2 = 5^2 + 3^2 = 25 + 9 = 34, \quad \text{donc} \quad SB = \sqrt{34} \approx 5{,}83 \text{ cm}. \]
\end{enumerate}
\emph{Remarque de rédaction :} c'est la propriété fondamentale (étape 2) qui autorise à utiliser Pythagore dans $SAB$. Sans citer « $(SA) \perp (ABC)$ donc $(SA) \perp (AB)$ », le calcul de $SB$ serait injustifié.
\end{exoresolu}

\begin{exercice}[À l'atelier : vérifier la verticalité d'un montant]
Un montant métallique $[SH]$ est fixé sur une plaque horizontale $(P)$ en $H$. Pour vérifier qu'il est bien vertical, un technicien mesure : $SH = 1{,}20$~m ; pour deux points $M$ et $N$ de la plaque tels que $HM = 0{,}50$~m, $HN = 0{,}90$~m avec $(HM) \perp (HN)$, il trouve $SM = 1{,}30$~m et $SN = 1{,}50$~m.
\begin{enumerate}
\item Montrer, à l'aide de la réciproque du théorème de Pythagore, que $(SH) \perp (HM)$ et $(SH) \perp (HN)$.
\item Conclure que le montant est perpendiculaire à la plaque.
\end{enumerate}
\end{exercice}

\begin{corrige}[Exercice 1]
\begin{enumerate}
\item Dans le triangle $SHM$ : $SH^2 + HM^2 = 1{,}2^2 + 0{,}5^2 = 1{,}44 + 0{,}25 = 1{,}69 = 1{,}3^2 = SM^2$. D'après la réciproque du théorème de Pythagore, $SHM$ est rectangle en $H$, donc $(SH) \perp (HM)$.
Dans le triangle $SHN$ : $SH^2 + HN^2 = 1{,}44 + 0{,}81 = 2{,}25 = 1{,}5^2 = SN^2$, donc $SHN$ est rectangle en $H$ et $(SH) \perp (HN)$.
\item $(HM)$ et $(HN)$ sont deux droites \textbf{sécantes} du plan $(P)$ (elles se coupent en $H$ et ne sont pas parallèles puisque perpendiculaires). La droite $(SH)$ est orthogonale à ces deux droites : d'après le \textbf{critère de perpendicularité}, $(SH) \perp (P)$. Le montant est bien vertical.
\end{enumerate}
\end{corrige}

\begin{savaistu}[Les grues de chantier et la perpendiculaire au plan]
Les grues à tour que l'on voit sur les chantiers de Douala ou de Yaoundé appliquent directement le critère de cette leçon : le mât de la grue doit être orthogonale au plan du sol (ou à sa plateforme de calage). Les techniciens vérifient la verticalité dans \textbf{deux directions perpendiculaires} — exactement « deux droites sécantes du plan » — à l'aide d'un théodolite. Si le mât n'est perpendiculaire qu'à une seule direction du sol, il est penché et la grue risque de basculer sous la charge. La géométrie de cette leçon est donc, littéralement, une question de sécurité sur les chantiers.
\end{savaistu}

\begin{aretenir}[L'essentiel de la section]
\begin{itemize}
\item $(d) \perp (P)$ en $H$ signifie : $(d)$ orthogonale à \textbf{toute} droite de $(P)$ passant par $H$ ; $H$ est le \cle{pied de la perpendiculaire}.
\item \textbf{Critère} (pour prouver) : il suffit de \textbf{deux droites sécantes} de $(P)$ orthogonales à $(d)$.
\item \textbf{Propriété fondamentale} (pour utiliser) : si $(d) \perp (P)$, alors $(d)$ est perpendiculaire à toute droite de $(P)$ passant par $H$ — ce qui ouvre la voie au théorème de Pythagore.
\item La \cle{distance} d'un point $A$ à un plan $(P)$ est la longueur $AH$, où $H$ est le projeté orthogonal de $A$ sur $(P)$ ; c'est la plus courte distance de $A$ au plan.
\end{itemize}
\end{aretenir}

\coursec{Plans orthogonaux : Définition, Critère et Propriétés}

Dans la section précédente, nous avons appris à reconnaître qu'une droite est orthogonale à un plan : il suffit qu'elle soit orthogonale à deux droites sécantes de ce plan. Nous allons maintenant monter d'un niveau : au lieu de comparer une \cle{droite} et un \cle{plan}, nous allons comparer \textbf{deux plans entre eux}. Quand dit-on que deux murs, qu'une toiture et un pignon, qu'une façade et un sol sont \og perpendiculaires \fg{} ? C'est toute la question de l'\cle{orthogonalité de deux plans}, omniprésente dans les métiers du bâtiment, de la menuiserie et de la mécanique.

\courssub{Intuition et définition}

\textbf{Intuition.} Regardez le coin d'une salle de classe : le mur du fond et le mur latéral se coupent le long d'une verticale, et chacun semble \og à angle droit \fg{} par rapport à l'autre. De même, le couvercle d'une malle ouvert à angle droit est \og perpendiculaire \fg{} à la base de la malle. Comment traduire cela mathématiquement ? L'idée est de se ramener à ce que nous savons déjà faire : l'orthogonalité d'une \emph{droite} et d'un \emph{plan}.

\begin{definition}[Plans orthogonaux]
Deux plans sécants $(P)$ et $(Q)$ sont dits \cle{orthogonaux} (ou perpendiculaires) lorsque l'un d'eux contient une droite orthogonale à l'autre plan. On note alors :
\[ (P) \perp (Q). \]
\end{definition}

\begin{attention}[Sécants d'abord !]
La définition ne concerne que des plans \textbf{sécants}. Deux plans strictement parallèles ne sont jamais dits orthogonaux. Avant de parler d'orthogonalité, vérifiez toujours que les deux plans se coupent selon une droite $(D)$, appelée \cle{ligne d'intersection}.
\end{attention}

\textbf{Lien avec l'angle dièdre.} Lorsque $(P)$ et $(Q)$ se coupent selon $(D)$, prenez un point $A$ de $(D)$ et menez, dans chaque plan, la perpendiculaire à $(D)$ en $A$ : on obtient deux demi-droites qui forment un angle. Cet angle est l'\cle{angle dièdre} des deux plans. Dire que $(P)$ et $(Q)$ sont orthogonaux revient exactement à dire que cet angle dièdre est \textbf{droit} (il mesure $90^\circ$).

\begin{popfigure}
\begin{tikzpicture}[scale=1.05, line cap=round, line join=round]
  % Plan (Q) : horizontal (parallélogramme)
  \fill[popblueL, opacity=0.55] (-3.4,-1.2) -- (3.4,-1.2) -- (4.6,0.6) -- (-2.2,0.6) -- cycle;
  \draw[popblue] (-3.4,-1.2) -- (3.4,-1.2) -- (4.6,0.6) -- (-2.2,0.6) -- cycle;
  \node[popblue] at (3.9,-0.9) {$(Q)$};
  % Plan (P) : vertical (parallélogramme)
  \fill[poporangeL, opacity=0.6] (-2.6,-0.3) -- (2.6,-0.3) -- (2.6,2.9) -- (-2.6,2.9) -- cycle;
  \draw[poporange] (-2.6,-0.3) -- (2.6,-0.3) -- (2.6,2.9) -- (-2.6,2.9) -- cycle;
  \node[poporange] at (-2.1,2.5) {$(P)$};
  % Droite d'intersection (D)
  \draw[popdark, very thick] (-2.9,-0.3) -- (2.9,-0.3);
  \node[popdark] at (3.15,-0.3) {$(D)$};
  % Droite (d) dans (P), orthogonale à (D) et à (Q)
  \draw[popgreen, very thick] (0,-0.3) -- (0,2.6);
  \node[popgreen] at (0.25,2.45) {$(d)$};
  % Point A
  \fill[popdark] (0,-0.3) circle (1.6pt);
  \node[popdark, below left] at (0,-0.32) {$A$};
  % Angle droit (angle dièdre droit)
  \draw[poppurple, thick] (0,0.25) -- (0.55,0.25) -- (0.55,-0.3);
  \node[poppurple] at (1.35,0.35) {angle dièdre droit};
  % Droite dans (Q) perpendiculaire à (D) en A (trace de l'angle dièdre)
  \draw[poppurple, thick, dashed] (0,-0.3) -- (1.7,-1.0);
\end{tikzpicture}
\legende{Les plans $(P)$ et $(Q)$ se coupent selon $(D)$. La droite $(d)$, contenue dans $(P)$, est orthogonale à $(D)$ et au plan $(Q)$ : les deux plans sont orthogonaux, l'angle dièdre est droit.}
\end{popfigure}

\courssub{Le critère d'orthogonalité de deux plans}

Voici l'outil principal de cette leçon. C'est lui que vous utiliserez dans toutes les démonstrations du Probatoire et du Baccalauréat technique.

\begin{propriete}[Critère plan/plan (théorème)]
Si un plan $(P)$ contient une droite $(d)$ orthogonale à un plan $(Q)$, alors les plans $(P)$ et $(Q)$ sont orthogonaux :
\[ \left. \begin{array}{r} (d) \subset (P) \\ (d) \perp (Q) \end{array} \right\} \Longrightarrow (P) \perp (Q). \]
Cette démonstration est \textbf{formatrice} : sa rédaction guidée est exigible en devoir.
\end{propriete}

\begin{experience}[Démonstration guidée du critère]
\textbf{Hypothèses :} $(d) \subset (P)$ et $(d) \perp (Q)$. \textbf{Conclusion visée :} $(P) \perp (Q)$.
\begin{enumerate}
  \item Puisque $(d) \perp (Q)$, la droite $(d)$ coupe $(Q)$ en un point $A$ et $(d)$ est orthogonale à \emph{toute} droite de $(Q)$ passant par $A$.
  \item Les plans $(P)$ et $(Q)$ contiennent tous deux le point $A$ : ils sont donc sécants selon une droite $(D)$ passant par $A$.
  \item Comme $(D) \subset (Q)$ et $(d) \perp (Q)$, on a $(d) \perp (D)$.
  \item Dans $(Q)$, menons la droite $(\Delta)$ perpendiculaire à $(D)$ en $A$. Comme $(d) \perp (Q)$, on a aussi $(d) \perp (\Delta)$.
  \item L'angle dièdre de $(P)$ et $(Q)$ le long de $(D)$ est l'angle formé par $(d)$ et $(\Delta)$ : or cet angle est droit d'après l'étape 4.
  \item Donc les plans $(P)$ et $(Q)$ sont orthogonaux. \hfill $\square$
\end{enumerate}
\end{experience}

\begin{methode}[Prouver que deux plans sont orthogonaux]
Pour démontrer que $(P) \perp (Q)$ :
\begin{enumerate}
  \item \textbf{Repérer} une droite $(d)$ contenue dans $(P)$ (souvent une arête, une hauteur, une médiane de la figure).
  \item \textbf{Prouver} que $(d) \perp (Q)$ grâce au critère droite/plan : montrer que $(d)$ est orthogonale à \emph{deux droites sécantes} de $(Q)$.
  \item \textbf{Conclure} par le critère plan/plan : $(P)$ contient une droite orthogonale à $(Q)$, donc $(P) \perp (Q)$.
\end{enumerate}
Rédaction type : \og La droite $(d)$ est orthogonale aux droites $(u)$ et $(v)$, sécantes dans le plan $(Q)$, donc $(d) \perp (Q)$. Or $(d) \subset (P)$, donc $(P) \perp (Q)$. \fg
\end{methode}

\courssub{Propriété réciproque}

Le critère admet une réciproque, très utile pour \emph{construire} des droites orthogonales à un plan.

\begin{propriete}[Réciproque]
Si deux plans $(P)$ et $(Q)$ sont orthogonaux et se coupent selon la droite $(D)$, alors \textbf{toute droite de $(P)$ perpendiculaire à $(D)$ est orthogonale à $(Q)$}. En particulier, $(P)$ contient une \textbf{infinité} de droites orthogonales à $(Q)$ : toutes les perpendiculaires à $(D)$ menées dans $(P)$. Conséquence : tout plan passant par une droite orthogonale à $(Q)$ est orthogonal à $(Q)$.
\end{propriete}

\begin{experience}[Démonstration guidée de la réciproque]
\textbf{Hypothèses :} $(P) \perp (Q)$, $(P) \cap (Q) = (D)$, et $(d) \subset (P)$ avec $(d) \perp (D)$ en un point $A$. \textbf{Conclusion visée :} $(d) \perp (Q)$.
\begin{enumerate}
  \item Dans $(Q)$, menons la droite $(\Delta)$ perpendiculaire à $(D)$ en $A$.
  \item L'angle formé par $(d)$ et $(\Delta)$ est l'angle dièdre de $(P)$ et $(Q)$ : il est droit car $(P) \perp (Q)$. Donc $(d) \perp (\Delta)$.
  \item La droite $(d)$ est ainsi orthogonale à deux droites sécantes de $(Q)$ : $(D)$ (par hypothèse) et $(\Delta)$ (étape 2).
  \item Par le critère droite/plan, $(d) \perp (Q)$. \hfill $\square$
\end{enumerate}
\end{experience}

\begin{aretenir}[À retenir absolument]
\begin{itemize}
  \item \textbf{Critère :} une droite de $(P)$ orthogonale à $(Q)$ suffit pour conclure $(P) \perp (Q)$.
  \item \textbf{Réciproque :} si $(P) \perp (Q)$, alors dans $(P)$, toutes les perpendiculaires à la ligne d'intersection $(D)$ sont orthogonales à $(Q)$.
  \item Toute la démonstration passe par le \textbf{critère droite/plan} de la section 2 : orthogonalité à \emph{deux droites sécantes}.
\end{itemize}
\end{aretenir}

\courssub{Un piège majeur : la fausse transitivité}

\begin{attention}[Erreur fréquente : croire à la transitivité]
\textbf{Faux :} \og Si $(P) \perp (Q)$ et $(Q) \perp (R)$, alors $(P) \perp (R)$. \fg{} C'est une erreur classique au Probatoire ! Deux plans orthogonaux à un même troisième plan ne sont \textbf{pas nécessairement} orthogonaux entre eux : ils peuvent être \emph{parallèles} ou \emph{sécants quelconques}.

\textbf{Contre-exemple concret :} dans une chambre, le mur Est $(P)$ et le mur Ouest $(R)$ sont tous deux orthogonaux au sol $(Q)$... et pourtant ils sont \textbf{parallèles} entre eux ! De même, deux murs mitoyens $(P)$ et $(R)$ peuvent être tous deux perpendiculaires au sol $(Q)$ et se couper selon un angle quelconque.

\textbf{Bonne attitude :} ne jamais \og transiter \fg{} l'orthogonalité des plans. Chaque orthogonalité se démontre séparément, par le critère plan/plan.
\end{attention}

\courssub{Applications : pavé droit et prisme droit}

\begin{exemplebox}[Plans orthogonaux dans un pavé droit]
Soit $ABCDEFGH$ un pavé droit (parallélépipède rectangle) : $ABCD$ est la face du dessous, $EFGH$ la face du dessus, et les arêtes $(AE)$, $(BF)$, $(CG)$, $(DH)$ sont verticales. Montrons que les faces latérales $(ABFE)$ et $(BCGF)$ sont orthogonales.
\begin{itemize}
  \item L'arête $(AB)$ est contenue dans le plan $(ABFE)$.
  \item Dans la face $ABFE$ (rectangle), $(AB) \perp (BF)$.
  \item Dans la base $ABCD$ (rectangle), $(AB) \perp (BC)$.
  \item Les droites $(BF)$ et $(BC)$ sont sécantes en $B$ et appartiennent au plan $(BCGF)$.
  \item Donc $(AB) \perp (BCGF)$ (critère droite/plan).
  \item Comme $(AB) \subset (ABFE)$, le critère plan/plan donne : $(ABFE) \perp (BCGF)$.
\end{itemize}
\textbf{Conclusion métier :} deux faces latérales voisines d'un pavé droit sont toujours orthogonales — c'est exactement la situation de deux murs adjacents d'un bâtiment.
\end{exemplebox}

\begin{popfigure}
\begin{tikzpicture}[scale=1.15, line cap=round, line join=round]
  % Sommets du pavé
  \coordinate (A) at (0,0);
  \coordinate (B) at (3.4,0);
  \coordinate (C) at (4.5,1.3);
  \coordinate (D) at (1.1,1.3);
  \coordinate (E) at (0,2.6);
  \coordinate (F) at (3.4,2.6);
  \coordinate (G) at (4.5,3.9);
  \coordinate (H) at (1.1,3.9);
  % Face (BCGF) coloriée
  \fill[popblueL, opacity=0.5] (B) -- (C) -- (G) -- (F) -- cycle;
  % Face (ABFE) coloriée
  \fill[poporangeL, opacity=0.55] (A) -- (B) -- (F) -- (E) -- cycle;
  % Arêtes cachées
  \draw[popmuted, dashed] (A) -- (D);
  \draw[popmuted, dashed] (D) -- (C);
  \draw[popmuted, dashed] (D) -- (H);
  % Arêtes visibles
  \draw[popdark] (A) -- (B) -- (C);
  \draw[popdark] (A) -- (E);
  \draw[popdark] (B) -- (F);
  \draw[popdark] (C) -- (G);
  \draw[popdark] (E) -- (F) -- (G) -- (H) -- (E);
  % Arête (AB) mise en valeur : orthogonale au plan (BCGF)
  \draw[popgreen, very thick] (A) -- (B);
  % Angle droit en B entre (AB) et (BC)
  \draw[poppurple, thick] ($(B)+(-0.35,0)$) -- ($(B)+(-0.35,0.13)$) -- ($(B)+(0.13,0.15)$);
  % Noms des sommets
  \node[below left] at (A) {$A$};
  \node[below] at (B) {$B$};
  \node[right] at (C) {$C$};
  \node[above right] at (D) {$D$};
  \node[left] at (E) {$E$};
  \node[above left] at (F) {$F$};
  \node[right] at (G) {$G$};
  \node[above] at (H) {$H$};
  % Noms des plans
  \node[poporange] at (1.7,1.3) {$(ABFE)$};
  \node[popblue] at (4.15,2.6) {$(BCGF)$};
  \node[popgreen, below] at ($(A)!0.5!(B)$) {$(AB)$};
\end{tikzpicture}
\legende{Pavé droit $ABCDEFGH$ : les plans $(ABFE)$ et $(BCGF)$ sont orthogonaux. La droite $(AB)$, contenue dans $(ABFE)$, est orthogonale au plan $(BCGF)$ car orthogonale aux droites sécantes $(BF)$ et $(BC)$.}
\end{popfigure}

\begin{exoresolu}[Plan de symétrie orthogonal à une face et à la base dans une pyramide régulière]
\textbf{Énoncé.} Soit $SABCD$ une pyramide régulière : sa base $ABCD$ est un carré de centre $O$, de côté $a = 6$ cm, et le sommet $S$ est à la verticale de $O$ avec $SO = 8$ cm. On note $I$ le milieu de $[AB]$.
Montrer que le plan $(SOI)$ est orthogonal au plan $(SAB)$ et au plan $(ABCD)$.

\tcblower
\textbf{Solution détaillée.}

\textbf{1. Étude du plan $(SOI)$.}
Le plan $(SOI)$ est le plan de symétrie de la pyramide qui contient l'altitude $(SO)$ et la médiane $(SI)$ de la face $SAB$. Il coupe la base selon la droite $(OI)$, médiatrice de $[AB]$.

\textbf{2. Démonstration de $(SOI) \perp (ABCD)$.}
\begin{itemize}
  \item $(SO) \perp (ABCD)$ par définition de la pyramide régulière (hauteur).
  \item $(SO) \subset (SOI)$.
  \item Donc, par le critère plan/plan : $(SOI) \perp (ABCD)$.
\end{itemize}

\textbf{3. Démonstration de $(SOI) \perp (SAB)$.}
\begin{itemize}
  \item Dans la base carrée, $(OI)$ est la médiatrice de $[AB]$, donc $(OI) \perp (AB)$.
  \item Dans la face latérale $SAB$ (triangle isocèle en $S$), $(SI)$ est la médiane relative à $[AB]$, donc $(SI) \perp (AB)$.
  \item La droite $(AB)$ est ainsi orthogonale à deux droites sécantes du plan $(SOI)$ : $(OI)$ et $(SI)$, qui se coupent en $I$.
  \item Par le critère droite/plan : $(AB) \perp (SOI)$.
  \item Or $(AB) \subset (SAB)$.
  \item Par le critère plan/plan : $(SOI) \perp (SAB)$.
\end{itemize}

\textbf{Vérification numérique (culture du calcul) :} $OI = 3$ cm (moitié du côté), $SO = 8$ cm, donc $SI = \sqrt{SO^2 + OI^2} = \sqrt{64+9} = \sqrt{73} \approx 8{,}54$ cm : c'est l'apothème de la pyramide, utile pour les calculs d'aire latérale en atelier (tôlerie, couverture).
\end{exoresolu}

\begin{savaistu}[Orthogonalité et métiers du BTP au Cameroun]
Lors du montage des murs d'une maison, le maçon vérifie avec l'\textbf{équerre} et le \textbf{fil à plomb} que chaque mur est orthogonal à la dalle (droite verticale orthogonale au plan horizontal), puis que deux murs adjacents sont orthogonaux entre eux. Une erreur d'orthogonalité de quelques degrés se répercute sur toute la charpente ! En tôlerie et en chaudronnerie, le pliage d'une tôle crée exactement un \emph{angle dièdre} : quand ce pliage est à $90^\circ$, les deux faces de tôle sont des plans orthogonaux.
\end{savaistu}

\begin{exercice}[À faire]
Soit $ABCDEFGH$ un cube d'arête $5$ cm.
\begin{enumerate}
  \item Citer quatre paires de plans orthogonaux formés par des faces du cube.
  \item Montrer que le plan $(ACGE)$ (plan diagonal) est orthogonal au plan $(ABCD)$.
  \item Les plans $(ABFE)$ et $(DCGH)$ sont tous deux orthogonaux au plan $(ABCD)$. Sont-ils orthogonaux entre eux ? Justifier.
\end{enumerate}
\end{exercice}

\begin{corrige}[Exercice]
\textbf{1.} Par exemple : $(ABCD) \perp (ABFE)$ ; $(ABCD) \perp (BCGF)$ ; $(ABFE) \perp (BCGF)$ ; $(EFGH) \perp (ADHE)$. Toute paire de faces adjacentes convient.

\textbf{2.} La droite $(AE)$ est contenue dans $(ACGE)$. Or $(AE)$ est orthogonale à $(AB)$ et à $(AD)$ (faces carrées), deux droites sécantes du plan $(ABCD)$ : donc $(AE) \perp (ABCD)$. Comme $(AE) \subset (ACGE)$, le critère plan/plan donne $(ACGE) \perp (ABCD)$.

\textbf{3.} Non ! $(ABFE)$ et $(DCGH)$ sont deux faces \textbf{opposées} du cube : elles sont strictement parallèles. C'est l'illustration exacte du piège de la fausse transitivité : deux plans orthogonaux à un même plan peuvent être parallèles entre eux.
\end{corrige}

\begin{aretenir}[Bilan de la section]
\begin{itemize}
  \item $(P) \perp (Q)$ signifie : $(P)$ contient une droite orthogonale à $(Q)$ (angle dièdre droit).
  \item \textbf{Critère :} trouver dans l'un des plans une droite orthogonale à l'autre plan (via deux droites sécantes).
  \item \textbf{Réciproque :} si $(P) \perp (Q)$, toute droite de $(P)$ perpendiculaire à l'intersection $(D)$ est orthogonale à $(Q)$.
  \item \textbf{Piège :} l'orthogonalité des plans n'est \emph{pas transitive} : $(P) \perp (Q)$ et $(R) \perp (Q)$ n'impliquent rien entre $(P)$ et $(R)$.
\end{itemize}
\end{aretenir}

\coursec{Construction et Calculs : Applications techniques}

Dans la section précédente, nous avons appris à \emph{reconnaître} et à \emph{justifier} l'orthogonalité entre plans : un plan contenant une droite orthogonale à un autre plan lui est orthogonal. Mais dans l'atelier, sur le chantier ou au bureau d'études, il ne suffit pas de constater : il faut \textbf{construire} et \textbf{calculer}. Comment tracer sur une feuille la perpendiculaire à un plan ? Comment calculer la longueur d'un hauban de pylône, le volume d'une fondation pyramidale, la distance d'un point à un mur ? C'est l'objet de cette section : transformer nos connaissances théoriques en outils techniques opérationnels.

\courssub{La perspective cavalière : conventions de représentation}

Pour travailler dans l'espace sur une feuille plane, on utilise la \cle{perspective cavalière}. C'est le langage graphique commun du dessinateur technique et du mathématicien.

\begin{aretenir}[Conventions de la perspective cavalière]
\begin{itemize}
\item Un \cle{plan horizontal} (sol, table) est représenté par un parallélogramme : le parallélisme des droites qu'il contient est conservé sur le dessin.
\item Les \cle{verticales} sont représentées par des segments verticaux, à leur \textbf{vraie longueur}.
\item Les \cle{lignes de fuite} (direction de la profondeur) sont tracées avec un \textbf{angle de fuite} de $30^\circ$ ou $45^\circ$ par rapport à l'horizontale.
\item Les longueurs dans la direction de la profondeur sont multipliées par un \cle{coefficient de réduction}, en général $\frac{1}{2}$ ou $\frac{2}{3}$.
\item Les droites parallèles dans l'espace restent parallèles sur le dessin ; les milieux sont conservés.
\end{itemize}
\end{aretenir}

\begin{attention}[Ce que la perspective déforme]
En perspective cavalière, le parallélisme et les rapports de longueurs \emph{dans une même direction} sont conservés, mais \textbf{les angles ne le sont pas}. Un angle droit entre une verticale et une ligne de fuite apparaît sur le dessin comme un angle de $45^\circ$ ou $135^\circ$. Conséquences pratiques :
\begin{itemize}
\item on \textbf{ne mesure jamais} un angle sur une figure en perspective cavalière ;
\item on \textbf{ne mesure jamais} une longueur dans la direction de fuite sans tenir compte du coefficient de réduction ;
\item un angle droit se \emph{signale} par un codage et se \emph{justifie} par un théorème, jamais par la mesure.
\end{itemize}
\end{attention}

\courssub{Construction de la perpendiculaire à un plan par un point}

\textbf{Intuition.} Pour planter un piquet bien vertical dans un sol horizontal, le maçon utilise un fil à plomb : la verticale est la seule direction perpendiculaire à \emph{toutes} les directions horizontales. En géométrie, le problème est le même : par un point $A$, il passe une \textbf{unique} droite perpendiculaire à un plan $(P)$. Il s'agit de la construire.

Rappelons le critère fondamental vu dans la section 2 : une droite est orthogonale à un plan si et seulement si elle est orthogonale à \textbf{deux droites sécantes} de ce plan. Toute la construction repose sur cette idée.

\begin{methode}[Perpendiculaire à un plan $(P)$ passant par un point $A$ hors du plan]
\begin{enumerate}
\item Tracer dans le plan $(P)$ deux droites \textbf{sécantes} $\Delta_1$ et $\Delta_2$ (on les choisit de directions bien distinctes).
\item Construire le plan $(Q_1)$ passant par $A$ et perpendiculaire à $\Delta_1$ : il contient $A$ et toutes les droites passant par $A$ orthogonales à $\Delta_1$.
\item Construire de même le plan $(Q_2)$ passant par $A$ et perpendiculaire à $\Delta_2$.
\item L'intersection des plans $(Q_1)$ et $(Q_2)$ est une droite $(d)$ passant par $A$, orthogonale à la fois à $\Delta_1$ et à $\Delta_2$ : d'après le critère fondamental, c'est \textbf{la perpendiculaire cherchée}.
\item Le point d'intersection $H$ de $(d)$ avec $(P)$ est le \cle{projeté orthogonal} de $A$ sur $(P)$. La longueur $AH$ est la \cle{distance du point $A$ au plan $(P)$}.
\end{enumerate}
\end{methode}

\textbf{Cas où $A$ appartient au plan $(P)$.} La démarche est identique : on construit par $A$ une droite orthogonale à deux droites sécantes de $(P)$. En pratique, sur le dessin, si $(P)$ est horizontal, on trace par $A$ une droite de direction verticale.

\begin{attention}[Piège classique]
Une droite orthogonale à \textbf{une seule} droite du plan n'est pas nécessairement perpendiculaire au plan ! Il faut impérativement \textbf{deux droites sécantes} de $(P)$. C'est l'erreur la plus fréquente au Probatoire : écrire « $(d) \perp \Delta$ et $\Delta \subset (P)$ donc $(d) \perp (P)$ » est faux et coûte des points.
\end{attention}

\begin{popfigure}
\begin{tikzpicture}[scale=1.05, line join=round]
% Plan (P) : parallélogramme
\fill[popblueL, opacity=0.5] (0,0) -- (5.2,0) -- (6.4,1.8) -- (1.2,1.8) -- cycle;
\draw[popblue, thick] (0,0) -- (5.2,0) -- (6.4,1.8) -- (1.2,1.8) -- cycle;
\node[popblue] at (5.9,0.35) {$(P)$};
% Deux droites sécantes dans (P)
\draw[popdark] (0.7,0.35) -- (5.6,1.45) node[pos=0.9, above, sloped] {\footnotesize $\Delta_1$};
\draw[popdark] (1.3,1.6) -- (4.9,0.15) node[pos=0.85, below, sloped] {\footnotesize $\Delta_2$};
% Point A et perpendiculaire
\coordinate (A) at (3.1,4.2);
\coordinate (H) at (3.1,0.9);
\draw[poporange, very thick] (A) -- (H);
\fill[popink] (A) circle (1.6pt) node[above right] {$A$};
\fill[popink] (H) circle (1.6pt) node[below right] {$H$};
% Codage angle droit
\draw[poporange] (3.1,1.15) -- (3.35,1.15) -- (3.35,0.9);
% Plan auxiliaire suggéré
\draw[popgreen, dashed, thick] (2.2,3.4) -- (4.0,3.4) -- (3.55,0.62) -- (2.65,1.18) -- cycle;
\node[popgreen] at (1.95,3.1) {\footnotesize $(Q_1)$};
\end{tikzpicture}
\legende{La perpendiculaire $(AH)$ au plan $(P)$ est l'intersection de deux plans $(Q_1)$ et $(Q_2)$ passant par $A$ et perpendiculaires à deux droites sécantes $\Delta_1$ et $\Delta_2$ de $(P)$.}
\end{popfigure}

\courssub{Le théorème de Pythagore dans l'espace}

\textbf{Intuition.} Pour mesurer la diagonale d'une pièce rectangulaire, on applique Pythagore une fois au sol. Pour mesurer la \emph{grande diagonale} d'une caisse (d'un coin du sol au coin opposé du plafond), on applique Pythagore \textbf{deux fois} : d'abord dans le plan horizontal, puis dans un plan vertical. C'est le principe général des calculs de longueurs dans l'espace : \emph{une suite de deux Pythagore dans deux plans orthogonaux}.

\begin{propriete}[Pythagore dans le tétraèdre tri-rectangle]
Soit $OABC$ un tétraèdre \cle{tri-rectangle} en $O$, c'est-à-dire tel que les trois arêtes $OA$, $OB$, $OC$ soient deux à deux orthogonales. Alors :
\[ AB^2 = OA^2 + OB^2, \qquad AC^2 = OA^2 + OC^2, \qquad BC^2 = OB^2 + OC^2. \]
De plus, en notant $a = OA$, $b = OB$, $c = OC$, on a la relation :
\[ BC^2 + OA^2 = AC^2 + OB^2 = AB^2 + OC^2 = a^2 + b^2 + c^2. \]
En particulier, pour un pavé droit de dimensions $a$, $b$, $c$, la grande diagonale $d$ vérifie :
\[ d^2 = a^2 + b^2 + c^2. \]
\end{propriete}

\begin{experience}[Démonstration guidée : la grande diagonale du pavé droit]
Soit $ABCDEFGH$ un pavé droit, avec $AB = a$, $AD = b$, $AE = c$. On veut calculer la grande diagonale $AG$.
\begin{enumerate}
\item \textbf{Pythagore dans le plan horizontal.} Le triangle $ABC$ est rectangle en $B$ (car la base $ABCD$ est un rectangle) : $AC^2 = AB^2 + BC^2 = a^2 + b^2$.
\item \textbf{Justifier l'orthogonalité.} L'arête $CG$ est orthogonale au plan $(ABC)$ (propriété du pavé droit : les arêtes latérales sont perpendiculaires à la base), donc $CG$ est orthogonale à toute droite de ce plan, en particulier à $AC$.
\item \textbf{Pythagore dans le plan vertical.} Le triangle $ACG$ est donc rectangle en $C$ :
\[ AG^2 = AC^2 + CG^2 = (a^2 + b^2) + c^2. \]
\item \textbf{Conclusion.} $AG = \sqrt{a^2 + b^2 + c^2}$. La clé de la démonstration est l'étape 2 : c'est l'orthogonalité d'une droite au plan qui autorise le second Pythagore.
\end{enumerate}
\end{experience}

\begin{popfigure}
\begin{tikzpicture}[scale=1.5, line join=round]
% Tétraèdre tri-rectangle OABC
\coordinate (O) at (0,0);
\coordinate (A) at (0,2.6);
\coordinate (B) at (3.4,0);
\coordinate (C) at (1.1,1.5);
% Faces
\fill[popblueL, opacity=0.45] (O) -- (A) -- (C) -- cycle;
\fill[poporangeL, opacity=0.35] (O) -- (B) -- (C) -- cycle;
\fill[popgreenL, opacity=0.35] (O) -- (A) -- (B) -- cycle;
% Arêtes
\draw[popink, thick] (O) -- (A) node[midway, left] {$a$};
\draw[popink, thick] (O) -- (B) node[midway, below] {$b$};
\draw[popink, thick] (O) -- (C) node[midway, above left=-1pt] {\footnotesize $c$};
\draw[popink, thick] (A) -- (B);
\draw[popink, thick] (A) -- (C);
\draw[popink, thick] (B) -- (C);
% Codages angles droits en O
\draw[poporange, thick] (0,0.28) -- (0.2,0.28) -- (0.2,0);
\draw[poporange, thick] (0.24,0.02) -- (0.36,0.18) -- (0.16,0.26);
\draw[poporange, thick] (0.05,0.27) -- (0.16,0.42) -- (0.3,0.33);
% Sommets
\fill[popink] (O) circle (1.4pt) node[below left] {$O$};
\fill[popink] (A) circle (1.4pt) node[above] {$A$};
\fill[popink] (B) circle (1.4pt) node[right] {$B$};
\fill[popink] (C) circle (1.4pt) node[above right] {$C$};
\end{tikzpicture}
\legende{Tétraèdre tri-rectangle $OABC$ : les trois angles en $O$ sont droits (codés). Chaque face issue de $O$ est un triangle rectangle, ce qui permet trois applications du théorème de Pythagore.}
\end{popfigure}

\begin{exoresolu}[Tétraèdre tri-rectangle : calculs de longueurs]
Soit $OABC$ un tétraèdre tri-rectangle en $O$ avec $OA = 6$ cm, $OB = 8$ cm et $OC = 5$ cm. Calculer $BC$, $AC$ et $AB$.
\tcblower
\textbf{Solution.} Les arêtes $OB$ et $OC$ sont orthogonales, donc le triangle $OBC$ est rectangle en $O$. D'après le théorème de Pythagore :
\[ BC^2 = OB^2 + OC^2 = 8^2 + 5^2 = 64 + 25 = 89 \quad \text{donc} \quad BC = \sqrt{89} \approx 9{,}4 \text{ cm}. \]
De même :
\[ AC^2 = OA^2 + OC^2 = 36 + 25 = 61 \quad \text{donc} \quad AC = \sqrt{61} \approx 7{,}8 \text{ cm}, \]
\[ AB^2 = OA^2 + OB^2 = 36 + 64 = 100 \quad \text{donc} \quad AB = 10 \text{ cm}. \]
Vérification avec la relation de la propriété : $AB^2 + OC^2 = 100 + 25 = 125$ et $OA^2 + OB^2 + OC^2 = 36 + 64 + 25 = 125$. Les résultats sont cohérents.
\end{exoresolu}

\courssub{Volume du tétraèdre et de la pyramide}

\begin{propriete}[Volume d'une pyramide ou d'un tétraèdre]
Le volume d'une pyramide (et en particulier d'un tétraèdre) est donné par :
\[ V = \frac{1}{3} \times \mathcal{A}_{\text{base}} \times h, \]
où $h$ est la \cle{hauteur}, c'est-à-dire la \textbf{distance du sommet au plan de la base} (longueur de la perpendiculaire abaissée du sommet sur ce plan).
\end{propriete}

\begin{attention}[La hauteur n'est pas toujours une arête]
La hauteur est la \emph{perpendiculaire} au plan de la base, pas nécessairement une arête du solide. Dans un tétraèdre tri-rectangle en $O$, si l'on choisit la face $OBC$ comme base, alors $OA$ \emph{est} la hauteur (car $OA \perp (OBC)$, puisque $OA$ est orthogonale aux deux droites sécantes $OB$ et $OC$ de ce plan) : c'est le cas favorable. Mais si l'on choisit la face $ABC$ comme base, la hauteur est la distance de $O$ au plan $(ABC)$, qu'il faut calculer. Choisir astucieusement la base simplifie toujours les calculs.
\end{attention}

\begin{exemplebox}[Volume d'un tétraèdre tri-rectangle]
Avec le tétraèdre précédent ($OA = 6$, $OB = 8$, $OC = 5$ cm), prenons le triangle $OBC$ comme base. Ce triangle est rectangle en $O$, donc :
\[ \mathcal{A}_{OBC} = \frac{OB \times OC}{2} = \frac{8 \times 5}{2} = 20 \text{ cm}^2, \qquad V = \frac{1}{3} \times 20 \times 6 = 40 \text{ cm}^3. \]
Formule pratique pour le tétraèdre tri-rectangle : $V = \dfrac{OA \times OB \times OC}{6}$.
\end{exemplebox}

\courssub{Résolution de problèmes concrets}

\begin{methode}[Stratégie générale face à un problème concret]
\begin{enumerate}
\item \textbf{Modéliser} : identifier les points, les plans (sol, mur, plan vertical) et les longueurs connues.
\item \textbf{Justifier un angle droit} : utiliser une propriété d'orthogonalité (droite orthogonale à un plan, plans orthogonaux) pour exhiber un triangle rectangle.
\item \textbf{Appliquer Pythagore} dans ce triangle ; si nécessaire, recommencer dans un second plan orthogonal au premier.
\item \textbf{Conclure} avec la valeur exacte, puis une valeur approchée assortie de l'unité.
\end{enumerate}
\end{methode}

\begin{exoresolu}[Le pylône haubané]
Un pylône vertical de $20$ m est haubané par $4$ câbles fixés au sommet $S$ du pylône et ancrés au sol aux quatre sommets d'un carré de côté $15$ m, centré au pied $O$ du pylône. Calculer la longueur d'un câble.
\tcblower
\textbf{Solution.} Soit $M$ un point d'ancrage. Le pylône est vertical, donc $SO$ est orthogonale au plan du sol ; par conséquent $SO \perp OM$ : le triangle $SOM$ est rectangle en $O$.

\textbf{Étape 1 : distance au sol.} $OM$ est la moitié de la diagonale du carré de côté $15$ m. La diagonale du carré vaut $15\sqrt{2}$, donc :
\[ OM = \frac{15\sqrt{2}}{2} \approx 10{,}61 \text{ m}. \]

\textbf{Étape 2 : Pythagore dans le plan vertical $(SOM)$.}
\[ SM^2 = SO^2 + OM^2 = 20^2 + \left(\frac{15\sqrt{2}}{2}\right)^2 = 400 + \frac{225 \times 2}{4} = 400 + 112{,}5 = 512{,}5. \]
\[ SM = \sqrt{512{,}5} \approx 22{,}6 \text{ m}. \]
Chaque câble mesure environ $\SI{22,6}{m}$ ; il faudra donc environ $\SI{91}{m}$ de câble pour les quatre haubans (plus les marges de tension).
\end{exoresolu}

\begin{exercice}[Inclinaison d'un toit]
Un hangar a une base rectangulaire de $12$ m sur $8$ m. Le faîtage (arête supérieure du toit), parallèle aux murs de $12$ m, mesure $6$ m et est situé à $3$ m au-dessus du plan du sol, dans le plan vertical médiateur des murs de $8$ m. Calculer la longueur d'un rampant (distance d'un point du faîtage au mur de $8$ m, mesurée dans le plan vertical perpendiculaire au faîtage), puis la longueur d'un arêtier (segment joignant une extrémité du faîtage au coin du bâtiment).
\end{exercice}

\begin{corrige}[Exercice : inclinaison d'un toit]
\textbf{Rampant.} Dans le plan vertical perpendiculaire au faîtage, le rampant est l'hypoténuse d'un triangle rectangle de côtés $3$ m (hauteur) et $4$ m (demi-largeur $8/2$) :
\[ r^2 = 3^2 + 4^2 = 9 + 16 = 25 \quad \text{donc} \quad r = 5 \text{ m}. \]

\textbf{Arêtier.} On procède en deux Pythagore.

\emph{Étape 1 (plan horizontal).} Projetons l'extrémité du faîtage sur le sol : sa distance horizontale au coin du bâtiment vaut
\[ \sqrt{3^2 + 4^2} = \sqrt{25} = 5 \text{ m}, \]
car la demi-longueur restante vaut $(12-6)/2 = 3$ m et la demi-largeur vaut $4$ m.

\emph{Étape 2 (plan vertical).} La verticale du faîtage est orthogonale au plan du sol, donc perpendiculaire à cette distance horizontale. L'arêtier $\ell$ est l'hypoténuse d'un triangle rectangle de côtés $5$ m et $3$ m :
\[ \ell^2 = 5^2 + 3^2 = 25 + 9 = 34 \quad \text{donc} \quad \ell = \sqrt{34} \approx 5{,}83 \text{ m}. \]
On a bien utilisé deux Pythagore dans deux plans orthogonaux.
\end{corrige}

\begin{savaistu}[Orthogonalité et toitures du Cameroun]
Les charpentiers camerounais, de Foumban à Bafoussam, utilisent depuis des générations l'orthogonalité des \textbf{fermes} (structures triangulaires) par rapport au plan du sol pour garantir la stabilité des toitures. La ferme est un triangle placé dans un plan vertical, orthogonal au plan horizontal : c'est exactement la situation « plan orthogonal à un plan » étudiée dans la section 3 ! Le fil à plomb du charpentier matérialise la perpendiculaire au plan du sol, et l'équerre de maçon concrétise le critère des deux droites sécantes. Les mathématiques de cette leçon sont littéralement au-dessus de nos têtes.
\end{savaistu}

\begin{savaistu}[Le fil à plomb et l'équerre : des théorèmes entre les mains]
Le \cle{fil à plomb} du maçon matérialise la droite orthogonale au plan horizontal du sol : un mur monté au fil à plomb est un mur contenu dans un plan orthogonal au sol. L'\cle{équerre} vérifie l'orthogonalité de deux droites sécantes : en contrôlant deux directions du sol, le maçon s'assure, par le critère fondamental, que son poteau est bien perpendiculaire à tout le plan. Ces gestes professionnels sont des applications directes des théorèmes de cette leçon.
\end{savaistu}

\begin{aretenir}[L'essentiel de la section]
\begin{itemize}
\item En perspective cavalière : verticales vraies, profondeur réduite (coefficient $\frac{1}{2}$ ou $\frac{2}{3}$, angle de fuite $30^\circ$ ou $45^\circ$), et \textbf{jamais de mesure d'angle} sur le dessin.
\item La perpendiculaire à un plan par un point $A$ s'obtient comme intersection de deux plans passant par $A$, perpendiculaires à deux droites sécantes du plan ; son pied $H$ est le projeté orthogonal de $A$, et $AH$ est la distance de $A$ au plan.
\item Tout calcul de longueur dans l'espace se ramène à \textbf{une ou deux applications du théorème de Pythagore} dans des plans orthogonaux ; pour le pavé droit : $d^2 = a^2 + b^2 + c^2$.
\item Volume d'une pyramide : $V = \frac{1}{3} \times \mathcal{A}_{\text{base}} \times h$, où $h$ est la distance du sommet au plan de la base ; pour le tétraèdre tri-rectangle : $V = \dfrac{OA \times OB \times OC}{6}$.
\item Dans tout problème concret (pylône, toit, distance point/plan), la première étape est de \textbf{repérer le triangle rectangle} et de \textbf{justifier l'angle droit} par une propriété d'orthogonalité.
\end{itemize}
\end{aretenir}

\coursec{Synthèse et Résolution de problèmes complexes}

Dans les sections précédentes, nous avons appris à \emph{construire} et à \emph{calculer} avec l'orthogonalité : sections planes, distances, angles dans des solides concrets (poutres, murs, toitures). Il reste une compétence décisive pour l'épreuve du Probatoire technique : \textbf{enchaîner plusieurs raisonnements d'orthogonalité dans un même problème}, sans se perdre. C'est l'objet de cette synthèse : récapituler les trois grands critères, adopter une stratégie de résolution, puis s'entraîner sur des problèmes complets rédigés comme à l'examen.

\courssub{Récapitulatif : les trois grands critères d'orthogonalité}

Toute la leçon tient dans trois situations. Pour chacune, il faut connaître la \cle{définition} et le \cle{critère} (l'outil de démonstration).

\begin{aretenir}[Les trois critères fondamentaux]
\begin{enumerate}
\item \textbf{Droite $\perp$ Droite.} Deux droites $(d)$ et $(d')$ de l'espace sont \cle{orthogonales} si la parallèle à $(d')$ menée par un point de $(d)$ est perpendiculaire à $(d)$. En pratique : on se ramène à un angle droit \emph{dans un plan} qui contient les deux droites (ou une parallèle à l'une d'elles).
\item \textbf{Droite $\perp$ Plan.} Une droite $(d)$ est orthogonale à un plan $(P)$ si et seulement si $(d)$ est orthogonale à \textbf{deux droites sécantes} de $(P)$. Conséquence : $(d)$ est alors orthogonale à \emph{toute} droite de $(P)$.
\item \textbf{Plan $\perp$ Plan.} Deux plans $(P)$ et $(Q)$ sont perpendiculaires si et seulement si l'un d'eux contient une droite \textbf{orthogonale à l'autre}. Pour prouver $(P)\perp(Q)$, il suffit de trouver une droite $(d)\subset(P)$ telle que $(d)\perp(Q)$.
\end{enumerate}
\end{aretenir}

\begin{attention}[Le critère plan/plan n'est pas une équivalence automatique]
Piège classique : on trouve une droite $(d)$ orthogonale à $(Q)$, et on conclut $(P)\perp(Q)$ \emph{sans vérifier que $(d)$ est bien contenue dans $(P)$}. Le critère exige \textbf{deux} conditions : $(d)\subset(P)$ \emph{et} $(d)\perp(Q)$. Oublier la première condition invalide toute la démonstration. De même, pour le critère droite/plan, les deux droites du plan doivent être \textbf{sécantes} : deux droites parallèles du plan ne suffisent pas.
\end{attention}

\courssub{Stratégie globale : le « Détective de l'angle droit »}

Face à un problème d'orthogonalité, le bon réflexe n'est pas de chercher la solution d'un coup, mais de \emph{faire parler les hypothèses}.

\begin{methode}[Stratégie « Détective de l'angle droit »]
\begin{enumerate}
\item \textbf{Lister} toutes les orthogonalités données par l'énoncé (angles droits, hauteurs, diagonales perpendiculaires, arêtes d'un pavé\dots) et celles offertes par la nature du solide (losange, carré, pavé droit, pyramide régulière).
\item \textbf{Propager} : utiliser les propriétés du cours pour créer de nouvelles orthogonalités. Par exemple : si $(d)\perp(P)$, alors $(d)$ est orthogonale à toute droite de $(P)$ ; si deux plans sont perpendiculaires, une droite de l'un perpendiculaire à leur intersection est orthogonale à l'autre ; le parallélisme transporte l'orthogonalité.
\item \textbf{Cibler} l'orthogonalité demandée et choisir le bon critère : pour $(P)\perp(Q)$, chercher une droite de $(P)$ orthogonale à $(Q)$ ; pour $(d)\perp(P)$, chercher deux droites sécantes de $(P)$ orthogonales à $(d)$.
\item \textbf{Rédiger} la démonstration dans l'ordre logique, en citant chaque propriété utilisée, puis conclure clairement.
\end{enumerate}
\end{methode}

\begin{popfigure}
\begin{center}
\begin{tikzpicture}[
  font=\small,
  quest/.style={rectangle, rounded corners, draw=popdark, fill=popblueL, text width=4.2cm, align=center, minimum height=1cm},
  act/.style={rectangle, draw=popdark, fill=poporangeL, text width=3.6cm, align=center, minimum height=0.9cm},
  fin/.style={rectangle, rounded corners, draw=popdark, fill=popgreenL, text width=3.6cm, align=center, minimum height=0.9cm},
  ->, >=Stealth, thick]
\node[quest] (q) at (0,0) {Comment prouver que $(P)\perp(Q)$ ?};
\node[act] (a1) at (-4.2,-2.2) {Chercher une droite $(d)$ contenue dans $(P)$};
\node[act] (a2) at (0,-2.2) {Prouver $(d)\perp(Q)$ : $(d)$ orthogonale à deux droites sécantes de $(Q)$};
\node[act] (a3) at (4.2,-2.2) {Vérifier : $(d)\subset(P)$ \textbf{et} $(d)\perp(Q)$};
\node[fin] (c) at (0,-4.4) {Critère plan/plan : conclure $(P)\perp(Q)$};
\draw (q) -- (a1);
\draw (q) -- (a2);
\draw (q) -- (a3);
\draw (a1) -- (c);
\draw (a2) -- (c);
\draw (a3) -- (c);
\end{tikzpicture}
\end{center}
\legende{Arbre de décision : démarche pour prouver que deux plans sont perpendiculaires.}
\end{popfigure}

\courssub{Étude de cas complète : la pyramide à base losangique}

\begin{exoresolu}[Pyramide $SABCD$ à base losangique]
\textbf{Énoncé.} $SABCD$ est une pyramide de sommet $S$, dont la base $ABCD$ est un \textbf{losange} de centre $O$. On suppose de plus que $S$ se projette orthogonalement en $O$ sur le plan de la base, c'est-à-dire $(SO)\perp(ABCD)$. Démontrer que les plans $(SAC)$ et $(SBD)$ sont perpendiculaires.
\tcblower
\textbf{Solution détaillée (stratégie du détective).}

\emph{Étape 1 : lister les orthogonalités.}
\begin{itemize}
\item $ABCD$ est un losange : ses diagonales sont perpendiculaires, donc $(AC)\perp(BD)$ en $O$.
\item $(SO)\perp(ABCD)$ par hypothèse.
\end{itemize}

\emph{Étape 2 : propager.} Comme $(SO)\perp(ABCD)$, la droite $(SO)$ est orthogonale à \emph{toute} droite du plan $(ABCD)$, en particulier à $(BD)$ :
\[ (SO)\perp(BD). \]

\emph{Étape 3 : cibler.} On veut $(SAC)\perp(SBD)$. Cherchons une droite de $(SBD)$ orthogonale à $(SAC)$ : la droite $(BD)$ est candidate. Elle est orthogonale à $(AC)$ (diagonales du losange) et à $(SO)$ (étape 2). Or $(AC)$ et $(SO)$ sont deux droites \textbf{sécantes} du plan $(SAC)$ (elles se coupent en $O$). Donc, par le critère droite/plan :
\[ (BD)\perp(SAC). \]

\emph{Étape 4 : conclure.} La droite $(BD)$ est contenue dans le plan $(SBD)$ et est orthogonale au plan $(SAC)$. Par le critère plan/plan :
\[ \boxed{(SBD)\perp(SAC)}. \]
\end{exoresolu}

\begin{popfigure}
\begin{center}
\begin{tikzpicture}[scale=1.15, line join=round, font=\small]
% Base losange ABCD, centre O
\coordinate (O) at (0,0);
\coordinate (A) at (-2.6,0.2);
\coordinate (C) at (2.6,-0.2);
\coordinate (B) at (0.6,1.15);
\coordinate (D) at (-0.6,-1.15);
\coordinate (S) at (0,3.1);
% Faces arrière en pointillés
\draw[dashed, popmuted] (A)--(B)--(C);
\draw[dashed, popmuted] (A)--(D);
\draw[dashed, popmuted] (S)--(B);
\draw[dashed, popmuted] (A)--(C);
\draw[dashed, popmuted] (B)--(D);
% Arêtes visibles
\draw[thick, popdark] (D)--(C);
\draw[thick, popdark] (S)--(A);
\draw[thick, popdark] (S)--(C);
\draw[thick, popdark] (S)--(D);
% Hauteur SO
\draw[thick, poporange] (S)--(O);
% Angle droit en O entre SO et la base
\draw[poporange] (0,0) -- (0.18,0.02) -- (0.18,0.2) -- (0,0.18);
% Plans (SAC) et (SBD) suggérés
\fill[popblue, opacity=0.18] (S)--(A)--(C)--cycle;
\fill[popgreen, opacity=0.15] (S)--(B)--(D)--cycle;
% Points et labels
\foreach \p/\pos in {S/above, A/left, C/right, B/above right, D/below left, O/below right}
  \fill[popdark] (\p) circle (1.4pt) node[\pos] {$\p$};
\end{tikzpicture}
\end{center}
\legende{Pyramide $SABCD$ à base losangique : $(SO)\perp(ABCD)$ et $(AC)\perp(BD)$, donc $(BD)\perp(SAC)$, d'où $(SAC)\perp(SBD)$.}
\end{popfigure}

\courssub{Problème inverse : retrouver la nature d'un solide}

Jusqu'ici, on connaissait le solide et on prouvait des orthogonalités. Le \emph{problème inverse} consiste à \textbf{déduire la nature du solide à partir des orthogonalités données}. C'est un exercice très formateur, fréquent au Probatoire.

\begin{exemplebox}[Tétraèdre orthocentrique]
Soit $ABCD$ un tétraèdre tel que $(AB)\perp(CD)$ et $(AC)\perp(BD)$. Que peut-on dire ?
\begin{itemize}
\item On peut montrer que la troisième paire d'arêtes opposées est aussi orthogonale : $(AD)\perp(BC)$. Un tel tétraèdre est dit \cle{orthocentrique}.
\item Conséquence remarquable : les quatre hauteurs du tétraèdre sont concourantes (l'analogue spatial de l'orthocentre d'un triangle).
\item Réciproquement, si un tétraèdre a ses trois paires d'arêtes opposées orthogonales, alors chaque sommet se projette orthogonalement sur l'orthocentre de la face opposée.
\end{itemize}
\textbf{Idée de démonstration} (dans le cas où $A$ se projette en $H$, orthocentre de $BCD$) : $(AH)\perp(BCD)$ donc $(AH)\perp(CD)$ ; de plus $(BH)\perp(CD)$ car $H$ est l'orthocentre. Donc $(CD)$ est orthogonale à deux droites sécantes du plan $(ABH)$ : $(CD)\perp(ABH)$, et en particulier $(CD)\perp(AB)$. On retrouve une orthogonalité d'arêtes opposées.
\end{exemplebox}

\begin{savaistu}[Orthogonalité et infographie 3D]
En infographie 3D (jeux vidéo, conception assistée par ordinateur), chaque surface est munie de sa \cle{normale} : un vecteur orthogonal au plan tangent en chaque point. Ces normales sont essentielles pour calculer l'éclairage : dans le modèle d'ombrage de Phong, la luminosité d'un point dépend de l'angle entre la normale et la direction de la lumière. Sans orthogonalité, pas de reflets réalistes sur les carrosseries de voitures ni d'ombres crédibles dans les jeux vidéo !
\end{savaistu}

\courssub{Exemple type « Épreuve du Probatoire » rédigé}

L'épreuve attend une copie structurée : \textbf{énoncé lu et données repérées, figure codée, démonstration rédigée avec les propriétés citées, calculs soignés, conclusion encadrée}.

\begin{exoresolu}[Pavé droit : orthogonalité de plans et angle dièdre]
\textbf{Énoncé.} Soit $ABCDEFGH$ un pavé droit tel que $AB=6$, $BC=8$ et $AE=10$. La base est $ABCD$, le sommet $E$ est au-dessus de $A$ (arêtes verticales $AE$, $BF$, $CG$, $DH$). Soit $M$ le milieu de $[AE]$.
\begin{enumerate}
\item Prouver que le plan $(MCD)$ est perpendiculaire au plan $(ABFE)$.
\item Calculer l'angle entre le plan $(MCD)$ et le plan de base $(ABCD)$.
\end{enumerate}
\tcblower
\textbf{Solution rédigée comme à l'examen.}

\emph{1. Orthogonalité des plans $(MCD)$ et $(ABFE)$.}

Dans le pavé droit, la face latérale $(ABFE)$ est un rectangle. L'arête $(AD)$ est orthogonale à cette face : en effet $(AD)\perp(AB)$ (angle droit du rectangle $ABCD$) et $(AD)\perp(AE)$ (arête latérale perpendiculaire à la base), et $(AB)$ et $(AE)$ sont deux droites \textbf{sécantes} du plan $(ABFE)$. Donc, par le critère droite/plan :
\[ (AD)\perp(ABFE). \]

Montrons maintenant que $(AD)$ est contenue dans le plan $(MCD)$.
Le plan $(MCD)$ contient les points $M$, $C$, $D$. Il contient donc la droite $(DC)$ et la droite $(DM)$.
\begin{itemize}
\item $(DC)\parallel(AB)$ (pavé droit).
\item $\overrightarrow{DM} = \overrightarrow{DA} + \overrightarrow{AM}$. Or $\overrightarrow{AM}$ est collinéaire à $\overrightarrow{AE}$ (verticale) et $\overrightarrow{DA}$ est le vecteur de la droite $(AD)$.
\end{itemize}
Les vecteurs directeurs de $(MCD)$ sont $\overrightarrow{DC}$ (direction $(AB)$) et $\overrightarrow{DM}$. Comme $\overrightarrow{DA} = \overrightarrow{DM} - \overrightarrow{AM}$ et que $\overrightarrow{AM}$ est collinéaire à $\overrightarrow{DC}$ (car $(AE)\perp(ABCD)$ et $(DC)\subset(ABCD)$, mais plus simplement : $M\in(AE)$, $A\in(ABCD)$, $D\in(ABCD)$...), on voit que la direction de $(AD)$ est combinaison linéaire des deux directions du plan. Puisque le plan $(MCD)$ contient le point $D$, il contient la droite $(AD)$ passant par $D$ de direction $\overrightarrow{DA}$.
\emph{(Argument simplifié : le plan $(MCD)$ contient $D$ et la droite $(DC)\parallel(AB)$ ; il contient aussi $M$ et $D$ donc la droite $(MD)$. Dans le plan $(ADM)$, $(AD)$ et $(AM)$ sont deux droites sécantes. Or $(AM)\parallel(AE)\perp(ABCD)$ donc $(AM)\perp(DC)$. Le plan $(MCD)$ contient $(DC)$ et $(DM)$, donc il contient la droite $(AD)$ qui est dans le plan $(ADM)$ et passe par $D$.)}

Ainsi $(AD)\subset(MCD)$ et $(AD)\perp(ABFE)$. Par le critère plan/plan :
\[ \boxed{(MCD)\perp(ABFE)}. \]

\emph{2. Angle entre $(MCD)$ et $(ABCD)$.}

Les plans $(MCD)$ et $(ABCD)$ se coupent selon la droite $(CD)$ (car $C,D\in(ABCD)$ et $M\notin(ABCD)$ car $M$ est sur l'arête verticale $[AE]$).
L'angle dièdre entre ces deux plans se mesure dans un plan perpendiculaire à leur intersection $(CD)$.
\begin{itemize}
\item Dans le plan de base $(ABCD)$, la droite $(DA)$ est perpendiculaire à $(CD)$ (rectangle $ABCD$).
\item Dans le plan $(MCD)$, la droite $(DM)$ est perpendiculaire à $(CD)$ : en effet $(CD)\perp(AD)$ et $(CD)\perp(AE)$ (arête verticale perpendiculaire à la base), donc $(CD)$ est orthogonale au plan $(ADE)$ qui contient $(DM)$ ; par conséquent $(CD)\perp(DM)$.
\end{itemize}
L'angle cherché est donc l'angle $\widehat{ADM}$ (ou son supplémentaire, ici aigu).
Le triangle $ADM$ est rectangle en $A$ car $(AD)\subset(ABCD)$ et $(AM)\subset(AE)\perp(ABCD)$.
On a $AD = BC = 8$ et $AM = \frac{AE}{2} = 5$.
\[ \tan(\widehat{ADM}) = \frac{AM}{AD} = \frac{5}{8} = 0,625. \]
\[ \widehat{ADM} = \arctan(0,625) \approx 32,0^{\circ}. \]

\textbf{Conclusion :} $(MCD)\perp(ABFE)$ et l'angle entre $(MCD)$ et $(ABCD)$ vaut environ $32,0^{\circ}$.
\end{exoresolu}

\begin{attention}[Bien identifier l'intersection des deux plans]
Pour mesurer un angle entre deux plans (angle dièdre), il faut d'abord déterminer leur \textbf{droite d'intersection}, puis construire deux demi-droites perpendiculaires à cette intersection, une dans chaque plan. Une erreur fréquente est de mesurer l'angle entre deux droites quelconques des plans : le résultat est alors faux. Vérifiez toujours que les droites utilisées pour mesurer l'angle sont bien perpendiculaires à la ligne d'intersection.
\end{attention}

\courssub{Vers la Terminale : trigonométrie et produit scalaire}

Les notions de cette leçon seront approfondies en Terminale technique :
\begin{itemize}
\item \textbf{Angle droite/plan} : l'angle entre une droite et un plan est le complémentaire de l'angle entre la droite et la normale au plan ; il se calcule par trigonométrie dans un triangle rectangle.
\item \textbf{Angle dièdre} : mesuré comme ci-dessus, il intervient dans les calculs de toitures (pente des pans) et d'usinage (angles de coupe).
\item \textbf{Produit scalaire} : l'orthogonalité se généralise analytiquement par la condition $\vec{u}\cdot\vec{v}=0$. Muni d'un repère, on pourra alors prouver toutes les orthogonalités de cette leçon par le calcul : équations de plans, vecteurs normaux, cosinus directeurs.
\end{itemize}

\begin{savaistu}[Pourquoi le produit scalaire change tout]
Avec le produit scalaire, prouver que $(d)\perp(P)$ revient à vérifier que le vecteur directeur de $(d)$ est colinéaire au vecteur normal de $(P)$ : un simple calcul de coordonnées remplace une longue démonstration géométrique. C'est cette méthode analytique qui est utilisée dans les logiciels de CAO et de simulation employés dans les ateliers et bureaux d'études.
\end{savaistu}

\begin{exercice}[Pour s'entraîner]
Soit $SABC$ un tétraèdre tel que $(SA)\perp(ABC)$ et le triangle $ABC$ rectangle en $B$.
\begin{enumerate}
\item Démontrer que $(BC)\perp(SAB)$.
\item En déduire que $(SBC)\perp(SAB)$.
\item On donne $SA=5$, $AB=12$, $BC=6$. Calculer $SB$ puis l'aire du triangle $SBC$.
\end{enumerate}
\end{exercice}

\begin{corrige}[Exercice 1]
\begin{enumerate}
\item $(SA)\perp(ABC)$ donc $(SA)\perp(BC)$. De plus $(BC)\perp(AB)$ (triangle rectangle en $B$). Or $(SA)$ et $(AB)$ sont deux droites sécantes du plan $(SAB)$ : donc $(BC)\perp(SAB)$.
\item $(BC)\subset(SBC)$ et $(BC)\perp(SAB)$ : par le critère plan/plan, $(SBC)\perp(SAB)$.
\item Triangle $SAB$ rectangle en $A$ : $SB=\sqrt{SA^{2}+AB^{2}}=\sqrt{25+144}=\sqrt{169}=13$. Comme $(BC)\perp(SAB)$, on a $(BC)\perp(SB)$ : le triangle $SBC$ est rectangle en $B$, d'aire $\dfrac{SB\times BC}{2}=\dfrac{13\times6}{2}=39$ unités d'aire.
\end{enumerate}
\end{corrige}

\newpage

\coursec{Activité d'intégration}

\coursec{Leçon 19 : Orthogonalité dans l'espace}

\courssub{Objectifs de la leçon}
À l'issue de cette leçon, l'élève sera capable de :
\begin{itemize}
    \item définir et repérer des droites orthogonales dans l'espace (coplanaires ou non) ;
    \item énoncer et appliquer le \cle{critère fondamental} : une droite orthogonale à deux droites sécantes d'un plan est orthogonale à ce plan ;
    \item énoncer et appliquer la \cle{propriété fondamentale} : une droite orthogonale à un plan est orthogonale à toute droite de ce plan ;
    \item définir l'orthogonalité de deux plans (angle dièdre droit) et appliquer le critère : un plan contenant une droite orthogonale à un autre plan lui est orthogonal ;
    \item modéliser une situation technique réelle (charpente, abri, échafaudage) en utilisant ces propriétés pour justifier la rigidité et calculer des longueurs (théorème de Pythagore dans l'espace) ;
    \item rédiger un rapport technique structuré, argumenté et conforme aux attendus du Probatoire.
\end{itemize}

\courssub{Droites orthogonales dans l'espace}

\begin{definition}[Droites orthogonales]
Deux droites $(d)$ et $(d')$ de l'espace sont \cle{orthogonales} si l'on peut traduire l'une d'elles pour qu'elles deviennent perpendiculaires dans un plan. On note $(d) \perp (d')$.
\end{definition}

\begin{propriete}[Orthogonalité et parallélisme]
Si $(d) \perp (d')$ et $(d'') \parallel (d')$, alors $(d) \perp (d'')$.
\end{propriete}

\begin{attention}[Piège : droites non coplanaires]
Dans l'espace, deux droites peuvent être orthogonales \emph{sans se couper} (droites non coplanaires). L'orthogonalité ne suppose pas l'intersection, contrairement à la perpendicularité dans le plan.
\end{attention}

\courssub{Droite orthogonale à un plan : Critère et Propriété fondamentale}

\begin{definition}[Droite orthogonale à un plan]
Une droite $(d)$ est orthogonale à un plan $(P)$ si elle est orthogonale à \emph{toute} droite de $(P)$ passant par son pied (point d'intersection). On note $(d) \perp (P)$.
\end{definition}

\begin{propriete}[Critère fondamental de l'orthogonalité droite/plan] \label{critere-droite-plan}
Si une droite $(d)$ est orthogonale à \textbf{deux droites sécantes} d'un plan $(P)$, alors $(d)$ est orthogonale au plan $(P)$.
\begin{center}
\textit{Condition nécessaire et suffisante (au programme).}
\end{center}
\end{propriete}

\begin{propriete}[Propriété fondamentale] \label{prop-fondamentale}
Si $(d) \perp (P)$, alors $(d)$ est orthogonale à \textbf{toute} droite de $(P)$ passant par son pied.
\end{propriete}

\begin{propriete}[Droite orthogonale et plans parallèles]
Si $(d) \perp (P)$ et $(P) \parallel (Q)$, alors $(d) \perp (Q)$.
\end{propriete}

\begin{methode}[Démontrer $(d) \perp (P)$]
\begin{enumerate}
    \item Identifier le plan $(P)$ et le pied $H$ de la droite $(d)$ sur $(P)$ (ou un point commun).
    \item Trouver \textbf{deux droites} $(d_1)$ et $(d_2)$ dans $(P)$, \textbf{sécantes} (idéalement en $H$).
    \item Prouver $(d) \perp (d_1)$ et $(d) \perp (d_2)$ (souvent par construction : vertical/horizontal, ou propriétés du pavé droit).
    \item Conclure par le \textbf{critère fondamental} : $(d) \perp (P)$.
    \item Si besoin, utiliser la \textbf{propriété fondamentale} pour déduire $(d) \perp (d_3)$ pour toute autre droite $(d_3)$ de $(P)$.
\end{enumerate}
\end{methode}

\courssub{Plans orthogonaux : Définition, Critère et Propriétés}

\begin{definition}[Plans orthogonaux]
Deux plans $(P)$ et $(Q)$ sont \cle{orthogonaux} si l'angle dièdre qu'ils forment est droit. On note $(P) \perp (Q)$.
\end{definition}

\begin{propriete}[Critère des plans orthogonaux] \label{critere-plans}
Si un plan $(Q)$ contient une droite $(d)$ orthogonale à un plan $(P)$, alors $(Q) \perp (P)$.
\end{propriete}

\begin{propriete}[Propriétés des plans orthogonaux]
\begin{itemize}
    \item Si $(P) \perp (Q)$, leur intersection est une droite.
    \item Dans un plan $(Q) \perp (P)$, il existe une infinité de droites orthogonales à $(P)$ (toutes parallèles entre elles).
    \item Si $(P) \perp (Q)$ et $(Q) \parallel (R)$, alors $(P) \perp (R)$.
\end{itemize}
\end{propriete}

\begin{methode}[Démontrer $(P) \perp (Q)$]
\begin{enumerate}
    \item Identifier une droite $(d)$ contenue dans $(Q)$.
    \item Prouver que $(d) \perp (P)$ (souvent via le critère fondamental droite/plan).
    \item Conclure par le \textbf{critère des plans} : $(Q) \perp (P)$.
\end{enumerate}
\end{methode}

\courssub{Construction et Calculs : Applications techniques}

\begin{methode}[Calculer une longueur dans l'espace par orthogonalité]
\begin{enumerate}
    \item Modéliser la situation : repérer les points, droites, plans.
    \item Prouver qu'un triangle est \textbf{rectangle} (grâce à l'orthogonalité droite/plan ou plan/plan).
    \item Appliquer le \textbf{théorème de Pythagore} dans ce triangle rectangle.
    \item Exprimer le résultat avec l'unité adaptée (m, cm) et la précision demandée.
\end{enumerate}
\end{methode}

\begin{aretenir}[Boîte à outils : Orthogonalité et Calculs]
\begin{itemize}
    \item \textbf{Critère droite $\perp$ plan :} $(d) \perp (d_1)$ et $(d) \perp (d_2)$ avec $(d_1), (d_2) \subset (P)$ et $(d_1)$ sécante à $(d_2)$ $\implies (d) \perp (P)$.
    \item \textbf{Propriété fondamentale :} $(d) \perp (P) \implies (d) \perp$ toute droite de $(P)$ passant par le pied.
    \item \textbf{Critère plans $\perp$ :} $(d) \subset (Q)$ et $(d) \perp (P) \implies (Q) \perp (P)$.
    \item \textbf{Plans parallèles :} $(d) \perp (P)$ et $(P) \parallel (Q) \implies (d) \perp (Q)$.
    \item \textbf{Pythagore :} Dans $\triangle ABC$ rectangle en $A$, $BC^2 = AB^2 + AC^2$.
\end{itemize}
\end{aretenir}

\courssub{Synthèse : Activité d'intégration — Abri de marché modulaire à Mokolo}

\courssub{Situation}
\textbf{Conception d'un abri de marché modulaire à Mokolo (Yaoundé).}

La coopérative des commerçants du marché de Mokolo souhaite faire construire par un atelier de soudure local un abri métallique modulaire pour protéger les étals du soleil et de la pluie. Le soudeur dispose du plan suivant.

L'abri repose sur quatre \cle{poteaux} verticaux plantés aux quatre coins d'une dalle rectangulaire de \SI{4}{m} de long et \SI{3}{m} de large. Chaque poteau mesure \SI{2,5}{m} de hauteur. La \cle{toiture} est plane et horizontale : elle est soutenue par quatre \cle{poutres} reliant les sommets des poteaux deux à deux, formant un rectangle identique à la dalle.

Pour rigidifier l'ensemble, le soudeur prévoit des \cle{haubans} (câbles d'acier tendus) : chaque hauban relie le \emph{milieu} d'une poutre longue (celles de \SI{4}{m}) au \emph{pied du poteau d'extrémité de la poutre courte adjacente}, situé dans le même plan vertical que ce poteau et le milieu de la poutre longue.

Le chef de la coopérative, méfiant après l'effondrement d'un hangar voisin pendant une tornade, exige un \textbf{rapport technique} prouvant la stabilité de la structure avant de payer l'atelier.

\textbf{Données numériques :}
\begin{itemize}
    \item Dalle rectangulaire : longueur $AB = \SI{4}{m}$, largeur $AD = \SI{3}{m}$ ;
    \item Hauteur des poteaux : $AA' = \SI{2,5}{m}$ ;
    \item Hauban : du milieu $M$ de la poutre $A'B'$ au pied $D$ du poteau $D'$ (extrémité de la poutre courte $A'D'$). Le plan vertical contenant le hauban est le plan $(MID)$ où $I$ est le milieu de $[AB]$.
\end{itemize}

\begin{popfigure}
\begin{tikzpicture}[scale=0.7, every node/.style={font=\small}, >=Stealth]
% Coordonnées 3D projetées en cavalier (angle 45°, rapport 0.5)
% A(0,0,0), B(4,0,0), D(0,3,0), A'(0,0,2.5)
% Projection : x -> x, y -> y*cos45, z -> z + y*sin45 (approx cavalier standard)
% Simplification visuelle : x horizontal, z vertical, y en profondeur (cavalier)
\def\cx{1} % coeff x
\def\cy{0.5} % coeff y (profondeur)
\def\cz{1} % coeff z
\coordinate (A) at (0,0);
\coordinate (B) at (4*\cx, 0);
\coordinate (D) at (0*\cx + 3*\cy, 3*\cy);
\coordinate (C) at (4*\cx + 3*\cy, 3*\cy);
\coordinate (A') at (0, 2.5*\cz);
\coordinate (B') at (4*\cx, 2.5*\cz);
\coordinate (D') at (3*\cy, 3*\cy + 2.5*\cz);
\coordinate (C') at (4*\cx + 3*\cy, 3*\cy + 2.5*\cz);
\coordinate (M) at (2*\cx, 2.5*\cz); % Milieu A'B'
\coordinate (I) at (2*\cx, 0); % Milieu AB

% Faces cachées (pointillés)
\draw[thin, densely dashed, popmuted] (A) -- (D) -- (C) -- (B) -- cycle;
\draw[thin, densely dashed, popmuted] (A) -- (A');
\draw[thin, densely dashed, popmuted] (D) -- (D');
\draw[thin, densely dashed, popmuted] (C) -- (C');
\draw[thin, densely dashed, popmuted] (A') -- (D') -- (C') -- (B') -- cycle;
\draw[thin, densely dashed, popmuted] (B) -- (B');

% Faces visibles
\draw[thick, popdark] (A) -- (B) -- (B') -- (A') -- cycle; % Face avant
\draw[thick, popdark] (A) -- (D) -- (D') -- (A') -- cycle; % Face latérale gauche
\draw[thick, popdark] (A') -- (B') -- (C') -- (D') -- cycle; % Toiture

% Arêtes verticales visibles
\draw[thick, popdark] (B) -- (B');
\draw[thick, popdark] (D) -- (D');

% Hauban MD et plan vertical (MID)
\draw[very thick, poporange] (M) -- (D);
\draw[thick, popblue, densely dashed] (M) -- (I) -- (D); % Triangle MID
\draw[thin, popblue] (I) -- (M); % Vertical MI

% Points et labels
\foreach \p/\pos/\name in {A/below left/$A$, B/below right/$B$, D/left/$D$, C/right/$C$, A'/above left/$A'$, B'/above right/$B'$, D'/above left/$D'$, C'/above right/$C'$, M/above/$M$, I/below/$I$}
    \node[\pos] at (\p) {\name};

% Cotes
\draw[|<->|, thin, popmuted] (0,-0.5) -- (4*\cx, -0.5) node[midway, below] {4 m};
\draw[|<->|, thin, popmuted] (4*\cx + 0.2, 0) -- (4*\cx + 0.2 + 3*\cy, 3*\cy) node[midway, right] {3 m};
\draw[|<->|, thin, popmuted] (-0.5, 0) -- (-0.5, 2.5*\cz) node[midway, left] {2,5 m};

% Angle droit au sol
\draw[popgreen, thick] (A) ++(0.3,0) -- ++(0,0.3) -- ++(-0.3,0);
\draw[popgreen, thick] (A) ++(0,0.3) -- ++(0.3*0.707,0.3*0.707) -- ++(-0.3*0.707,0.3*0.707); % approx pour AD
% Angle droit MI / ID (dans le plan vertical)
% Projection de I(2,0), M(2,2.5), D(1.5, 1.5) approx.
% On marque l'angle droit en I sur le schéma 2D (plan MID)
\end{tikzpicture}
\legende{Modélisation en perspective cavalière de l'abri : pavé droit $ABCDA'B'C'D'$, hauban $[MD]$ dans le plan vertical $(MID)$ ($I$ milieu de $[AB]$).}
\end{popfigure}

\courssub{Consignes}

\textbf{Partie A — Modélisation.}
\begin{enumerate}
    \item Dessiner la structure en perspective cavalière (voir figure ci-dessus). Nommer $ABCD$ la dalle au sol et $A'B'C'D'$ le rectangle de la toiture, $A'$ étant le sommet du poteau de pied $A$, etc. Placer $I$ milieu de $[AB]$ et $M$ milieu de $[A'B']$.
    \item Placer un repère de l'espace $(A\,;\vec{AB},\vec{AD},\vec{AA'})$. Donner les coordonnées des points $A$, $B$, $C$, $D$, $A'$, $B'$, $C'$, $D'$, $I$ et $M$.
\end{enumerate}

\textbf{Partie B — Orthogonalité droite/plan.}
\begin{enumerate}
    \setcounter{enumi}{2}
    \item Justifier que le poteau $(AA')$ est orthogonal à deux droites sécantes du plan du sol $(ABC)$. Conclure.
    \item En déduire que $(AA')$ est orthogonal à toute droite du sol passant par $A$, en particulier à la diagonale $[AC]$.
    \item Les poteaux sont-ils orthogonaux au plan de la toiture $(A'B'C')$ ? Justifier en utilisant le parallélisme des plans $(ABC)$ et $(A'B'C')$.
\end{enumerate}

\textbf{Partie C — Orthogonalité plan/plan.}
\begin{enumerate}
    \setcounter{enumi}{5}
    \item On considère le plan $(AA'D')$ (façade latérale) contenant le poteau $(AA')$. Montrer que le plan $(AA'D')$ est orthogonal au plan du sol $(ABC)$.
    \item Le hauban $[MD]$ est contenu dans le plan vertical $(MID)$ où $I$ est le milieu de $[AB]$. Montrer que ce plan contient une droite orthogonale au plan de la toiture $(A'B'C')$, puis conclure que le plan $(MID)$ (plan du hauban) est orthogonal au plan de la toiture.
\end{enumerate}

\textbf{Partie D — Calculs.}
\begin{enumerate}
    \setcounter{enumi}{7}
    \item Calculer la longueur $A'M$ (moitié de la poutre longue).
    \item Justifier que le triangle $MID$ est rectangle en $I$, puis calculer la longueur du hauban $MD$ au centimètre près.
    \item Le câble d'acier coûte \num{1500} FCFA le mètre et l'atelier pose 4 haubans identiques. Calculer le coût total des câbles.
\end{enumerate}

\textbf{Partie E — Rapport technique.}
\begin{enumerate}
    \setcounter{enumi}{10}
    \item Rédiger en dix lignes maximum un rapport technique destiné au chef de la coopérative, justifiant la rigidité de l'abri par l'orthogonalité des plans porteurs (poteaux $\perp$ sol, plan des haubans $\perp$ toiture).
\end{enumerate}

\courssub{Ressources à mobiliser}

\begin{aretenir}[Boîte à outils de l'activité]
\begin{itemize}
    \item \textbf{Critère droite $\perp$ plan :} si une droite $(d)$ est orthogonale à deux droites \emph{sécantes} d'un plan $(P)$, alors $(d)$ est orthogonale au plan $(P)$.
    \item \textbf{Propriété fondamentale :} si $(d) \perp (P)$, alors $(d)$ est orthogonale à \emph{toute} droite de $(P)$ passant par son pied.
    \item \textbf{Critère plans $\perp$ :} si un plan $(Q)$ contient une droite orthogonale à un plan $(P)$, alors $(Q) \perp (P)$.
    \item \textbf{Plans parallèles :} une droite orthogonale à l'un de deux plans parallèles est orthogonale à l'autre.
    \item \textbf{Théorème de Pythagore} dans un triangle rectangle : $c^2 = a^2 + b^2$.
\end{itemize}
\end{aretenir}

\courssub{Démarche et corrigé commenté}

\begin{exoresolu}[Corrigé complet de l'activité d'intégration]
Résoudre les cinq parties de l'activité : modélisation, orthogonalités droite/plan, orthogonalité plan/plan, calcul du hauban, rapport technique.
\tcblower
\textbf{Partie A — Modélisation.}

La structure est un \textbf{pavé droit} $ABCDA'B'C'D'$.
\begin{itemize}
    \item Base $ABCD$ : rectangle, $AB = 4$, $AD = 3$.
    \item Arêtes latérales $AA', BB', CC', DD'$ : verticales, hauteur $2,5$.
    \item Toit $A'B'C'D'$ : rectangle congru à la base, horizontal.
    \item $I$ milieu de $[AB]$, $M$ milieu de $[A'B']$.
\end{itemize}

Le repère $(A\,;\vec{AB},\vec{AD},\vec{AA'})$ est un \textbf{repère orthogonal} (axes deux à deux perpendiculaires) car le pavé est droit. On prend comme unités : le mètre.

\begin{center}
\begin{tabularx}{\linewidth}{|c|X|X|X|X|X|}\hline
\rowcolor{popblueL} Point & $A$ & $B$ & $D$ & $C$ & $I$ \\ \hline
Coordonnées & $(0;0;0)$ & $(4;0;0)$ & $(0;3;0)$ & $(4;3;0)$ & $(2;0;0)$ \\ \hline
\rowcolor{popblueL} Point & $A'$ & $B'$ & $D'$ & $C'$ & $M$ \\ \hline
Coordonnées & $(0;0;2,5)$ & $(4;0;2,5)$ & $(0;3;2,5)$ & $(4;3;2,5)$ & $(2;0;2,5)$ \\ \hline
\end{tabularx}
\end{center}

\textbf{Partie B — Poteaux orthogonaux au sol et à la toiture.}

\begin{enumerate}[label=\arabic*., ref=\arabic*]
    \item Dans le plan du sol $(ABC)$, les droites $(AB)$ et $(AD)$ sont \textbf{sécantes} en $A$.
    Par construction du pavé droit, le poteau $(AA')$ est vertical, donc orthogonal à l'horizontale :
    \[ (AA') \perp (AB) \quad \text{et} \quad (AA') \perp (AD). \]
    D'après le \textbf{critère fondamental} (\ref{critere-droite-plan}) :
    \[ \boxed{(AA') \perp (ABC)}. \]
    \item Par la \textbf{propriété fondamentale} (\ref{prop-fondamentale}), puisque $(AA') \perp (ABC)$, la droite $(AA')$ est orthogonale à \emph{toute} droite du plan $(ABC)$ passant par $A$.
    En particulier, $(AA') \perp (AC)$. Le triangle $AA'C$ est rectangle en $A$.
    \item Le plan de la toiture $(A'B'C')$ est \textbf{parallèle} au plan du sol $(ABC)$ (translation selon le vecteur $\vec{AA'}$).
    Comme $(AA') \perp (ABC)$ et $(ABC) \parallel (A'B'C')$, la propriété des plans parallèles donne :
    \[ \boxed{(AA') \perp (A'B'C')}. \]
    Il en va de même pour $(BB')$, $(CC')$, $(DD')$.
\end{enumerate}

\textbf{Partie C — Plans orthogonaux.}

\begin{enumerate}[label=\arabic*., ref=\arabic*]
    \setcounter{enumi}{5}
    \item Le plan $(AA'D')$ (façade latérale) contient la droite $(AA')$.
    Nous venons de prouver $(AA') \perp (ABC)$.
    D'après le \textbf{critère des plans orthogonaux} (\ref{critere-plans}) :
    \[ \boxed{(AA'D') \perp (ABC)}. \]
    \item Soit $I$ le milieu de $[AB]$. $I(2;0;0)$. Le plan $(MID)$ contient les points $M(2;0;2,5)$, $I(2;0;0)$ et $D(0;3;0)$.
    \begin{itemize}
        \item La droite $(MI)$ est verticale (même abscisse et ordonnée pour $M$ et $I$), donc $(MI) \parallel (AA')$.
        \item Comme $(AA') \perp (A'B'C')$ (Partie B.3) et $(MI) \parallel (AA')$, on a $(MI) \perp (A'B'C')$.
        \item La droite $(MI)$ est contenue dans le plan $(MID)$.
    \end{itemize}
    D'après le \textbf{critère des plans orthogonaux} (\ref{critere-plans}), le plan $(MID)$ contenant une droite $(MI)$ orthogonale au plan $(A'B'C')$ est orthogonal à ce plan :
    \[ \boxed{(MID) \perp (A'B'C')}. \]
    Le hauban $[MD]$ appartient à ce plan, ce qui assure son ancrage rigide.
\end{enumerate}

\begin{popfigure}
\begin{tikzpicture}[scale=1.2, >=Stealth, every node/.style={font=\small}]
% Vue de dessus (plan horizontal) pour le calcul ID
\coordinate (A) at (0,0);
\coordinate (B) at (4,0);
\coordinate (D) at (0,3);
\coordinate (C) at (4,3);
\coordinate (I) at (2,0);
\coordinate (Mproj) at (2,0); % Projection de M

\draw[thick, popdark] (A) -- (B) -- (C) -- (D) -- cycle;
\draw[thick, poporange] (Mproj) -- (D);
\draw[thick, popblue, dashed] (I) -- (D);
\draw[thick, popblue, dashed] (Mproj) -- (I);

\foreach \p/\pos/\name in {A/below left/$A$, B/below right/$B$, D/above left/$D$, C/above right/$C$, I/below/$I$, Mproj/above right/$M$ (proj.)}
    \node[\pos] at (\p) {\name};

% Angle droit en I
\draw[popgreen, thick] (I) ++(0.2,0) -- ++(0,0.2) -- ++(-0.2,0);

% Cotes
\draw[|<->|, thin, popmuted] (0,-0.4) -- (4,-0.4) node[midway, below] {4 m};
\draw[|<->|, thin, popmuted] (-0.4,0) -- (-0.4,3) node[midway, left] {3 m};
\draw[|<->|, thin, popmuted] (2,-0.4) -- (2,0) node[midway, right] {2 m};
\draw[|<->|, thin, popmuted] (2,0) -- (0,3) node[midway, above, sloped] {$\sqrt{13}$ m};
\end{tikzpicture}
\legende{Vue de dessus (plan horizontal) : le triangle $MID$ vu en projection. $MI$ est vertical (sort du plan), $ID$ est horizontal. Le triangle est rectangle en $I$.}
\end{popfigure}

\textbf{Partie D — Longueur du hauban.}

\begin{enumerate}[label=\arabic*., ref=\arabic*]
    \setcounter{enumi}{7}
    \item $M$ est le milieu de $[A'B']$, et $A'B' = AB = \SI{4}{m}$.
    \[ A'M = \frac{A'B'}{2} = \frac{4}{2} = \boxed{\SI{2}{m}}. \]
    \item \textbf{Justification du triangle rectangle :}
    Dans le plan vertical $(MID)$ :
    \begin{itemize}
        \item $(MI)$ est verticale (poteau fictif au point $I$).
        \item $(ID)$ est horizontale (dans le plan du sol $(ABC)$).
    \end{itemize}
    Une droite verticale est orthogonale à toute droite horizontale : $(MI) \perp (ID)$.
    Donc le triangle $MID$ est \textbf{rectangle en $I$}.
    
    \textbf{Calcul :}
    \begin{align*}
        MI &= AA' = \SI{2,5}{m}. \\
        ID &= \text{distance entre } I(2;0;0) \text{ et } D(0;3;0) \\
           &= \sqrt{(2-0)^2 + (0-3)^2 + (0-0)^2} \\
           &= \sqrt{4 + 9} = \sqrt{13} \approx \SI{3,606}{m}.
    \end{align*}
    Dans $\triangle MID$ rectangle en $I$, d'après Pythagore :
    \begin{align*}
        MD^2 &= MI^2 + ID^2 \\
             &= 2,5^2 + (\sqrt{13})^2 \\
             &= 6,25 + 13 = 19,25. \\
        MD &= \sqrt{19,25} \approx \SI{4,387}{m} \approx \boxed{\SI{4,39}{m}} \text{ (au cm près)}.
    \end{align*}
    \item Longueur totale de câble : $4 \times 4,39 = \SI{17,56}{m}$.
    Coût total : $17,56 \times \num{1500} = \boxed{\num{26340} \text{ FCFA}}$.
\end{enumerate}

\textbf{Partie E — Rapport technique (modèle).}

\begin{quote}
\textbf{Rapport technique — Stabilité de l'abri modulaire Mokolo}
\vspace{0.2cm}

La structure est un pavé droit $ABCDA'B'C'D'$. Chaque poteau (ex. $AA'$) est orthogonal à deux droites sécantes de la dalle $(AB)$ et $(AD)$, donc orthogonal au plan du sol $(ABC)$ (critère fondamental). Il en résulte l'orthogonalité à toute diagonale (ex. $AC$), garantissant un transfert pur des charges verticales sans moment de flexion.

Le plan de la toiture $(A'B'C')$ est parallèle au sol, donc les poteaux lui sont aussi orthogonaux. Chaque hauban (ex. $[MD]$) appartient au plan vertical $(MID)$ qui contient la verticale $(MI)$. Ce plan est donc orthogonal à la toiture (critère plan/plan). Les quatre haubans triangulent les façades et empêchent le dévers latéral (vent). L'ensemble forme une structure hyperstatique rigide, conforme aux règles de l'art.
\end{quote}
\end{exoresolu}

\begin{attention}[Pièges fréquents au Probatoire]
\begin{itemize}
    \item \textbf{Critère mal appliqué :} Affirmer $(d) \perp (P)$ car $(d) \perp (d_1)$ avec $(d_1) \subset (P)$ \emph{une seule droite}. \textbf{Il faut deux droites sécantes.}
    \item \textbf{Ordre logique :} On ne peut pas utiliser la propriété fondamentale ($(d) \perp$ toute droite) \emph{avant} d'avoir prouvé $(d) \perp (P)$ par le critère.
    \item \textbf{Pythagore sans justification :} Calculer une longueur dans l'espace sans avoir prouvé au préalable que le triangle est rectangle (via l'orthogonalité).
    \item \textbf{Vocabulaire :} Confondre « plans parallèles » (sol et toiture) et « plans orthogonaux » (façade/sol, hauban/toiture).
    \item \textbf{Repérage :} Oublier de préciser que le repère est orthogonal (pavé droit) pour donner des coordonnées simples.
\end{itemize}
\end{attention}

\courssub{Bilan de la leçon}

Cette leçon a permis de structurer les savoirs sur l'orthogonalité dans l'espace autour de trois piliers :
\begin{itemize}
    \item \textbf{Droite/Plan :} Le couple Critère / Propriété fondamentale est le moteur de toute démonstration.
    \item \textbf{Plan/Plan :} Ramène toujours au cas Droite/Plan (un plan contient une droite).
    \item \textbf{Applications techniques :} L'orthogonalité n'est pas abstraite ; elle modélise la \cle{verticalité} (poteaux, haubans) et la \cle{rigidité} (triangulation, équerrage).
\end{itemize}
La rédaction d'un rapport technique (compétence « Communiquer ») exige de \textbf{nommer les objets} (droites, plans, points), de \textbf{citer le théorème} (critère, propriété) et de \textbf{conclure} explicitement.

\begin{savaistu}[Pour aller plus loin : Le fil à plomb, ancêtre du critère fondamental]
Les bâtisseurs de l'Antiquité (Égypte, Mésopotamie) et les artisans camerounais d'aujourd'hui vérifient la verticalité d'un poteau au \emph{fil à plomb} : un fil tendu par une masse (le plomb). Ce fil matérialise une droite orthogonale au plan horizontal (le sol). En pratique, on vérifie que le poteau est parallèle au fil : c'est l'application concrète de la propriété « droite orthogonale à un plan $\implies$ orthogonale à toutes les droites du plan ». Le critère fondamental (deux droites sécantes) est la traduction mathématique rigoureuse de ce geste millénaire.
\end{savaistu}

\newpage

\coursec{Méthodes et exercices résolus}

\coursec{Orthogonalité dans l'espace}

\courssub{Droites orthogonales dans l'espace}

\begin{definition}[Droites orthogonales dans l'espace]
Deux droites $(d)$ et $(d')$ de l'espace sont \cle{orthogonales} si l'on peut tracer, par un point quelconque, une parallèle à $(d')$ qui est \cle{perpendiculaire} à $(d)$.
\par
\textbf{Notation :} $(d) \perp (d')$.
\par
\textbf{Remarque :} Contrairement aux droites perpendiculaires (qui se coupent), deux droites orthogonales peuvent être \emph{non sécantes} (droites non coplanaires).
\end{definition}

\begin{propriete}[Critère d'orthogonalité de deux droites]
Soient $(d)$ et $(d')$ deux droites de l'espace. Si there existe une droite $(\Delta)$ telle que :
\begin{itemize}
\item $(\Delta) \parallel (d')$,
\item $(\Delta)$ est sécante à $(d)$,
\item $(d) \perp (\Delta)$ (perpendiculaires dans le plan qui les contient),
\end{itemize}
alors $(d) \perp (d')$.
\end{propriete}

\begin{methode}[Prouver que deux droites de l'espace sont orthogonales]
Pour montrer que deux droites $(d)$ et $(d')$ de l'espace sont orthogonales :
\begin{enumerate}
\item Chercher une droite $(\Delta)$ parallèle à $(d')$ et \emph{sécante} à $(d)$ (souvent $(\Delta)$ est confondue avec $(d')$ si les droites sont déjà sécantes).
\item Montrer que $(d)$ et $(\Delta)$ sont \cle{perpendiculaires} dans un plan qui les contient toutes les deux (utiliser : triangle rectangle, diagonales d'un carré, hauteur d'un triangle, Pythagore, propriété du pavé droit).
\item Conclure : $(d)$ et $(d')$ sont orthogonales car $(d)$ est perpendiculaire à une parallèle à $(d')$.
\end{enumerate}
\textbf{Réflexe technique :} dans un pavé droit, une arête est orthogonale à toute arête ou diagonale d'une face qui ne la coupe pas, dès lors qu'elle est perpendiculaire à une parallèle de cette arête dans la face.
\end{methode}

\begin{exemplebox}[Arêtes orthogonales dans un pavé droit]
Un coffret électrique a la forme d'un pavé droit $ABCDEFGH$ tel que $ABCD$ et $EFGH$ soient des rectangles horizontaux, $E$ au-dessus de $A$. On donne $AB = \SI{40}{cm}$, $AD = \SI{30}{cm}$, $AE = \SI{50}{cm}$.
\par
Démontrer que les droites $(AE)$ et $(BD)$ sont orthogonales.
\end{exemplebox}

\begin{exoresolu}[Arêtes orthogonales dans un pavé droit]
Un coffret électrique a la forme d'un pavé droit $ABCDEFGH$ tel que $ABCD$ et $EFGH$ soient des rectangles horizontaux, $E$ au-dessus de $A$. On donne $AB = \SI{40}{cm}$, $AD = \SI{30}{cm}$, $AE = \SI{50}{cm}$.\\
Démontrer que les droites $(AE)$ et $(BD)$ sont orthogonales.
\tcblower
\textbf{Étape 1 : chercher une droite sécante à $(AE)$ et parallèle à $(BD)$.}\\
La droite $(BD)$ est la diagonale de la face $ABCD$. La droite $(AE)$ est l'arête verticale issue de $A$.

\textbf{Étape 2 : se placer dans un plan contenant $(AE)$.}\\
Le solide est un pavé droit : l'arête $(AE)$ est perpendiculaire à la face $ABCD$. En particulier, $(AE)$ est perpendiculaire à toute droite de la face $ABCD$ passant par $A$.

Traçons dans la face $ABCD$ la parallèle à $(BD)$ passant par $A$ : appelons-la $(\Delta)$. Comme $(\Delta)$ est une droite du plan $(ABCD)$ passant par $A$, on a :
\[ (AE) \perp (\Delta). \]

\textbf{Étape 3 : conclure.}\\
Les droites $(AE)$ et $(\Delta)$ sont sécantes en $A$ et perpendiculaires, et $(\Delta)$ est parallèle à $(BD)$. Par définition des droites orthogonales dans l'espace :
\[ (AE) \text{ et } (BD) \text{ sont orthogonales.} \]

\textbf{Conclusion :} l'arête verticale $(AE)$ du coffret est orthogonale à la diagonale $(BD)$ du fond, bien que ces deux droites ne se coupent pas.
\end{exoresolu}

\courssub{Droite orthogonale à un plan : Critère et Propriété fondamentale}

\begin{definition}[Droite orthogonale à un plan]
Une droite $(d)$ est \cle{orthogonale} à un plan $(P)$ si elle est orthogonale à \emph{toute} droite de $(P)$ passant par son point d'intersection avec $(P)$ (le \cle{pied} de la perpendiculaire).
\par
On note $(d) \perp (P)$.
\end{definition}

\begin{propriete}[Critère fondamental : droite orthogonale à un plan]
Si une droite $(d)$ est perpendiculaire à \textbf{deux droites sécantes} $(d_1)$ et $(d_2)$ d'un plan $(P)$, alors $(d)$ est orthogonale au plan $(P)$.
\par
\textbf{Condition essentielle :} les deux droites $(d_1)$ et $(d_2)$ doivent être \cle{sécantes} (donc non parallèles).
\par
\textit{Démonstration exigible au Probatoire.}
\end{propriete}

\begin{propriete}[Propriété fondamentale de la droite orthogonale à un plan]
Si une droite $(d)$ est orthogonale à un plan $(P)$, alors elle est orthogonale à \emph{toute} droite de $(P)$, même celles qui ne la coupent pas.
\end{propriete}

\begin{methode}[Prouver qu'une droite est orthogonale à un plan]
Pour montrer qu'une droite $(d)$ est orthogonale à un plan $(P)$ :
\begin{enumerate}
\item Trouver dans $(P)$ \textbf{deux droites sécantes} $(d_1)$ et $(d_2)$ (attention : elles doivent se couper !).
\item Montrer que $(d) \perp (d_1)$ et que $(d) \perp (d_2)$ (en général $(d)$ coupe $(P)$ en un point, et on montre des perpendicularités dans des plans bien choisis).
\item Conclure par le \cle{critère fondamental} : une droite orthogonale à deux droites sécantes d'un plan est orthogonale à ce plan.
\end{enumerate}
\textbf{Conséquence utile :} si $(d) \perp (P)$, alors $(d)$ est orthogonale à \emph{toute} droite de $(P)$, même celles qui ne la coupent pas.
\end{methode}

\begin{exemplebox}[Pied d'une échelle et plan du sol]
Un pylône est modélisé par un segment $[SA]$ vertical, planté au point $A$ d'un sol horizontal plan. Au sol, deux câbles d'ancrage suivent les droites $(AB)$ et $(AC)$, sécantes en $A$, avec $AB = \SI{3}{m}$, $AC = \SI{4}{m}$ et $\widehat{BAC} = 90°$. Le sommet $S$ vérifie $SA = \SI{5}{m}$.
\end{exemplebox}

\begin{exoresolu}[Pied d'une échelle et plan du sol]
Un pylône est modélisé par un segment $[SA]$ vertical, planté au point $A$ d'un sol horizontal plan. Au sol, deux câbles d'ancrage suivent les droites $(AB)$ et $(AC)$, sécantes en $A$, avec $AB = \SI{3}{m}$, $AC = \SI{4}{m}$ et $\widehat{BAC} = 90°$. Le sommet $S$ vérifie $SA = \SI{5}{m}$.\\
1) Justifier que $(SA) \perp (AB)$ et $(SA) \perp (AC)$.\\
2) En déduire que $(SA)$ est orthogonale au plan du sol $(ABC)$.\\
3) Soit $M$ le milieu de $[BC]$. Montrer que $(SA)$ et $(AM)$ sont orthogonales, puis calculer $SM$.
\tcblower
\textbf{1)} Le pylône est vertical et le sol est horizontal : par construction, la verticale $[SA]$ est perpendiculaire en $A$ à toute direction horizontale passant par $A$, donc :
\[ (SA) \perp (AB) \quad \text{et} \quad (SA) \perp (AC). \]

\textbf{2)} Les droites $(AB)$ et $(AC)$ sont deux droites du plan $(ABC)$, \textbf{sécantes en $A$} (elles ne sont pas parallèles puisque $\widehat{BAC} = 90°$). La droite $(SA)$ est perpendiculaire à chacune d'elles. D'après le critère fondamental :
\[ (SA) \perp (ABC). \]

\textbf{3)} Puisque $(SA)$ est orthogonale au plan $(ABC)$, elle est orthogonale à \emph{toute} droite de ce plan, en particulier à $(AM)$ :
\[ (SA) \perp (AM). \]
Le triangle $SAM$ est donc rectangle en $A$. Calculons $AM$ : dans le triangle $ABC$ rectangle en $A$, le théorème de Pythagore donne
\[ BC^2 = AB^2 + AC^2 = 3^2 + 4^2 = 9 + 16 = 25 \quad \text{donc} \quad BC = \SI{5}{m}. \]
La médiane issue de l'angle droit vaut la moitié de l'hypoténuse :
\[ AM = \dfrac{BC}{2} = \dfrac{5}{2} = \SI{2,5}{m}. \]
Enfin, dans le triangle $SAM$ rectangle en $A$ :
\[ SM^2 = SA^2 + AM^2 = 5^2 + 2{,}5^2 = 25 + 6{,}25 = 31{,}25. \]
\[ SM = \sqrt{31{,}25} \approx \SI{5,59}{m}. \]

\textbf{Conclusion :} la distance du sommet du pylône au milieu de $[BC]$ est $SM = \sqrt{31{,}25} \approx \SI{5,59}{m}$.
\end{exoresolu}

\courssub{Plans orthogonaux : Définition, Critère et Propriétés}

\begin{definition}[Plans orthogonaux]
Deux plans $(P)$ et $(Q)$ sont \cle{orthogonaux} si l'un d'eux contient une droite orthogonale à l'autre plan.
\par
On note $(P) \perp (Q)$.
\end{definition}

\begin{propriete}[Critère d'orthogonalité de deux plans]
Si un plan $(P)$ contient une droite $(d)$ orthogonale à un plan $(Q)$, alors les plans $(P)$ et $(Q)$ sont orthogonaux.
\end{propriete}

\begin{propriete}[Propriétés de transmission de l'orthogonalité]
\begin{enumerate}
\item Si deux plans sont orthogonaux, toute droite de l'un \textbf{perpendiculaire à leur droite d'intersection} est orthogonale à l'autre plan.
\item Si $(d) \perp (P)$, alors tout plan contenant $(d)$ est orthogonal à $(P)$.
\item Deux droites orthogonales à un même plan sont \textbf{parallèles} entre elles.
\item Deux plans orthogonaux à une même droite sont \textbf{parallèles} entre eux.
\end{enumerate}
\end{propriete}

\begin{methode}[Prouver que deux plans sont orthogonaux]
Pour montrer que deux plans $(P)$ et $(Q)$ sont orthogonaux :
\begin{enumerate}
\item Trouver une droite $(d)$ contenue dans l'un des plans, par exemple $(d) \subset (P)$.
\item Montrer que $(d)$ est orthogonale à l'autre plan $(Q)$ (avec le critère : $(d)$ perpendiculaire à deux droites sécantes de $(Q)$).
\item Conclure : si un plan contient une droite orthogonale à un autre plan, alors les deux plans sont \cle{orthogonaux}.
\end{enumerate}
\textbf{Réflexe chantier :} un mur est perpendiculaire au sol parce qu'il contient une droite verticale (le fil à plomb), orthogonale au plan du sol.
\end{methode}

\begin{exemplebox}[Mur et dalle d'un atelier]
Sur un chantier, la dalle horizontale est le plan $(ABC)$ et le mur vertical est le plan $(ABD)$, sécants suivant l'arête $(AB)$. L'ouvrier vérifie avec un fil à plomb que l'arête $(AD)$ du mur est verticale, donc orthogonale au sol : $(AD) \perp (ABC)$.
\end{exemplebox}

\begin{exoresolu}[Mur et dalle d'un atelier]
Sur un chantier, la dalle horizontale est le plan $(ABC)$ et le mur vertical est le plan $(ABD)$, sécants suivant l'arête $(AB)$. L'ouvrier vérifie avec un fil à plomb que l'arête $(AD)$ du mur est verticale, donc orthogonale au sol : $(AD) \perp (ABC)$.\\
1) Montrer que les plans $(ABD)$ et $(ABC)$ sont orthogonaux.\\
2) On donne $AB = \SI{6}{m}$, $AD = \SI{2,5}{m}$ et $AC = \SI{4}{m}$ avec $(AC) \perp (AB)$. Calculer la longueur $DC$ du câble reliant le haut du mur $D$ au point $C$ de la dalle.
\tcblower
\textbf{1)} La droite $(AD)$ est contenue dans le plan $(ABD)$ (le mur) et, d'après la vérification au fil à plomb, $(AD)$ est orthogonale au plan $(ABC)$ (la dalle).\\
Or : \emph{si un plan contient une droite orthogonale à un autre plan, alors les deux plans sont orthogonaux.} Donc :
\[ (ABD) \perp (ABC). \]

\textbf{2)} Comme $(AD) \perp (ABC)$, la droite $(AD)$ est orthogonale à toute droite du plan $(ABC)$ passant par $A$, en particulier :
\[ (AD) \perp (AC). \]
Le triangle $ADC$ est donc \textbf{rectangle en $A$}. Le théorème de Pythagore donne :
\begin{align*}
DC^2 &= AD^2 + AC^2 \\
&= 2{,}5^2 + 4^2 \\
&= 6{,}25 + 16 = 22{,}25.
\end{align*}
\[ DC = \sqrt{22{,}25} \approx \SI{4,72}{m}. \]

\textbf{Conclusion :} le mur est bien orthogonal à la dalle, et le câble $[DC]$ mesure environ $\SI{4,72}{m}$.
\end{exoresolu}

\courssub{Construction et Calculs : Applications techniques}

\begin{methode}[Utiliser les propriétés de transmission de l'orthogonalité]
Réflexes à mobiliser :
\begin{enumerate}
\item Si deux plans sont orthogonaux, toute droite de l'un \textbf{perpendiculaire à leur droite d'intersection} est orthogonale à l'autre plan.
\item Si $(d) \perp (P)$, alors tout plan contenant $(d)$ est orthogonal à $(P)$.
\item Deux droites orthogonales à un même plan sont \textbf{parallèles} entre elles.
\item Deux plans orthogonaux à une même droite sont \textbf{parallèles} entre eux.
\end{enumerate}
\textbf{Réflexe rédaction :} citer la propriété utilisée, vérifier ses hypothèses (contenance, perpendicularité à l'intersection), puis conclure.
\end{methode}

\begin{exemplebox}[Problème de synthèse : la charpente]
Une charpente métallique repose sur un sol horizontal. Le poteau $[SH]$ est vertical, $H$ étant un point du sol. On donne $SH = \SI{3}{m}$. Au sol, $B$ et $C$ sont deux points tels que $HB = \SI{4}{m}$, $HC = \SI{4}{m}$ et $(HB) \perp (HC)$.
\end{exemplebox}

\begin{exoresolu}[Problème de synthèse : la charpente]
Une charpente métallique repose sur un sol horizontal. Le poteau $[SH]$ est vertical, $H$ étant un point du sol. On donne $SH = \SI{3}{m}$. Au sol, $B$ et $C$ sont deux points tels que $HB = \SI{4}{m}$, $HC = \SI{4}{m}$ et $(HB) \perp (HC)$.\\
1) Montrer que $(SH)$ est orthogonale au plan $(HBC)$.\\
2) Calculer $SB$ et $SC$.\\
3) Montrer que le triangle $SBC$ est isocèle en $S$.\\
4) Calculer l'aire du triangle $HBC$, puis celle du triangle $SBC$ (on admettra que la hauteur issue de $S$ dans $SBC$ mesure $\sqrt{SB^2 - \dfrac{BC^2}{4}}$).
\tcblower
\textbf{1)} Le poteau $[SH]$ est vertical et le sol est horizontal : $(SH)$ est perpendiculaire en $H$ à toute droite du sol passant par $H$, donc $(SH) \perp (HB)$ et $(SH) \perp (HC)$. Les droites $(HB)$ et $(HC)$ sont deux droites \textbf{sécantes en $H$} du plan $(HBC)$ (elles sont perpendiculaires, donc non parallèles). D'après le critère fondamental :
\[ (SH) \perp (HBC). \]

\textbf{2)} Puisque $(SH) \perp (HBC)$, on a $(SH) \perp (HB)$ : le triangle $SHB$ est rectangle en $H$. Pythagore :
\[ SB^2 = SH^2 + HB^2 = 3^2 + 4^2 = 25 \quad \text{donc} \quad SB = \SI{5}{m}. \]
De même, le triangle $SHC$ est rectangle en $H$ :
\[ SC^2 = SH^2 + HC^2 = 3^2 + 4^2 = 25 \quad \text{donc} \quad SC = \SI{5}{m}. \]

\textbf{3)} On a $SB = SC = \SI{5}{m}$ : le triangle $SBC$ est \textbf{isocèle en $S$}.

\textbf{4)} Le triangle $HBC$ est rectangle en $H$ :
\[ \mathcal{A}_{HBC} = \dfrac{HB \times HC}{2} = \dfrac{4 \times 4}{2} = \SI{8}{m^2}. \]
Dans $HBC$ rectangle en $H$ : $BC^2 = HB^2 + HC^2 = 16 + 16 = 32$, donc $BC = \sqrt{32} = 4\sqrt{2}$ m. La hauteur issue de $S$ dans le triangle isocèle $SBC$ mesure :
\[ h = \sqrt{SB^2 - \dfrac{BC^2}{4}} = \sqrt{25 - \dfrac{32}{4}} = \sqrt{25 - 8} = \sqrt{17} \approx \SI{4,12}{m}. \]
\[ \mathcal{A}_{SBC} = \dfrac{BC \times h}{2} = \dfrac{4\sqrt{2} \times \sqrt{17}}{2} = 2\sqrt{34} \approx \SI{11,66}{m^2}. \]

\textbf{Conclusion :} $SB = SC = \SI{5}{m}$, le triangle $SBC$ est isocèle en $S$, $\mathcal{A}_{HBC} = \SI{8}{m^2}$ et $\mathcal{A}_{SBC} = 2\sqrt{34} \approx \SI{11,66}{m^2}$.
\end{exoresolu}

\begin{exemplebox}[Deux poteaux verticaux]
Deux poteaux d'un hangar sont modélisés par les segments $[AE]$ et $[BF]$, tous deux verticaux, plantés au sol en $A$ et $B$. Le sol est le plan $(P)$.
\end{exemplebox}

\begin{exoresolu}[Deux poteaux verticaux]
Deux poteaux d'un hangar sont modélisés par les segments $[AE]$ et $[BF]$, tous deux verticaux, plantés au sol en $A$ et $B$. Le sol est le plan $(P)$.\\
1) Justifier que $(AE) \perp (P)$ et $(BF) \perp (P)$.\\
2) En déduire la position relative des droites $(AE)$ et $(BF)$.\\
3) Soit $(Q)$ le plan vertical contenant $(AE)$. Montrer que $(Q) \perp (P)$.
\tcblower
\textbf{1)} Chaque poteau est vertical et le sol est horizontal : un poteau vertical est orthogonal au plan du sol. Donc :
\[ (AE) \perp (P) \quad \text{et} \quad (BF) \perp (P). \]

\textbf{2)} Propriété : \emph{deux droites orthogonales à un même plan sont parallèles entre elles.} Comme $(AE)$ et $(BF)$ sont toutes deux orthogonales au plan $(P)$ :
\[ (AE) \parallel (BF). \]
Les deux poteaux sont donc parallèles, ce qui garantit la régularité de la charpente.

\textbf{3)} Le plan $(Q)$ contient la droite $(AE)$, et $(AE)$ est orthogonale au plan $(P)$. Propriété : \emph{si un plan contient une droite orthogonale à un autre plan, alors les deux plans sont orthogonaux.} Donc :
\[ (Q) \perp (P). \]

\textbf{Conclusion :} les poteaux sont parallèles entre eux, et tout plan vertical contenant un poteau est orthogonal au sol.
\end{exoresolu}

\courssub{Synthèse et Résolution de problèmes complexes}

\begin{methode}[Rédiger et justifier une démarche complète]
Pour rédiger une solution de géométrie dans l'espace au Probatoire :
\begin{enumerate}
\item \textbf{Nommer} précisément les objets : droites $(AB)$, plans $(ABC)$, points d'intersection.
\item \textbf{Annoncer le théorème utilisé} avant de l'appliquer : « d'après le critère d'orthogonalité d'une droite à un plan... ».
\item \textbf{Vérifier les hypothèses} explicitement : les deux droites du plan sont-elles bien sécantes ? Le point d'appartenance est-il justifié ?
\item \textbf{Conclure par une phrase} complète qui répond à la question posée, avec les unités pour les grandeurs.
\item Pour les calculs de longueurs : repérer d'abord le \emph{triangle rectangle} (justifier l'angle droit par l'orthogonalité démontrée), \emph{puis} appliquer Pythagore.
\end{enumerate}
\end{methode}

\begin{attention}[Les 5 erreurs qui coûtent des points au Probatoire]
\begin{enumerate}
\item \textbf{Oublier la condition « deux droites sécantes »} dans le critère d'orthogonalité d'une droite à un plan. Une droite perpendiculaire à deux droites \emph{parallèles} d'un plan n'est pas nécessairement orthogonale à ce plan. Écris toujours : « $(d_1)$ et $(d_2)$ sont deux droites \textbf{sécantes} du plan $(P)$ ».
\item \textbf{Confondre perpendiculaires et orthogonales} : deux droites orthogonales de l'espace ne se coupent pas forcément. Ne rédige jamais « $(AE)$ et $(BD)$ se coupent à angle droit » si elles ne sont pas sécantes.
\item \textbf{Appliquer Pythagore sans justifier l'angle droit} : avant d'écrire $SB^2 = SH^2 + HB^2$, il faut avoir démontré (ou rappelé) que le triangle est rectangle, en citant l'orthogonalité utilisée.
\item \textbf{Oublier les unités et la conclusion} : un résultat numérique sans unité ($\SI{}{m}$, $\SI{}{m^2}$) et sans phrase de conclusion fait perdre des points de présentation.
\item \textbf{Mal citer les plans} : un plan est désigné par trois points non alignés, par exemple $(ABC)$. Écrire « le plan $(AB)$ » n'a aucun sens et est sanctionné.
\end{enumerate}
\end{attention}

\begin{savaistu}[L'orthogonalité dans les métiers techniques]
En construction métallique, en charpente et en électricité, l'orthogonalité garantit la stabilité et la répartition des charges. Un poteau vertical (orthogonal au sol) transmet les efforts de compression directement au sol sans moment fléchissant. Les câbles d'ancrage (haubans) forment des triangles rectangles avec le mât : connaître l'orthogonalité permet de calculer les longueurs de câbles et les angles de tension. Au Probatoire, les sujets techniques (pylônes, hangars, coffrets, charpentes) sont des classiques : maîtrisez le vocabulaire (pied, arête, face, diagonale, plan horizontal/vertical) et les raisonnements types (critère, propriété fondamentale, transmission).
\end{savaistu}

\begin{exercice}[Entraînement : Orthogonalité dans un tétraèdre rectangle]
Soit $SABC$ un tétraèdre tel que $(SA) \perp (SB)$, $(SA) \perp (SC)$ et $(SB) \perp (SC)$. On donne $SA = \SI{6}{cm}$, $SB = \SI{8}{cm}$, $SC = \SI{10}{cm}$.
\begin{enumerate}
\item Montrer que $(SA) \perp (SBC)$.
\item Calculer $BC$, $AC$, $AB$.
\item Calculer le volume du tétraèdre $SABC$.
\item Soit $H$ le projeté orthogonal de $S$ sur le plan $(ABC)$. Exprimer le volume en fonction de $SH$ et de l'aire de $ABC$, puis en déduire $SH$.
\end{enumerate}
\end{exercice}

\begin{corrige}[Exercice n°1]
\textbf{1)} Dans le plan $(SBC)$, les droites $(SB)$ et $(SC)$ sont sécantes en $S$. On a $(SA) \perp (SB)$ et $(SA) \perp (SC)$ par hypothèse. D'après le critère fondamental : $(SA) \perp (SBC)$.

\textbf{2)} Triangles rectangles en $S$ :
\begin{itemize}
\item Dans $SBC$ : $BC^2 = SB^2 + SC^2 = 8^2 + 10^2 = 164 \Rightarrow BC = \sqrt{164} = 2\sqrt{41}$ cm.
\item Dans $SAC$ : $AC^2 = SA^2 + SC^2 = 6^2 + 10^2 = 136 \Rightarrow AC = \sqrt{136} = 2\sqrt{34}$ cm.
\item Dans $SAB$ : $AB^2 = SA^2 + SB^2 = 6^2 + 8^2 = 100 \Rightarrow AB = 10$ cm.
\end{itemize}

\textbf{3)} Volume $V = \dfrac{1}{3} \times \mathcal{A}_{SBC} \times SA$. $\mathcal{A}_{SBC} = \dfrac{SB \times SC}{2} = \dfrac{8 \times 10}{2} = 40$ cm$^2$.
$V = \dfrac{1}{3} \times 40 \times 6 = 80$ cm$^3$.

\textbf{4)} $V = \dfrac{1}{3} \times \mathcal{A}_{ABC} \times SH$. Calculons $\mathcal{A}_{ABC}$ par Héron ou en notant que $AB=10$, $AC=2\sqrt{34}$, $BC=2\sqrt{41}$.
Demi-périmètre $p = \frac{10 + 2\sqrt{34} + 2\sqrt{41}}{2} = 5 + \sqrt{34} + \sqrt{41}$.
Aire au carré : $\mathcal{A}^2 = p(p-AB)(p-AC)(p-BC)$. Calcul fastidieux.
Alternative : $\mathcal{A}_{ABC}^2 = \mathcal{A}_{SAB}^2 + \mathcal{A}_{SBC}^2 + \mathcal{A}_{SAC}^2$ (théorème de la face opposée au trièdre rectangle).
$\mathcal{A}_{SAB} = 24$, $\mathcal{A}_{SBC} = 40$, $\mathcal{A}_{SAC} = 30$.
$\mathcal{A}_{ABC}^2 = 24^2 + 40^2 + 30^2 = 576 + 1600 + 900 = 3076$.
$\mathcal{A}_{ABC} = \sqrt{3076} = 2\sqrt{769}$ cm$^2$.
$80 = \dfrac{1}{3} \times 2\sqrt{769} \times SH \Rightarrow SH = \dfrac{120}{\sqrt{769}} \approx \SI{4,33}{cm}$.
\end{corrige}

\newpage

\coursec{Exercices}

\courssub{Je vérifie mes connaissances}

\begin{exercice}[QCM : Orthogonalité droite/plan]
Soit un plan $\mathcal{P}$ et une droite $\Delta$ orthogonale à $\mathcal{P}$. Parmi les affirmations suivantes, laquelle est \textbf{fausse} ?
\begin{enumerate}
    \item Toute droite de $\mathcal{P}$ passant par le pied de $\Delta$ est orthogonale à $\Delta$.
    \item Toute droite parallèle à $\Delta$ est orthogonale à $\mathcal{P}$.
    \item Toute droite orthogonale à $\Delta$ est contenue dans $\mathcal{P}$.
    \item Le plan passant par $\Delta$ et orthogonal à une droite de $\mathcal{P}$ est orthogonal à $\mathcal{P}$.
\end{enumerate}
\end{exercice}

\begin{exercice}[Vrai / Faux Justifié : Plans orthogonaux]
On considère deux plans $\mathcal{P}$ et $\mathcal{Q}$ sécants selon la droite $\Delta$. Pour chaque affirmation, dire si elle est vraie ou fausse et justifier brièvement.
\begin{enumerate}
    \item Si une droite de $\mathcal{P}$ est orthogonale à $\mathcal{Q}$, alors $\mathcal{P} \perp \mathcal{Q}$.
    \item Si $\mathcal{P} \perp \mathcal{Q}$, alors toute droite de $\mathcal{P}$ est orthogonale à $\mathcal{Q}$.
    \item Si une droite $\Delta_1$ de $\mathcal{P}$ est orthogonale à $\Delta$ (l'intersection), alors $\mathcal{P} \perp \mathcal{Q}$.
    \item Si $\mathcal{P} \perp \mathcal{Q}$, il existe au moins une droite de $\mathcal{P}$ orthogonale à $\mathcal{Q}$.
\end{enumerate}
\end{exercice}

\begin{exercice}[Définitions et Critères : Rappels directs]
\begin{enumerate}
    \item Donner la définition : « Une droite est orthogonale à un plan si... »
    \item Énoncer le \textbf{critère} pour qu'une droite soit orthogonale à un plan.
    \item Énoncer la \textbf{propriété fondamentale} de la droite orthogonale à un plan.
    \item Donner la définition de deux plans orthogonaux.
    \item Énoncer le \textbf{critère} pour que deux plans sécants soient orthogonaux.
\end{enumerate}
\end{exercice}

\begin{exercice}[Calculs directs : Angle droite/plan et distance]
Dans un repère orthonormé $(\mathcal{O}; \vec{i}, \vec{j}, \vec{k})$, on donne le plan $\mathcal{P}$ d'équation $2x - y + 2z = 5$ et la droite $\Delta$ de système paramétrique $\begin{cases} x = 1 + 2t \\ y = -1 - t \\ z = 3 + 2t \end{cases}$.
\begin{enumerate}
    \item Montrer que $\Delta \perp \mathcal{P}$.
    \item Calculer la distance du point $A(1;-1;3)$ (sur $\Delta$) au plan $\mathcal{P}$.
    \item Déterminer les coordonnées du projeté orthogonal $H$ de $A$ sur $\mathcal{P}$.
\end{enumerate}
\end{exercice}

\courssub{Je m'entraîne}

\begin{exercice}[Application : Pyramide à base carrée]
Soit $SABCD$ une pyramide à base carrée de centre $O$, telle que $AB = 6$ cm et la hauteur $SO = 8$ cm.
\begin{enumerate}
    \item Montrer que la droite $(SO)$ est orthogonale au plan $(ABCD)$.
    \item Calculer la longueur $SA$.
    \item Déterminer l'angle entre la droite $(SA)$ et le plan $(ABCD)$ (arrondi au degré près).
    \item $M$ est le milieu de $[SA]$. Le plan $(MBC)$ coupe $(SD)$ en $N$. Déterminer la nature du quadrilatère $MBCN$ et montrer que le plan $(MBC)$ est orthogonal au plan $(SAD)$.
\end{enumerate}
\end{exercice}

\begin{exercice}[Application : Tétraèdre rectangle]
Soit $SABC$ un tétraèdre tel que la face $ABC$ est un triangle rectangle en $A$ ($AB=5$ cm, $AC=12$ cm) et la droite $(SA)$ est orthogonale au plan $(ABC)$ avec $SA = 4$ cm.
\begin{enumerate}
    \item Calculer le volume $V$ du tétraèdre $SABC$.
    \item Calculer l'aire de la face $SBC$.
    \item En déduire la distance du sommet $S$ au plan $(SBC)$ (hauteur relative à la base $SBC$).
    \item Calculer l'angle entre les plans $(SAB)$ et $(ABC)$.
\end{enumerate}
\end{exercice}

\begin{exercice}[Contexte technique : Escalier droit]
Un escalier droit est modélisé dans un repère. Les marches ont un giron (largeur) de $28$ cm et une hauteur de $17$ cm. La cage d'escalier est un parallélépipède rectangle. On considère la ligne de pente (limon) qui relie le nez de la première marche au nez de la dernière marche.
\begin{enumerate}
    \item Dans un repère adapté, donner les coordonnées d'un vecteur directeur $\vec{u}$ de la ligne de pente pour une marche.
    \item Calculer l'angle $\alpha$ que fait la ligne de pente avec le plan horizontal (le palier). Arrondir au degré près.
    \item Si l'escalier comporte $15$ marches, quelle est la longueur totale du limon (arrondi au cm) ?
    \item La norme impose un angle de pente compris entre $25^\circ$ et $42^\circ$. Cet escalier est-il conforme ?
\end{enumerate}
\end{exercice}

\begin{exercice}[Démonstration : Cube et orthogonalité]
Soit $ABCDEFGH$ un cube de côté $a$. $I$ est le milieu de $[CG]$.
\begin{enumerate}
    \item Montrer que la droite $(AI)$ est orthogonale au plan $(BDH)$.
    \item En déduire que $(AI) \perp (BD)$.
    \item Calculer la distance de $A$ au plan $(BDH)$ en fonction de $a$.
    \item Déterminer l'angle entre les plans $(AIC)$ et $(BDH)$.
\end{enumerate}
\end{exercice}

\courssub{Je me prépare au Probatoire}

\begin{exercice}[Sujet type Probatoire : Pyramide régulière et plans orthogonaux]
\textbf{Énoncé :} Soit $SABCD$ une pyramide régulière à base carrée de centre $O$. On a $AB = 6$ cm et la hauteur $SO = 8$ cm. $M$ est le milieu de $[SA]$. Le plan $\mathcal{P} = (MBC)$ coupe $(SD)$ en $N$ et $(SO)$ en $K$.
\begin{enumerate}
    \item Montrer que $(SO) \perp (ABCD)$. En déduire que $(SO) \perp (BC)$.
    \item Montrer que $(BC) \perp (SAD)$. En déduire que $(BC) \perp (MN)$.
    \item Montrer que $(BC) \perp \mathcal{P}$.
    \item Montrer que le plan $\mathcal{P}$ est orthogonal au plan $(SAD)$.
    \item Calculer le volume de la pyramide $S-MBCN$.
    \item Déterminer l'angle entre les plans $\mathcal{P}$ et $(ABCD)$ (arrondi au degré près).
\end{enumerate}
\end{exercice}

\begin{exercice}[Sujet type Probatoire : Tétraèdre et distances]
\textbf{Énoncé :} On considère un tétraèdre $ABCD$ tel que :
\begin{itemize}
    \item Le triangle $BCD$ est équilatéral de côté $6$ cm.
    \item La droite $(AD)$ est orthogonale au plan $(BCD)$ et $AD = 4$ cm.
\end{itemize}
$I$ est le centre de gravité du triangle $BCD$. $M$ est le milieu de $[AB]$.
\begin{enumerate}
    \item Montrer que $(DI) \perp (BC)$. En déduire que $(BC) \perp (ADI)$.
    \item Montrer que le plan $(ADI)$ est orthogonal au plan $(ABC)$.
    \item Calculer le volume du tétraèdre $ABCD$.
    \item Calculer la distance de $A$ au plan $(BCD)$.
    \item Calculer l'angle entre les droites $(AB)$ et $(CD)$.
    \item Le plan parallèle à $(BCD)$ passant par $M$ coupe $(AC)$ en $N$ et $(AD)$ en $P$. Calculer le volume du tétraèdre $A-MNP$.
\end{enumerate}
\end{exercice}

\begin{exercice}[Sujet type Probatoire : Parallélépipède rectangle et orthogonalité]
\textbf{Énoncé :} Dans un parallélépipède rectangle $ABCDEFGH$, on a $AB = 12$ cm, $BC = 5$ cm, $AE = 4$ cm. $I$ est le milieu de $[EH]$, $J$ le milieu de $[FG]$. On considère le plan $\mathcal{Q} = (AIJ)$.
\begin{enumerate}
    \item Montrer que $(BC) \perp (AEH)$. En déduire $(BC) \perp \mathcal{Q}$.
    \item Montrer que $\mathcal{Q} \perp (ABCD)$.
    \item Le plan $\mathcal{Q}$ coupe $(DH)$ en $K$ et $(CG)$ en $L$. Déterminer la nature du quadrilatère $AIJL$ et calculer son aire.
    \item Calculer l'angle entre les plans $\mathcal{Q}$ et $(ABFE)$.
    \item Calculer la distance du point $C$ au plan $\mathcal{Q}$.
\end{enumerate}
\end{exercice}

\newpage

\coursec{Corrigés des exercices}

\coursec{Orthogonalité dans l'espace — Corrigés détaillés}

\courssub{Exercices de révision : QCM, Vrai/Faux et Rappels de cours}

\begin{corrige}[Exercice 1 — QCM : Orthogonalité droite/plan]
L'affirmation \textbf{fausse} est la \textbf{c)}.

\begin{itemize}
\item \textbf{a) Vrai.} Par définition, une droite orthogonale à un plan est orthogonale à \emph{toute} droite de ce plan ; en particulier à toute droite du plan passant par son pied (droites orthogonales \emph{sécantes}, donc perpendiculaires).
\item \textbf{b) Vrai.} Si deux droites sont parallèles, tout plan orthogonal à l'une est orthogonal à l'autre (propriété de l'orthogonalité droite/plan).
\item \textbf{c) Faux.} Une droite orthogonale à $\Delta$ (où $\Delta \perp \mathcal{P}$) est seulement \emph{parallèle au plan} $\mathcal{P}$ (ou contenue dans $\mathcal{P}$). Contre-exemple : dans un repère, si $\Delta$ est l'axe vertical et $\mathcal{P}$ le plan horizontal, une droite horizontale située à \SI{5}{m} au-dessus de $\mathcal{P}$ est orthogonale à $\Delta$ sans être contenue dans $\mathcal{P}$.
\item \textbf{d) Vrai.} C'est le critère d'orthogonalité de deux plans : un plan contenant une droite orthogonale à $\mathcal{P}$ est orthogonal à $\mathcal{P}$.
\end{itemize}
\end{corrige}

\begin{corrige}[Exercice 2 — Vrai / Faux : Plans orthogonaux]
\begin{enumerate}
\item \textbf{Vrai.} C'est exactement le \cle{critère d'orthogonalité} de deux plans : si un plan $\mathcal{P}$ contient une droite orthogonale à $\mathcal{Q}$, alors $\mathcal{P} \perp \mathcal{Q}$.
\item \textbf{Faux.} Seules les droites de $\mathcal{P}$ \emph{orthogonales à la droite d'intersection} $\Delta$ sont orthogonales à $\mathcal{Q}$. Contre-exemple : $\Delta$ elle-même est contenue dans $\mathcal{P}$ et n'est pas orthogonale à $\mathcal{Q}$ (elle y est contenue).
\item \textbf{Faux.} Être orthogonale à la seule droite $\Delta$ ne suffit pas : il faudrait que $\Delta_1$ soit orthogonale à \emph{deux droites sécantes} de $\mathcal{Q}$ (c'est-à-dire orthogonale au plan $\mathcal{Q}$). Une droite de $\mathcal{P}$ peut être perpendiculaire à $\Delta$ sans être orthogonale à $\mathcal{Q}$ si les plans ne sont pas orthogonaux.
\item \textbf{Vrai.} Si $\mathcal{P} \perp \mathcal{Q}$, alors toute droite de $\mathcal{P}$ orthogonale à l'intersection $\Delta$ est orthogonale à $\mathcal{Q}$ (propriété des plans orthogonaux). Il en existe donc au moins une (en fait une infinité, toutes parallèles entre elles).
\end{enumerate}
\end{corrige}

\begin{corrige}[Exercice 3 — Définitions et critères : rappels directs]
\begin{enumerate}
\item \textbf{Définition.} Une droite est orthogonale à un plan si elle est orthogonale à \cle{toute droite} de ce plan (de manière équivalente : à toute droite du plan passant par son point d'intersection avec le plan).
\item \textbf{Critère.} Si une droite est orthogonale à \cle{deux droites sécantes} d'un plan, alors elle est orthogonale à ce plan.
\item \textbf{Propriété fondamentale.} Si une droite $\Delta$ est orthogonale à un plan $\mathcal{P}$ en un point $H$, alors $\Delta$ est orthogonale à toute droite de $\mathcal{P}$ ; de plus $H$ est le \cle{projeté orthogonal} sur $\mathcal{P}$ de tout point de $\Delta$, et la distance d'un point $M$ de $\Delta$ au plan $\mathcal{P}$ est $MH$.
\item \textbf{Définition.} Deux plans sécants sont orthogonaux lorsque l'angle entre eux (angle diédral) vaut $90^\circ$ : les droites de chacun des plans, perpendiculaires en un même point à la droite d'intersection, sont perpendiculaires entre elles.
\item \textbf{Critère.} Si l'un des deux plans contient une droite orthogonale à l'autre plan, alors les deux plans sont orthogonaux.
\end{enumerate}
\end{corrige}

\courssub{Applications directes : calculs, pyramide, tétraèdre, escalier et cube}

\begin{corrige}[Exercice 4 — Calculs directs : angle droite/plan et distance]
Données : Plan $\mathcal{P} : 2x - y + 2z = 5$ ; Droite $\Delta : \begin{cases} x = 1+2t \\ y = -1-t \\ z = 3+2t \end{cases}$ ; Point $A(1\,;-1\,;3) \in \Delta$.
\begin{enumerate}
\item Un vecteur normal à $\mathcal{P}$ est $\vec{n}\begin{pmatrix}2\\-1\\2\end{pmatrix}$ (coefficients de l'équation cartésienne). Un vecteur directeur de $\Delta$ est $\vec{u}\begin{pmatrix}2\\-1\\2\end{pmatrix}$. On a $\vec{u} = \vec{n}$ : la droite $\Delta$ est dirigée par un vecteur normal au plan, donc $\Delta \perp \mathcal{P}$.
\item La distance de $A$ à $\mathcal{P}$ :
\[ d(A,\mathcal{P}) = \frac{|2(1) - (-1) + 2(3) - 5|}{\sqrt{2^2+(-1)^2+2^2}} = \frac{|2+1+6-5|}{\sqrt{9}} = \frac{4}{3}. \]
Donc $d(A,\mathcal{P}) = \dfrac{4}{3}$.
\item Le projeté orthogonal $H$ de $A$ sur $\mathcal{P}$ est l'intersection de $\Delta$ (perpendiculaire à $\mathcal{P}$) avec $\mathcal{P}$. On injecte le paramétrage dans l'équation du plan :
\[ 2(1+2t) - (-1-t) + 2(3+2t) = 5 \iff 2+4t+1+t+6+4t = 5 \iff 9t + 9 = 5 \iff t = -\frac{4}{9}. \]
D'où $H\left(1-\dfrac{8}{9}\,;\,-1+\dfrac{4}{9}\,;\,3-\dfrac{8}{9}\right)$, c'est-à-dire $H\left(\dfrac{1}{9}\,;\,-\dfrac{5}{9}\,;\,\dfrac{19}{9}\right)$.

\emph{Vérification :} $AH = |t|\cdot\|\vec{u}\| = \dfrac{4}{9}\times 3 = \dfrac{4}{3}$ : cohérent avec la question 2.
\end{enumerate}
\end{corrige}

\begin{corrige}[Exercice 5 — Application : pyramide à base carrée]
Soit $SABCD$ une pyramide à base carrée $ABCD$ de côté $6$ cm, centre $O$, hauteur $SO = 8$ cm ($S$ à la verticale de $O$). $M$ est le milieu de $[SA]$. Le plan $(MBC)$ coupe $(SD)$ en $N$.
\begin{enumerate}
\item \textbf{$(SO) \perp (ABCD)$ :} Dans le carré $ABCD$, les diagonales $(AC)$ et $(BD)$ sont sécantes en $O$. Comme $SO$ est la hauteur, $(SO) \perp (AC)$ et $(SO) \perp (BD)$. D'après le critère (deux droites sécantes du plan), $(SO) \perp (ABCD)$.
\item \textbf{Calcul de $SA$ :} $O$ centre du carré $\implies OA = \frac{AC}{2} = \frac{6\sqrt{2}}{2} = 3\sqrt{2}$ cm. Triangle $SOA$ rectangle en $O$ (car $(SO)\perp(ABCD)$) :
\[ SA = \sqrt{SO^2 + OA^2} = \sqrt{8^2 + (3\sqrt{2})^2} = \sqrt{64 + 18} = \sqrt{82} \approx \num{9,06} \text{ cm}. \]
\item \textbf{Angle $(SA), (ABCD)$ :} Le projeté orthogonal de $S$ sur $(ABCD)$ est $O$, donc l'angle cherché est $\widehat{SAO}$. Dans $SOA$ rectangle en $O$ :
\[ \tan\widehat{SAO} = \frac{SO}{OA} = \frac{8}{3\sqrt{2}} \approx \num{1,886} \quad\Rightarrow\quad \widehat{SAO} \approx 62^\circ. \]
\item \textbf{Nature de $MBCN$ :} $M$ milieu de $[SA]$, $N$ intersection de $(MBC)$ avec $(SD)$. Dans le plan $(SAD)$, $M \in (SA)$ et la trace de $(MBC)$ sur $(SAD)$ est parallèle à $(AD)$ (car $(BC)\parallel(AD)$). Cette trace passe par $M$ et coupe $[SD]$ en $N$. D'après le théorème des milieux (ou de Thalès) dans le triangle $SAD$, $N$ est le milieu de $[SD]$. On a donc $(MN) \parallel (AD) \parallel (BC)$ et $MN = \frac{AD}{2} = 3$ cm, tandis que $BC = 6$ cm. Le quadrilatère $MBCN$ a deux côtés parallèles de longueurs différentes : c'est un \cle{trapèze}. De plus, par symétrie de la pyramide régulière, $MB = NC$ : c'est un \cle{trapèze isocèle}.
\end{enumerate}
\end{corrige}

\begin{attention}[Note importante sur l'Exercice 5]
Dans une pyramide régulière à base carrée, la droite $(BC)$ est \textbf{parallèle} au plan $(SAD)$ (car $(BC)\parallel(AD)\subset(SAD)$) et non orthogonale. Elle est orthogonale au plan $(SOJ)$ où $J$ est le milieu de $[BC]$, ce qui donne la hauteur du trapèze $MBCN$. Toute affirmation d'orthogonalité entre $(BC)$ et $(SAD)$ est fausse.
\end{attention}

\begin{corrige}[Exercice 6 — Application : tétraèdre rectangle]
$SABC$ tétraèdre tel que $(SA) \perp (ABC)$, triangle $ABC$ rectangle en $A$, $AB=5$ cm, $AC=12$ cm, $SA=4$ cm.
\begin{enumerate}
\item \textbf{Volume :} $\mathcal{A}_{ABC} = \frac{AB\times AC}{2} = \frac{5\times 12}{2} = 30$ cm$^2$. Hauteur $SA = 4$ cm.
\[ V = \frac{1}{3}\times 30 \times 4 = 40 \text{ cm}^3. \]
\item \textbf{Aire $\mathcal{A}_{SBC}$ :} $BC = \sqrt{AB^2+AC^2} = \sqrt{25+144} = 13$ cm. Hauteur $AH$ de $ABC$ issue de $A$ : $AH = \frac{AB\cdot AC}{BC} = \frac{60}{13}$ cm. Comme $(SA)\perp(ABC)$, le triangle $SAH$ est rectangle en $H$. Hauteur $SH$ de $SBC$ issue de $S$ :
\[ SH = \sqrt{SA^2 + AH^2} = \sqrt{16 + \frac{3600}{169}} = \sqrt{\frac{6304}{169}} = \frac{\sqrt{6304}}{13} = \frac{4\sqrt{394}}{13}. \]
\[ \mathcal{A}_{SBC} = \frac{1}{2}\times BC \times SH = \frac{1}{2}\times 13 \times \frac{4\sqrt{394}}{13} = 2\sqrt{394} \approx \num{39,70} \text{ cm}^2. \]
\item \textbf{Distance de $A$ au plan $(SBC)$ :} $S$ appartient au plan $(SBC)$, on calcule la distance du sommet $A$ à la face opposée $(SBC)$ (hauteur du tétraèdre relative à la base $SBC$) :
\[ h = \frac{3V}{\mathcal{A}_{SBC}} = \frac{120}{2\sqrt{394}} = \frac{60}{\sqrt{394}} \approx \num{3,02} \text{ cm}. \]
\item \textbf{Angle entre $(SAB)$ et $(ABC)$ :} Les plans se coupent selon $(AB)$. Comme $(SA)\perp(ABC)$ et $(SA)\subset(SAB)$, d'après le critère, $(SAB)\perp(ABC)$. L'angle diédral vaut $90^\circ$.
\end{enumerate}
\end{corrige}

\begin{corrige}[Exercice 7 — Contexte technique : escalier droit]
Un escalier droit a des marches de giron $28$ cm et de hauteur $17$ cm. La norme impose un angle de pente entre $25^\circ$ et $42^\circ$.
\begin{enumerate}
\item Dans un repère 2D (horizontal, vertical), un vecteur directeur de la ligne de foulée (limon) est $\vec{u}(28\,;\,17)$.
\item L'angle $\alpha$ de la ligne de pente avec l'horizontale vérifie :
\[ \tan\alpha = \frac{17}{28} \approx \num{0,607} \quad\Rightarrow\quad \alpha \approx 31^\circ. \]
\item Longueur du limon pour une marche : $\sqrt{28^2+17^2} = \sqrt{784+289} = \sqrt{1073} \approx \num{32,76}$ cm. Pour $15$ marches :
\[ L = 15\sqrt{1073} \approx \num{491,4} \text{ cm} \approx \num{4,91} \text{ m}. \]
\item On a $\alpha \approx 31^\circ \in [25^\circ\,;\,42^\circ]$ : l'escalier est \cle{conforme} à la norme.
\end{enumerate}
\end{corrige}

\begin{corrige}[Exercice 8 — Démonstration : cube et orthogonalité]
Soit $ABCDEFGH$ un cube de côté $a$. $I$ milieu de $[CG]$. On munit l'espace du repère $(A\,;\overrightarrow{AB},\overrightarrow{AD},\overrightarrow{AE})$ : $A(0;0;0)$, $B(a;0;0)$, $D(0;a;0)$, $E(0;0;a)$, $C(a;a;0)$, $H(0;a;a)$, $G(a;a;a)$, $I(a\,;\,a\,;\,\frac{a}{2})$.
\begin{enumerate}
\item \textbf{Montrer $(AC) \perp (BDH)$ :} $\overrightarrow{AC}=\begin{pmatrix}a\\a\\0\end{pmatrix}$. Dans $(BDH)$ : $\overrightarrow{BD}=\begin{pmatrix}-a\\a\\0\end{pmatrix}$ et $\overrightarrow{BH}=\begin{pmatrix}-a\\a\\a\end{pmatrix}$.
\[ \overrightarrow{AC}\cdot\overrightarrow{BD} = -a^2+a^2+0 = 0 \quad;\quad \overrightarrow{AC}\cdot\overrightarrow{BH} = -a^2+a^2+0 = 0. \]
$(AC)$ est orthogonale aux deux droites sécantes $(BD)$ et $(BH)$ du plan $(BDH)$ : d'après le critère, $(AC)\perp(BDH)$.
\item \textbf{Montrer $(AI) \perp (BD)$ :} $\overrightarrow{AI}=\begin{pmatrix}a\\a\\a/2\end{pmatrix}$, $\overrightarrow{BD}=\begin{pmatrix}-a\\a\\0\end{pmatrix}$. $\overrightarrow{AI}\cdot\overrightarrow{BD} = -a^2+a^2+0 = 0$, donc $(AI)\perp(BD)$.
\emph{Alternative géométrique :} $(BD)\perp(AC)$ et $(BD)\perp(CG)$ (directions orthogonales). $(AC)$ et $(CG)$ sont sécantes en $C$ dans le plan $(ACG)$. $(BD)$ coupe $(AC)$ en $O$ (centre du cube). Par translation, $(BD)$ est orthogonale à deux droites sécantes de $(ACG)$ passant par $O$, donc $(BD)\perp(ACG)$. Or $(AI)\subset(ACG)$ (car $I\in[CG]$), donc $(BD)\perp(AI)$.
\item \textbf{Distance de $A$ au plan $(BDH)$ :} Le plan $(BDH)$ a pour équation $x+y=a$ (vecteur normal $(1;1;0)$, passant par $B(a;0;0)$).
\[ d(A,(BDH)) = \frac{|0+0-a|}{\sqrt{1^2+1^2}} = \frac{a}{\sqrt{2}} = \frac{a\sqrt{2}}{2}. \]
\item \textbf{Angle entre $(AIC)$ et $(BDH)$ :} Le plan $(AIC)$ est le plan $(ACG)$ (car $I\in[CG]$). Il contient la droite $(AC)$ qui est orthogonale au plan $(BDH)$ (question 1). D'après le critère d'orthogonalité de deux plans, $(AIC)\perp(BDH)$. L'angle entre les plans vaut donc $90^\circ$.
\end{enumerate}
\end{corrige}

\courssub{Sujets type Probatoire : synthèse et résolution de problèmes complexes}

\begin{attention}[Exercice 9 — Corrections d'énoncé]
L'énoncé original contient des erreurs : 
\begin{itemize}
\item Q2 : « Montrer que $(BC)\perp(SAD)$ » est \textbf{faux}. $(BC)\parallel(AD)\subset(SAD) \implies (BC)\parallel(SAD)$.
\item Q3 : « Montrer que $(BC)\perp\mathcal{P}$ » est \textbf{faux} car $(BC)\subset\mathcal{P}$.
\item Q4 : « Montrer que $\mathcal{P}\perp(SAD)$ » est \textbf{faux} (calcul des normales donne un produit scalaire non nul).
\end{itemize}
Le corrigé ci-dessous traite les questions corrigées (parallélisme, orthogonalité $(BC)\perp(SOJ)$, volume, angle).
\end{attention}

\begin{corrige}[Exercice 9 — Sujet type Probatoire : pyramide régulière]
$SABCD$ pyramide régulière, base carrée $ABCD$ de côté $6$ cm, hauteur $SO=8$ cm ($O$ centre). $M$ milieu de $[SA]$, $\mathcal{P}=(MBC)$, $N = \mathcal{P}\cap(SD)$.
\begin{enumerate}
\item $(SO)\perp(ABCD)$ (hauteur de pyramide régulière). Comme $(BC)\subset(ABCD)$, $(SO)\perp(BC)$.
\item \textbf{Parallélisme $(BC)\parallel(SAD)$ :} $(BC)\parallel(AD)$ (carré) et $(AD)\subset(SAD)$, donc $(BC)\parallel(SAD)$.
\textbf{Orthogonalité $(BC)\perp(SOJ)$ :} Soit $J$ milieu de $[BC]$. $(BC)\perp(AB)$ (carré) et $(OJ)\parallel(AB)$ ($O,J$ milieux), donc $(BC)\perp(OJ)$. De plus $(SO)\perp(ABCD) \implies (SO)\perp(BC)$. $(OJ)$ et $(SO)$ sont sécantes en $O$ dans le plan $(SOJ)$, donc $(BC)\perp(SOJ)$. En particulier $(BC)\perp(SJ)$, ce qui donne la hauteur du trapèze $MBCN$.
\item $(BC)\subset\mathcal{P}$, donc $(BC)$ n'est pas orthogonale à $\mathcal{P}$.
\item \textbf{Plans $\mathcal{P}$ et $(SAD)$ non orthogonaux :} Coordonnées : $A(-3;-3;0)$, $B(3;-3;0)$, $C(3;3;0)$, $D(-3;3;0)$, $S(0;0;8)$, $M(-1,5;-1,5;4)$, $N(-1,5;1,5;4)$.
Vecteurs de $\mathcal{P}$ : $\overrightarrow{BC}=(0;6;0)$, $\overrightarrow{BM}=(4,5;1,5;4)$. Normale $\vec{n}_{\mathcal{P}} = \overrightarrow{BC}\times\overrightarrow{BM} = (24;0;27) \sim (8;0;9)$.
Vecteurs de $(SAD)$ : $\overrightarrow{AD}=(0;6;0)$, $\overrightarrow{AS}=(3;3;8)$. Normale $\vec{n}_{SAD} = \overrightarrow{AD}\times\overrightarrow{AS} = (48;0;-18) \sim (8;0;-3)$.
$\vec{n}_{\mathcal{P}}\cdot\vec{n}_{SAD} = 64 + 0 - 27 = 37 \neq 0$ : plans non orthogonaux.
\item \textbf{Aire $MBCN$ :} $MN=3$, $BC=6$, hauteur $h = SJ = \sqrt{SO^2+OJ^2} = \sqrt{8^2+3^2} = \sqrt{73}$? \emph{Attention :} $J(3;0;0)$, $S(0;0;8) \implies SJ = \sqrt{9+64}=\sqrt{73}$. Mais dans le trapèze, la hauteur est la distance entre les parallèles $(MN)$ et $(BC)$. $M(-1,5;-1,5;4)$, $B(3;-3;0)$. Vecteur $\overrightarrow{MB}=(4,5;-1,5;-4)$. Projection sur direction orthogonale à $(BC)$ (i.e. plan $xOz$) : distance entre droites $x=-1,5, z=4$ et $x=3, z=0$ : $\sqrt{(4,5)^2+4^2} = \sqrt{36,25} = \frac{\sqrt{145}}{2}$.
\[ \mathcal{A}_{MBCN} = \frac{MN+BC}{2}\times h = \frac{9}{2}\times\frac{\sqrt{145}}{2} = \frac{9\sqrt{145}}{4}. \]
\textbf{Volume $V_{S-MBCN}$ :} Plan $\mathcal{P} : 8x+9z=24$. Distance $d(S,\mathcal{P}) = \frac{|0+72-24|}{\sqrt{145}} = \frac{48}{\sqrt{145}}$.
\[ V = \frac{1}{3}\times\frac{9\sqrt{145}}{4}\times\frac{48}{\sqrt{145}} = 36 \text{ cm}^3. \]
\item \textbf{Angle $\mathcal{P}$ et $(ABCD)$ :} Normales $(8;0;9)$ et $(0;0;1)$.
\[ \cos\theta = \frac{|9|}{\sqrt{145}\times 1} = \frac{9}{\sqrt{145}} \approx \num{0,747} \quad\Rightarrow\quad \theta \approx 42^\circ.
\]
\end{enumerate}
\textbf{Barème indicatif (10 pts) :} Q1 : 2 pts ; Q2 : 2 pts ; Q3–Q4 (corrigées) : 2 pts ; Q5 : 2,5 pts ; Q6 : 1,5 pt.
\textbf{Vigilance :} citer explicitement le critère (« deux droites sécantes du plan ») ; justifier $N$ milieu par Thalès ; pour le volume, préciser base et hauteur ; arrondir l'angle au degré près.
\end{corrige}

\begin{corrige}[Exercice 10 — Sujet type Probatoire : tétraèdre et distances]
$ABCD$ tétraèdre tel que $(AD)\perp(BCD)$, $BCD$ équilatéral de côté $6$ cm, $AD=4$ cm. $I$ centre de gravité de $BCD$.
\begin{enumerate}
\item \textbf{$(BC)\perp(ADI)$ :} Dans $BCD$ équilatéral, $(DI)$ est médiane donc hauteur : $(DI)\perp(BC)$. De plus $(AD)\perp(BCD) \implies (AD)\perp(BC)$. $(DI)$ et $(AD)$ sont sécantes en $D$ dans le plan $(ADI)$. D'après le critère, $(BC)\perp(ADI)$.
\item \textbf{$(ABC)\perp(ADI)$ :} $(BC)\subset(ABC)$ et $(BC)\perp(ADI)$ (Q1). Le plan $(ABC)$ contient une droite orthogonale à $(ADI)$, donc $(ABC)\perp(ADI)$ (critère plans orthogonaux).
\item \textbf{Volume :} $\mathcal{A}_{BCD} = \frac{\sqrt{3}}{4}\times 6^2 = 9\sqrt{3}$ cm$^2$. Hauteur $AD=4$ cm.
\[ V = \frac{1}{3}\times 9\sqrt{3}\times 4 = 12\sqrt{3} \approx \num{20,78} \text{ cm}^3. \]
\item \textbf{Distance $A$ au plan $(BCD)$ :} Comme $(AD)\perp(BCD)$ en $D$, la distance est $AD = 4$ cm.
\item \textbf{Angle $(AB)$ et $(CD)$ :} Repère : $D(0;0;0)$, $B(6;0;0)$, $C(3;3\sqrt{3};0)$, $A(0;0;4)$.
$\overrightarrow{AB}=(6;0;-4)$, $\overrightarrow{CD}=(-3;-3\sqrt{3};0)$.
\[ \cos\theta = \frac{|\overrightarrow{AB}\cdot\overrightarrow{CD}|}{\|\overrightarrow{AB}\|\,\|\overrightarrow{CD}\|} = \frac{|-18|}{\sqrt{52}\times 6} = \frac{18}{12\sqrt{13}} = \frac{3}{2\sqrt{13}} \approx \num{0,416} \quad\Rightarrow\quad \theta \approx 65^\circ. \]
\item \textbf{Volume section par plan $\parallel(BCD)$ par $M$ milieu $[AB]$ :} Le plan coupe un tétraèdre $A-MNP$ \cle{réduction} de $A-BCD$ de rapport $k=\frac{1}{2}$ (milieu). Volumes dans le rapport $k^3 = \frac{1}{8}$.
\[ V_{A-MNP} = \frac{1}{8}\times 12\sqrt{3} = \frac{3\sqrt{3}}{2} \approx \num{2,60} \text{ cm}^3. \]
\end{enumerate}
\textbf{Barème indicatif (10 pts) :} Q1 : 2 pts ; Q2 : 1,5 pt ; Q3 : 1,5 pt ; Q4 : 1 pt ; Q5 : 2 pts ; Q6 : 2 pts.
\textbf{Vigilance :} Q1 justifier $(DI)\perp(BC)$ par « médiane = hauteur triangle équilatéral » ; Q2 citer critère plans orthogonaux ; Q6 écrire que rapport volumes = cube du rapport de réduction.
\end{corrige}

\begin{attention}[Exercice 11 — Correction de la section]
Avec $I$ milieu de $[EH]$ et $J$ milieu de $[FG]$, le plan $\mathcal{Q}=(AIJ)$ contient $A$ et la direction $\overrightarrow{IJ}\parallel\overrightarrow{AB}$, donc il contient la droite $(AB)$. La section du parallélépipède est le quadrilatère $ABJI$ (le plan ne rencontre pas $[DH]$ ni $[CG]$).
\end{attention}

\begin{corrige}[Exercice 11 — Sujet type Probatoire : parallélépipède rectangle]
$ABCDEFGH$ parallélépipède rectangle : $AB=12$ cm, $AD=5$ cm, $AE=4$ cm. $I$ milieu $[EH]$, $J$ milieu $[FG]$. $\mathcal{Q}=(AIJ)$. Repère $A(0;0;0)$, $\vec{\imath}=\overrightarrow{AB}$, $\vec{\jmath}=\overrightarrow{AD}$, $\vec{k}=\overrightarrow{AE}$.
$B(12;0;0)$, $D(0;5;0)$, $E(0;0;4)$, $H(0;5;4)$, $I(0;2,5;4)$, $J(12;2,5;4)$. $\mathcal{Q} : 8y-5z=0$.
\begin{enumerate}
\item \textbf{$\mathcal{Q}\perp(AEH)$ :} $(AB)\perp(AE)$ et $(AB)\perp(AD)$ (faces rectangulaires), $(AE)$ et $(AD)$ sécantes en $A$ dans $(AEH)$ : $(AB)\perp(AEH)$. Comme $(AB)\subset\mathcal{Q}$, $\mathcal{Q}$ contient une droite orthogonale à $(AEH)$, donc $\mathcal{Q}\perp(AEH)$.
\item \textbf{Angle $\mathcal{Q}$ et $(ABCD)$ :} Intersection : $(AB)$. Dans $\mathcal{Q}$, $(AI)\perp(AB)$ (car $(AB)\perp(AEH)$ et $(AI)\subset(AEH)$). Dans $(ABCD)$, $(AD)\perp(AB)$. Angle $\theta$ entre $(AI)$ et $(AD)$. Projeté $I'$ de $I$ sur $(ABCD)$ : $I'(0;2,5;0)$ (milieu $[AD]$). $AI'=2,5$, $II'=4$.
\[ \tan\theta = \frac{II'}{AI'} = \frac{4}{2,5} = 1,6 \quad\Rightarrow\quad \theta \approx 58^\circ. \]
\item \textbf{Section $ABJI$ :} $(AB)\parallel(IJ)$ (direction $\vec{\imath}$). $(AB)\perp(AI)$ (prouvé Q1). $ABJI$ est un \cle{rectangle}.
$AI = \sqrt{AE^2+EI^2} = \sqrt{4^2+2,5^2} = \sqrt{22,25} = \frac{\sqrt{89}}{2} \approx \num{4,72}$ cm.
$\mathcal{A}_{ABJI} = AB\times AI = 12\times\frac{\sqrt{89}}{2} = 6\sqrt{89} \approx \num{56,6}$ cm$^2$.
\item \textbf{Angle $\mathcal{Q}$ et $(ABFE)$ :} Intersection : $(AB)$. Dans $\mathcal{Q}$, $(AI)\perp(AB)$ ; dans $(ABFE)$, $(AE)\perp(AB)$. Angle = $\widehat{IAE}$.
\[ \tan\widehat{IAE} = \frac{EI}{AE} = \frac{2,5}{4} = 0,625 \quad\Rightarrow\quad \widehat{IAE} \approx 32^\circ. \]
\item \textbf{Distance $C$ au plan $\mathcal{Q}$ :} $C(12;5;0)$, $\mathcal{Q}: 8y-5z=0$.
\[ d(C,\mathcal{Q}) = \frac{|8\times 5 - 5\times 0|}{\sqrt{8^2+(-5)^2}} = \frac{40}{\sqrt{89}} = \frac{40\sqrt{89}}{89} \approx \num{4,24} \text{ cm}. \]
\end{enumerate}
\textbf{Barème indicatif (10 pts) :} Q1 : 2 pts ; Q2 : 2 pts ; Q3 : 2,5 pts ; Q4 : 1,5 pt ; Q5 : 2 pts.
\textbf{Vigilance :} pour angle diédral, prendre dans \emph{chaque} plan une droite $\perp$ à l'intersection \emph{au même point} ; justifier rectangle par orthogonalité prouvée ; formule distance point-plan : valeur absolue au numérateur, racine somme carrés au dénominateur.
\end{corrige}

\newpage

\coursec{Fiche bilan}

\begin{popfigure}
\centering
\begin{center}\tcbox[colback=popink,colframe=popink,coltext=white,arc=9pt,boxsep=4pt,left=6pt,right=6pt]{\bfseries\small Orthogonalité dans l'espace}\end{center}
\vspace{-2pt}
\begin{tcbraster}[raster columns=3,raster equal height=rows,raster column skip=6pt,raster row skip=6pt,enhanced,blankest]
\begin{tcolorbox}[enhanced,colback=white,colframe=popblue,boxrule=0.9pt,arc=3pt,title={Droites orthogonales},fonttitle=\bfseries\scriptsize,coltitle=white,colbacktitle=popblue,left=4pt,right=4pt,top=2pt,bottom=2pt,boxsep=1pt,toptitle=1.5pt,bottomtitle=1.5pt,halign=left]
\begin{itemize}[nosep,leftmargin=*,itemsep=1pt,font=\scriptsize]
\item $(d)\perp(\Delta)$ si parallèles sécantes perpendiculaires
\item Non nécessairement sécantes !
\end{itemize}
\end{tcolorbox}
\begin{tcolorbox}[enhanced,colback=white,colframe=poporange,boxrule=0.9pt,arc=3pt,title={Droite $\perp$ plan},fonttitle=\bfseries\scriptsize,coltitle=white,colbacktitle=poporange,left=4pt,right=4pt,top=2pt,bottom=2pt,boxsep=1pt,toptitle=1.5pt,bottomtitle=1.5pt,halign=left]
\begin{itemize}[nosep,leftmargin=*,itemsep=1pt,font=\scriptsize]
\item Critère : $\perp$ à 2 droites sécantes du plan
\item Alors $\perp$ à toute droite du plan
\end{itemize}
\end{tcolorbox}
\begin{tcolorbox}[enhanced,colback=white,colframe=popgreen,boxrule=0.9pt,arc=3pt,title={Plans orthogonaux},fonttitle=\bfseries\scriptsize,coltitle=white,colbacktitle=popgreen,left=4pt,right=4pt,top=2pt,bottom=2pt,boxsep=1pt,toptitle=1.5pt,bottomtitle=1.5pt,halign=left]
\begin{itemize}[nosep,leftmargin=*,itemsep=1pt,font=\scriptsize]
\item Un plan contient une droite $\perp$ à l'autre
\end{itemize}
\end{tcolorbox}
\begin{tcolorbox}[enhanced,colback=white,colframe=poppurple,boxrule=0.9pt,arc=3pt,title={Propriétés clés},fonttitle=\bfseries\scriptsize,coltitle=white,colbacktitle=poppurple,left=4pt,right=4pt,top=2pt,bottom=2pt,boxsep=1pt,toptitle=1.5pt,bottomtitle=1.5pt,halign=left]
\begin{itemize}[nosep,leftmargin=*,itemsep=1pt,font=\scriptsize]
\item Unicité de la perpendiculaire par un point
\item Parallélisme et orthogonalité
\end{itemize}
\end{tcolorbox}
\begin{tcolorbox}[enhanced,colback=white,colframe=poppink,boxrule=0.9pt,arc=3pt,title={Applications techniques},fonttitle=\bfseries\scriptsize,coltitle=white,colbacktitle=poppink,left=4pt,right=4pt,top=2pt,bottom=2pt,boxsep=1pt,toptitle=1.5pt,bottomtitle=1.5pt,halign=left]
\begin{itemize}[nosep,leftmargin=*,itemsep=1pt,font=\scriptsize]
\item Mur $\perp$ sol, poteau, équerre
\item Distances, hauteurs
\end{itemize}
\end{tcolorbox}
\end{tcbraster}

\legende{Carte mentale de la leçon : l'orthogonalité dans l'espace.}
\end{popfigure}

\begin{bilan}
\textbf{Définitions clés}
\begin{itemize}
  \item \cle{Droites orthogonales} : deux droites de l'espace sont orthogonales si les parallèles à ces droites menées par un même point sont \textbf{perpendiculaires}. Elles ne sont pas forcément sécantes.
  \item \cle{Droite orthogonale à un plan} : une droite $(d)$ est orthogonale à un plan $(P)$ si elle est orthogonale à \textbf{toute} droite de $(P)$. On note $(d)\perp(P)$.
  \item \cle{Plans orthogonaux} : deux plans sont orthogonaux (perpendiculaires) si l'un contient une droite orthogonale à l'autre. On note $(P)\perp(Q)$.
\end{itemize}

\textbf{Théorèmes et formules à connaître par cœur}
\begin{center}
\begin{tabularx}{\linewidth}{|l|X|}
\hline
\rowcolor{popblueL} \textbf{Résultat} & \textbf{Énoncé} \\
\hline
Critère droite $\perp$ plan & Si $(d)$ est orthogonale à \textbf{deux droites sécantes} de $(P)$, alors $(d)\perp(P)$. \\
\hline
Propriété fondamentale & Si $(d)\perp(P)$, alors $(d)$ est orthogonale à \textbf{toute} droite de $(P)$. \\
\hline
Unicité & Par un point donné, il passe une \textbf{unique} droite orthogonale à un plan donné, et un \textbf{unique} plan orthogonal à une droite donnée. \\
\hline
Critère plans $\perp$ & Si un plan $(P)$ contient une droite orthogonale à $(Q)$, alors $(P)\perp(Q)$. \\
\hline
Propriété des plans $\perp$ & Si $(P)\perp(Q)$, toute droite de $(P)$ orthogonale à l'intersection $(P)\cap(Q)$ est orthogonale à $(Q)$. \\
\hline
Parallélisme & Deux droites orthogonales au même plan sont parallèles ; deux plans orthogonaux à la même droite sont parallèles. \\
\hline
\end{tabularx}
\end{center}

\textbf{Méthodes express}
\begin{itemize}
  \item Prouver $(d)\perp(P)$ : exhiber \textbf{deux droites sécantes} de $(P)$ orthogonales à $(d)$ (souvent dans deux faces différentes d'un solide).
  \item Prouver $(P)\perp(Q)$ : trouver dans $(P)$ une droite orthogonale à $(Q)$.
  \item Prouver $(d)\perp(d')$ dans l'espace : montrer que $(d)$ est orthogonale à un plan contenant $(d')$.
  \item Calculs : se ramener à un \textbf{triangle rectangle} dans un plan bien choisi, puis Pythagore ou trigonométrie.
\end{itemize}
\end{bilan}

\begin{aretenir}[Checklist avant le Probatoire]
\begin{itemize}
  \item[$\square$] Je sais que deux droites orthogonales de l'espace ne sont pas forcément sécantes.
  \item[$\square$] Je sais définir une droite orthogonale à un plan.
  \item[$\square$] Je connais le critère : orthogonalité à \textbf{deux droites sécantes} du plan (une seule ne suffit pas !).
  \item[$\square$] Je sais utiliser la propriété fondamentale : $(d)\perp(P)$ implique $(d)$ orthogonale à toute droite de $(P)$.
  \item[$\square$] Je sais prouver que deux plans sont orthogonaux en exhibant une droite de l'un orthogonale à l'autre.
  \item[$\square$] Je connais les théorèmes d'unicité (droite $\perp$ plan, plan $\perp$ droite par un point).
  \item[$\square$] Je sais que deux droites orthogonales au même plan sont parallèles.
  \item[$\square$] Je sais repérer un plan pertinent pour ramener un calcul de l'espace à un triangle rectangle.
  \item[$\square$] Je sais appliquer le théorème de Pythagore dans l'espace (diagonale d'un pavé : $d=\sqrt{a^2+b^2+c^2}$).
  \item[$\square$] Je sais rédiger une justification complète : hypothèses, théorème cité, conclusion.
  \item[$\square$] Je sais faire le lien avec des situations concrètes : poteau vertical sur un sol horizontal, mur perpendiculaire au sol.
\end{itemize}
\end{aretenir}

\findecours

\end{document}
